
%% ── 1-Introduccion/1-Introduccion.tex
\subsection*{1 Introduccion}
% !TeX root = ../tesis.tex

\section*{Introducción}
\label{ch:introduccion}

Debido al crecimiento demográfico exponencial, la Ciudad de México enfrenta
desafíos estructurales complejos, entre los cuales destaca la movilidad intraurbana y
metropolitana como el eje central de la presente investigación. Esta problemática
tiene raíces históricas profundas; si bien desde el Porfiriato el automóvil fue
concebido como un símbolo de modernidad, su consolidación como eje de la
planeación urbana ocurrió entre las décadas de 1920 y 1930. Durante este periodo,
la expansión urbana y la motorización transformaron aceleradamente la morfología
de la capital \parencite{franco2022}.

En consecuencia, durante la segunda mitad del siglo XX, se priorizó un modelo de
ciudad volcado hacia el autotransporte mediante la construcción de arterias viales
para vehículos particulares y de carga. Este desarrollo de infraestructura y servicios
de transporte público fue el factor determinante en la ubicación y conectividad de los
nuevos fraccionamientos con el centro de la ciudad \parencite{franco2022}.

En este contexto, la dinámica de movilidad respondió a la constante necesidad
de desplazamiento entre distintos puntos de la capital. Una consecuencia directa
de ello fue la construcción de vialidades de mayor capacidad para sostener el
creciente volumen de traslados, lo cual propició, a su vez, la expansión de la
mancha urbana en sus dimensiones habitacional, laboral y de servicios.

\begin{quote}
``Entre éstas destacaron la Avenida de los Insurgentes, que se extendió hasta el
municipio de San Ángel como parte de los festejos del centenario de la consumación
de la Independencia en 1921 [\ldots] Hacia el norte, sobre la Calzada de Guadalupe
se construyeron barrios populares y para trabajadores industriales, tales como las
colonias Maza, Peralvillo, Valle Gómez, Vallejo y del Rastro.'' \parencite[p.~XX]{franco2022}
\end{quote}

Sin embargo, todo este avance de movilidad tuvo su auge para las rutas
concesionadas de camiones y transportes motorizados. A diferencia de los tranvías,
cuya implementación implicaba mayores costos y tiempos de instalación, los
camiones concesionados ofrecieron mayor versatilidad y cobertura \parencite{franco2022}.

Las carreteras representaron la principal vía de comunicación por toda la república
mexicana. Este crecimiento urbano evidenció la necesidad de ampliar y robustecer
las vialidades, derivando en la expansión de la mancha urbana. Diversos estudios
han señalado que las carreteras fungen como catalizadores de nuevos desarrollos
urbanos \parencite{garcia2019allroads}. Según \textcite{garcia2019allroads}, las autopistas
expanden las ciudades con mayor desarrollo urbanístico y reducen su densidad
residencial general, generando nuevos desarrollos más discontinuos, fragmentados,
rodeados de terreno sin urbanizar y más aislados de otras áreas artificiales
(traducción propia).

De este mismo texto podemos extraer una de las figuras que da estructura a nuestra
investigación: los patrones espaciales alternativos de desarrollo de suelo en el
proceso de suburbanización de la población, representados en tres figuras:
\textit{Compact, linear} y \textit{leapfrog} (Compacto, lineal y saltón). Según
\textcite{jia2010}, pueden distinguirse tres patrones espaciales de desarrollo urbano:
el desarrollo saltón muestra mayor fragmentación y aislamiento; el lineal implica
menor fragmentación pero mayor aislamiento; y el compacto, aunque más eficiente,
también puede derivar en expansión si la densidad residencial disminuye (traducción
propia).

\begin{figure}[htbp]
    \centering
    \figinclude[width=0.85\textwidth]{Figures/Cap1/Sprawl}
    \caption[Patrones de desarrollo urbano]{Patrones espaciales de desarrollo de suelo:
    compacto, lineal y saltón (\textit{Compact, Linear and Leapfrog development}).}
    \label{fig:sprawl}
    \small\textit{Fuente: García-López (2019, p.~3).}
\end{figure}

En el caso de la Ciudad de México, la expansión en colonias del sur y oriente a lo
largo de ejes como Tlalpan o Periférico Sur corresponde a un patrón lineal; mientras
que los nuevos desarrollos en Milpa Alta y Tláhuac muestran rasgos de patrón
saltón, al estar aislados de la mancha urbana consolidada.

En este marco, la presente investigación analiza la viabilidad y el impacto de
implementar un corredor circular de transporte público masivo en la \zmvm{},
evaluando su efecto sobre la estructura topológica de la red mediante el modelado
de grafos matemáticos. El caso de validación central es la sección sur del Anillo
Periférico de la \cdmx{}, cuya configuración concentra de manera paradigmática
las deficiencias estructurales del sistema metropolitano en su conjunto.

%-----------------------------------------------------------------------
\section{Problemática}
\label{sec:problematica}
%-----------------------------------------------------------------------

El Periférico de la Ciudad de México fue concebido como una infraestructura
estratégica para conectar diversos puntos de la capital. Indirectamente, creó los
límites más visibles para la ciudad, ocasionando la centralización de actividades
comerciales, legislativas y, en particular, de transporte. Con el paso de los años,
algunas zonas han experimentado crecimientos acelerados producto del
desbordamiento poblacional, haciendo más notoria la ausencia de transporte público
masivo en la sección sur del anillo. Las alcaldías más afectadas son Xochimilco,
Milpa Alta, Tláhuac y Tlalpan, que carecen de opciones de traslados de alta
capacidad y dependen principalmente de transporte concesionado como camiones
y vagonetas, además del uso intensivo del automóvil particular.

Esta situación responde, en buena medida, a la falta histórica de inversión en
transporte público, lo que ha fomentado el uso individualizado del automóvil y
disminuido la eficiencia del sistema de movilidad en su conjunto. Tal como señalan
especialistas en movilidad urbana \parencite{cano2002}, el impacto del transporte
motorizado se refleja en un incremento del tráfico, de modo que la construcción o
ampliación de vialidades urbanas \DIFdelbegin \DIFdel{sólo }\DIFdelend \DIFaddbegin \DIFadd{solo }\DIFaddend ofrece una solución temporal: al cabo de tres
a cinco años, estas nuevas infraestructuras tienden a saturarse por el aumento de
viajes inducidos.

%-----------------------------------------------------------------------
\section{Planteamiento del problema}
\label{sec:planteamiento}
%-----------------------------------------------------------------------

La problemática expuesta anteriormente revela que la Zona Metropolitana del Valle de México (ZMVM) padece una desarticulación 
topográfica profunda, donde el sistema de transporte masivo estructurado presenta una cobertura insuficiente en las franjas 
periféricas que componen el \textbf{Tercer y Cuarto Contorno} metropolitano \parencite{semovi2020diagnostico}. Esta limitación
técnica genera una dependencia crítica del transporte público concesionado de baja capacidad y del automóvil particular, 
lo cual resulta ineficiente para mitigar una demanda de movilidad que desborda las fronteras administrativas de la capital.

Diversos diagnósticos urbanos coinciden en que la red masiva (Metro y Metrobús) se ha consolidado bajo una lógica de centralización,
obligando a los habitantes de zonas periféricas en el \textbf{Norte (Ecatepec, Naucalpan)}, \textbf{Oriente (Iztapalapa, Chalco)} y 
\textbf{Sur (Tlalpan, Xochimilco)} a realizar desplazamientos radiales que saturan los nodos del núcleo central \parencite{obras2025}. 
Esta configuración no solo incrementa los tiempos de traslado —superando en promedio los 122 minutos para los habitantes del 
cuarto contorno—\DIFaddbegin \DIFadd{, }\DIFaddend sino que eleva la \textit{vulnerabilidad funcional} del sistema ante fallas operativas en hubs críticos como Pantitlán,
Indios Verdes o Taxqueña \parencite{loterovélez2014}.

Aunque el Anillo Periférico fue proyectado como una infraestructura estratégica para conectar la capital, su diseño priorizó históricamente
la automovilidad privada, dejando una brecha topográfica para el transporte masivo transversal. De este modo, surge la necesidad científica 
de evaluar la pertinencia de un corredor circular en el Anillo Periférico como un motor de recalibración geométrica. La adopción de estas 
geometrías anillares ha demostrado ser una solución robusta en megaciudades para descomprimir el centro y articular las periferias de manera
directa \parencite{efe2023}, hipótesis que será validada en esta investigación mediante el motor analítico \texttt{VFTModel}.

%-----------------------------------------------------------------------
\subsection{La morfología y topología en la periferia de la ZMVM}
\label{subsec:morfologia}
%-----------------------------------------------------------------------

La estructura urbana de la periferia metropolitana no constituye una extensión continua y ordenada de la traza central. 
Siguiendo el análisis de \textcite{ortiz2007}, la morfología de la Ciudad de México se define como una jerarquía de centros que da 
lugar a un ``\textit{Patchwork of offset grids}'' (mosaico de retículas desfasadas). En este modelo, cada sector periférico 
(Norte, Oriente, Poniente o Sur) actúa como un ``parche'' geométricamente aislado, cuya trama interna no se alinea con 
la estructura vial primaria de la ciudad.

Esta fragmentación morfológica es el factor determinante de la ineficiencia en las alcaldías y municipios del cinturón metropolitano. 
A diferencia de la \textbf{Ciudad Central}, donde la conectividad en red facilita desplazamientos multidireccionales, estos parches 
periféricos operan como islas conectadas únicamente por arterias radiales que tienden al colapso crónico. Desde la perspectiva del 
motor \texttt{VFTModel}, esta desconexión implica que los caminos mínimos entre dos nodos periféricos aledaños tienen 
un \textbf{Índice de Ruta Directa ($DI$)} significativamente elevado, ya que la red vial obliga a los usuarios a converger hacia 
el centro para luego retornar a su zona de destino \parencite{fortin2016innovative}.

Complementariamente, la teoría de redes de \textcite{owais2020} sostiene que en ciudades con estructuras no ortogonales, 
la carencia de conexiones transversales entre estos parches fragmentados genera una pérdida de robustez sistémica. 
Esta teoría valida la tesis central de este trabajo: la problemática de la periferia de la ZMVM no es únicamente una cuestión de 
capacidad vehicular, sino una \textbf{desconexión geométrica estructural}. Mediante el procesamiento de datos en \texttt{Apimetro}, 
se demostrará que la implementación de una ruta anillar en el Tercer Contorno permite ``coser'' estos parches urbanos, 
optimizando la topología del grafo metropolitano y reduciendo la fricción vial global.

%-----------------------------------------------------------------------
\section{Objetivos de la Investigación}
\label{sec:objetivos}
%-----------------------------------------------------------------------

\subsection{Objetivo General}
\label{subsec:obj-general}

Analizar la viabilidad y robustez de la red de transporte multimodal de la Zona Metropolitana
del Valle de México (ZMVM) mediante el modelado de grafos matemáticos y la aplicación del motor
analítico \textit{VFTModel} \textcite{cabrera2026VFTModel} y 
el motor de datos \textit{Apimetro} \textcite{cabrera2025apimetro}. Con el \DIFdelbegin \DIFdel{fín }\DIFdelend \DIFaddbegin \DIFadd{fin }\DIFaddend de determinar el impacto teórico de
la implementación \DIFaddbegin \DIFadd{de }\DIFaddend un corredor circular de transporte público masivo en el Anillo Periférico de 
la Ciudad de México, con el fin de mejorar la movilidad urbana y reducir la dependencia del transporte
concesionado (RTP, Corredor concesionado) y del automóvil privado.

\subsection{Objetivos específicos}
\label{subsec:obj-especificos}

\begin{enumerate}
    \item Sistematizar y procesar los feeds de datos abiertos (GTFS) de la totalidad 
    de la red de Movilidad Integrada de la Ciudad de México y de los sistemas 
    de transporte del Estado de México (Mexibús, Mexicable) 
    utilizando el motor de datos \textit{Apimetro} \textcite{cabrera2025apimetro} 
    para la construcción de un grafo multimodal georreferenciado.

    \item Implementar el motor \texttt{VFTModel} para calcular los indicadores base de eficiencia urbana, 
    tales como el Índice de Ruta Directa ($DI$) y el Tiempo de Viaje Promedio ($T$), 
    sobre la configuración topológica actual de la red metropolitana.

    \item Identificar nodos críticos y vulnerabilidades estructurales en la ZMVM a través del cálculo 
    algorítmico de la Centralidad de Intermediación ($B(v)$) y métricas de Robustez Geométrica ($\Delta E$), 
    evaluando el impacto sistémico de posibles fallas en estaciones de alta transferencia.

    \item Identificar las principales limitaciones en términos de cobertura, capacidad y
    eficiencia del transporte público en la zona, a partir del análisis de datos
    espaciales y temporales.

    \item Evaluar y simular escenarios comparativos entre la situación actual y la posible
    implementación de un corredor circular masivo, estimando su impacto en la
    reducción de tiempos de traslado, así como validar la pertinencia de corredores circulares 
    frente al modelo radial predominante, cuantificando la reducción neta en los tiempos de 
    traslado y la mejora en la interconectividad entre contornos periféricos.

    \item Proponer lineamientos de integración operativa y tecnológica basados en los resultados del modelo, 
    que favorezcan una movilidad más equitativa y reduzcan la fragmentación del servicio en las alcaldías y 
    municipios del sur y oriente de la \cdmx y la \zmvm.

\end{enumerate}

\subsection{Pregunta de Investigación}
\label{subsec:pregunta}

¿En qué medida la optimización topológica de la red de transporte de la \cdmx y la \zmvm, 
evaluada mediante el motor analítico \texttt{VFTModel} \textcite{cabrera2026VFTModel}, demuestra que la implementación 
de corredores anillares reduce la vulnerabilidad funcional y mejora la eficiencia global 
del sistema en comparación con el modelo radial vigente?


De manera complementaria, se plantean las siguientes preguntas secundarias:

\begin{itemize}
    \item ¿De qué manera la alta dependencia de nodos centrales (\textit{hubs}) 
    en la Ciudad Central y el Primer Contorno \textcite{semovi2020diagnostico} 
    compromete la robustez del sistema ante perturbaciones operativas?

    \item ¿Cómo impacta la brecha de conectividad lateral en el 
    Tercer Contorno en el incremento de 
    los tiempos de traslado para los usuarios provenientes de 
    la periferia profunda o Cuarto Contorno \textcite{semovi2020diagnostico}?

\end{itemize}

%-----------------------------------------------------------------------
\section{Justificación}
\label{sec:justificacion}
%-----------------------------------------------------------------------

La pertinencia de esta investigación se sustenta en dar una alternativa y explicación a la fragmentación 
estructural de la movilidad en la Zona Metropolitana del Valle de México, la cual se manifiesta de forma 
diferenciada según el modelo de \textbf{Contornos Metropolitanos} definido en los diagnósticos técnicos de movilidad urbana.

En la \textbf{Ciudad Central} y el \textbf{Primer Contorno}, el sistema presenta una alta densidad de infraestructura 
de transporte masivo que, si bien garantiza una cobertura amplia, genera una centralización excesiva de los flujos. 
Esta configuración radial-concéntrica obliga a que la mayoría de los viajes metropolitanos atraviesen nodos 
críticos como Pantitlán, Indios Verdes o Taxqueña, saturando su capacidad operativa y elevando la vulnerabilidad 
del sistema ante cualquier falla en estos puntos neurálgicos.

Por el contrario, el \textbf{Tercer Contorno} (o \textit{Middle-Ring}), que alberga zonas de alta demanda y 
servicios críticos como la Zona de Hospitales y Ciudad Universitaria, carece de una conectividad transversal robusta. 
La ausencia de rutas perimetrales eficientes en esta franja obliga a los usuarios a realizar trayectos multimodales hacia el núcleo central
para luego retornar a zonas periféricas aledañas, lo que se traduce matemáticamente en un incremento del Índice de Ruta Directa ($DI$) 
y una pérdida de eficiencia geométrica.

Finalmente, el análisis se extiende al \textbf{Cuarto Contorno} (periferia profunda), donde se registran los tiempos de traslado más elevados del sistema, 
superando en promedio los 122 minutos por viaje. En este sector, la dependencia del transporte de baja capacidad y la fragmentación de los servicios alimentadores 
hacia la red masiva profundizan la desigualdad socioespacial \textcite{semovi2020diagnostico}. El uso de grafos matemáticos y el motor \texttt{VFTModel} 
permite demostrar que la inserción de una ruta anillar en el Tercer Contorno no representa únicamente una mejora local, sino una estrategia de ``recalibración urbana''
que beneficia sistémicamente a los contornos más alejados al redistribuir los flujos y reducir la presión sobre el centro de la metrópoli.


\begin{quote}
``Otro nivel de complejidad se debe a que los servicios que conforman el sistema de
movilidad no fueron diseñados de origen para integrar una red física de transporte
para la Ciudad de México, ni menos aún, para dar servicio a la Zona Metropolitana
del Valle de México. Por lo anterior, su integración es uno de los puntos prioritarios
a atender con el fin de garantizar un servicio de movilidad integrado, digno e
incluyente.'' \parencite[p.~XX]{semovi2020diagnostico}
\end{quote}


\begin{figure}[htbp]
    \centering
    \figinclude[width=0.80\textwidth]{Figures/Cap1/SEMOVITC}
    \caption[Contornos Metropolitanos]{Contornos Metropolitanos y red de transporte
    público integrada.}
    \label{fig:semovi-contornos}
    \small\textit{Fuente: Diagnóstico Técnico de Movilidad para el Programa Integral de
    Movilidad (2019--2024, p.~16).}
\end{figure}

Esta falta de integración se refleja en la calidad de los viajes y genera severos
problemas de desigualdad en el acceso a la capital. Conforme aumenta la distancia
entre la red de transporte masivo y las zonas habitacionales, más viajes multimodales
se tienen que realizar, lo que implica más tiempo y gastos de traslado
\parencite{semovi2020diagnostico}.

El diagnóstico de SEMOVITC evidencia que las demarcaciones ubicadas en el \textbf{Tercer Contorno}, 
con especial énfasis en las alcaldías del sur de la metrópoli, enfrentan un déficit estructural de 
accesibilidad derivado de la configuración radial-concéntrica de la red masiva. 
Dicha desconexión topográfica justifica la necesidad técnica de explorar alternativas perimetrales 
que articulen transversalmente estas zonas sin la obligatoriedad de converger en el núcleo central, 
reduciendo así la centralidad de intermediación de los nodos históricos. 

Actualmente, la carencia de infraestructura de alta capacidad en estos sectores delega la movilidad 
en el transporte público concesionado ---compuesto primordialmente por colectivos, autobuses y vagonetas---,
el cual es utilizado por el 48\% de la población en al menos uno de sus tramos de viaje \parencite{semovi2020diagnostico}. 
Este servicio opera mayoritariamente bajo el esquema ``hombre-camión'', lo que fomenta una competencia ineficiente por 
el pasaje y eleva los costos asociados a la seguridad vial. 
Desde la perspectiva del motor analítico \texttt{VFTModel}, esta dependencia de un servicio poco confiable y fragmentado 
constituye la principal fuente de \textit{vulnerabilidad funcional} en la periferia, validando la urgencia 
de una recalibración geométrica a través de corredores anillares.

Un elemento significativo es el uso de un medio de transporte que presente una
geometría diferente a la convencional. Las llamadas rutas radiales son el tipo más
común:

\begin{quote}
``[\ldots] En ciudades mayores a los 300,000 habitantes este tipo de rutas empieza
a ser ineficiente ya que concentra los movimientos y no considera las necesidades
que se presentan entre otras áreas urbanas. Esto induce a que la distribución del
servicio se encuentre limitada a ciertas áreas de la ciudad y concentre las terminales
en las zonas de mayor densidad.'' \parencite[p.~XX]{molinero2002}
\end{quote}

\begin{figure}[htbp]
    \centering
    \figinclude[width=0.55\textwidth]{Figures/Cap1/ETR}
    \caption[Esquema de transporte radial]{Esquema del sistema de transporte radial.}
    \label{fig:transporte-radial}
    \small\textit{Fuente: Elaboración propia.}
\end{figure}

En cambio, las rutas circulares sirven de rutas conectoras con las radiales,
permitiendo una mejor distribución de los usuarios y una mejor utilización del parque
vehicular. Casos típicos de este tipo son las líneas circulares de los metros de
Londres y Moscú, o los circuitos de los autobuses de Guadalajara y otras ciudades
mexicanas \parencite{molinero2002}.

\begin{figure}[htbp]
    \centering
    \figinclude[width=0.55\textwidth]{Figures/Cap1/ETC}
    \caption[Esquema de transporte circular]{Esquema del sistema de transporte
    circular.}
    \label{fig:transporte-circular}
    \small\textit{Fuente: Elaboración propia con base en \textcite{molinero2002}.}
\end{figure}

A la fecha de realización de este estudio (2025), no existe en la Ciudad de México alguna ruta circular 
que sirva de soporte para el transporte público masivo colectivo. A pesar de que la Ciudad cuenta 
con una vialidad de trazo circular —el Anillo Periférico—, \DIFdelbegin \DIFdel{ésta }\DIFdelend \DIFaddbegin \DIFadd{esta }\DIFaddend se ha orientado históricamente al automóvil 
privado como resultado de políticas de automovilidad hegemónicas. En este contexto, la presente investigación 
se justifica en la necesidad técnica de evaluar la pertinencia de un corredor circular de transporte público 
masivo en el Anillo Periférico para revertir la fragmentación del sistema.

\subsection{Tercer y Cuarto Contorno: Beneficiarios de la Recalibración Anillar} \label{subsec:beneficiarios-contornos}
La implementación de un corredor de transporte masivo en el Anillo Periférico tiene como objetivo principal mitigar las 
desigualdades estructurales de movilidad que afectan primordialmente al Tercer Contorno (Middle-Ring) y al Cuarto Contorno 
(periferia profunda) de la ZMVM. Mientras que la Ciudad Central, el Primer Contorno y el Segundo Contorno 
—que integran las alcaldías centrales— presentan la mayor densidad de infraestructura de transporte masivo (Metro y Metrobús), 
las zonas periféricas han sido históricamente desatendidas por la red estructurada. Esta \DIFdelbegin \DIFdel{"}\DIFdelend \DIFaddbegin \DIFadd{«}\DIFaddend suposición de cobertura\DIFdelbegin \DIFdel{" }\DIFdelend \DIFaddbegin \DIFadd{» }\DIFaddend centralizada se
abordará con rigor matemático en el Capítulo 3: Conceptos e Indicadores, donde métricas como el Nivel de Cobertura (C) y la Fuerza Nodal (Ci) 
permitirán cuantificar el desajuste entre la oferta de transporte y la localización de la población en la periferia.

Intervenir en la sección sur del Anillo Periférico responde a la necesidad de consolidar el llamado \textit{Middle-Ring} o \DIFdelbegin \DIFdel{Anillo Medio''}\DIFdelend \DIFaddbegin \DIFadd{«Anillo Medio»}\DIFaddend .
Según el análisis estratégico del \textit{Suburban Rail Loop} de Melbourne \parencite{suburbanrailloop2021}, estos anillos se definen 
como zonas de transición que albergan infraestructuras críticas como hospitales y universidades, pero que tienden a carecer de una conectividad
lateral robusta. Para la Ciudad de México, las alcaldías del sur (\textbf{Tlalpan, Xochimilco y Tláhuac}) conforman funcionalmente este tercer contorno,
operando como \DIFdelbegin \DIFdel{parches'' }\DIFdelend \DIFaddbegin \DIFadd{«parches» }\DIFaddend aislados en la trama urbana que obligan a los usuarios a realizar viajes radiales ineficientes hacia el centro para conectar 
con zonas vecinas.

Por otro lado, el Cuarto Contorno se posiciona como el beneficiario indirecto más crítico. En esta zona, los habitantes experimentan los tiempos de 
traslado más elevados de la metrópoli, superando los 122 minutos por trayecto laboral. Al no existir rutas perimetrales, los flujos provenientes de 
la periferia profunda saturan los nodos centrales del sistema, elevando la vulnerabilidad funcional de la red. La inserción de un corredor circular 
permitiría replicar un efecto de ``recalibración urbana'' \parencite{suburbanrailloop2021}, transformando suburbios desconectados en una red policéntrica 
y permitiendo que el Periférico Sur deje de ser una barrera para convertirse en una arteria integradora que detone el desarrollo económico y social transversal.
Esta mejora en la eficiencia geométrica será demostrada algorítmicamente mediante el Índice de Ruta Directa (DI) y el cálculo del Tiempo de Viaje Promedio (T) 
en los capítulos posteriores.

%-----------------------------------------------------------------------
\section{Hipótesis}
\label{sec:hipotesis}
%-----------------------------------------------------------------------

La red de transporte multimodal de la \zmvm{} presenta una configuración
radial-concéntrica que genera vulnerabilidades funcionales cuantificables: una
centralidad de intermediación ($B(v)$) elevada en los nodos del núcleo central,
un Índice de Ruta Directa ($DI$) alto en los contornos periféricos y tiempos de
viaje promedio ($T$) que superan los 122 minutos en el Cuarto Contorno
\parencite{semovi2020diagnostico}. Se postula que la incorporación de un corredor
anillar de transporte público masivo en el Tercer Contorno ---con la sección sur del
Anillo Periférico como caso de validación--- producirá una reducción estadísticamente
significativa de dichos indicadores, demostrable mediante el motor analítico
\textit{VFTModel} \parencite{cabrera2026VFTModel}, y favorecerá una movilidad
más equitativa en las alcaldías y municipios del sur y oriente de la \zmvm{}.

%-----------------------------------------------------------------------
\section{Método}
\label{sec:metodo}
%-----------------------------------------------------------------------

La presente investigación adopta un enfoque \textbf{cuantitativo}
fundamentado en la teoría de grafos y el modelado de redes complejas de transporte.
Este enfoque permite representar formalmente la red multimodal de la \zmvm{} como
un grafo dirigido y ponderado $\grafo = (\nodos, \aristas, \peso{i}{j})$, sobre el
cual se calculan indicadores de eficiencia, vulnerabilidad y conectividad de manera
objetiva y reproducible \parencite{allauca2023, owais2020}.

La elección de este método responde a tres condiciones estructurales de la
investigación: (1)~la disponibilidad de datos abiertos estandarizados en formato
\gtfs{} para la totalidad de los sistemas de transporte masivo de la \cdmx{};
(2)~la existencia de motores computacionales de autoría propia ---\textit{Apimetro}
\parencite{cabrera2025apimetro} y \textit{VFTModel} \parencite{cabrera2026VFTModel}---
que permiten procesar dichos datos a escala metropolitana; y (3)~la naturaleza
cuantificable de las preguntas de investigación, centradas en la reducción de tiempos
de traslado, la distribución de centralidad y la robustez geométrica de la red.

A diferencia de los enfoques tradicionales de modelación basados en matrices
origen-destino, el análisis mediante grafos no requiere calibración de encuestas de
movilidad, lo que lo hace especialmente adecuado para la evaluación estructural de
redes existentes y la simulación de escenarios topológicos alternativos
\parencite{fortin2016innovative}.

Las herramientas y fuentes de datos que articulan el proceso metodológico son:

\begin{itemize}
    \item \textbf{Feeds \gtfs{}} de la red de Movilidad Integrada de la \cdmx{} y
    los sistemas del Estado de México (Mexibús, Mexicable), como fuente primaria
    de datos estructurados de la red de transporte.
    \item \textit{Apimetro} \parencite{cabrera2025apimetro}: motor de adquisición,
    normalización y construcción del grafo multimodal georreferenciado a partir de
    datos \gtfs{}.
    \item \textit{VFTModel} \parencite{cabrera2026VFTModel}: motor analítico para el
    cálculo de indicadores de eficiencia urbana ($DI$, $T$), centralidad de
    intermediación ($B(v)$) y robustez geométrica ($\Delta E$) sobre el grafo
    construido.
    \item \textbf{Python y Go (Golang)}: lenguajes de programación empleados en el
    procesamiento de datos, análisis espacial y desarrollo de los motores
    computacionales.
    \item \textbf{Sistemas de Información Geográfica (\gis{})}: QGIS y GeoPandas
    para la visualización, validación espacial y análisis territorial de resultados.
\end{itemize}

%-----------------------------------------------------------------------
\section{Metodología}
\label{sec:metodologia}
%-----------------------------------------------------------------------

La investigación se estructura en cuatro fases secuenciales e interdependientes,
cada una orientada a cubrir uno o más de los objetivos específicos planteados. El
hilo conductor del proceso es la construcción progresiva del análisis: de la
adquisición y modelado de los datos hacia el diagnóstico estructural de la red, la
simulación de escenarios y, finalmente, la formulación de lineamientos de política
pública en materia de movilidad \parencite{allauca2023}.

\subsection{Fase 1: Adquisición y construcción del grafo multimodal}
\label{subsec:fase1}

Esta fase tiene una ponderación del 25\% sobre el total de la investigación y cubre
el \textbf{Objetivo Específico~1}.

\begin{description}
    \item[Objetivo de la Fase:] Construir el grafo multimodal georreferenciado de
    la \zmvm{} a partir del procesamiento sistemático de los feeds \gtfs{} de la
    totalidad de los sistemas de transporte masivo de la región.

    \item[Actividades a Desarrollar:]~
    \begin{itemize}
        \item \textit{Adquisición y normalización de datos \gtfs{}}: Descarga y
        limpieza de los feeds oficiales mediante el motor \textit{Apimetro}
        \parencite{cabrera2025apimetro}. Este proceso incluye la validación de la
        integridad estructural de los archivos (\texttt{stops.txt},
        \texttt{routes.txt}, \texttt{trips.txt}, \texttt{stop\_times.txt}) conforme
        al estándar \gtfs{} \parencite{fortin2016innovative}.
        \item \textit{Construcción del grafo dirigido y ponderado}: Transformación
        de los datos normalizados en un grafo $\grafo = (\nodos, \aristas)$, donde
        cada nodo $v \in \nodos$ representa una estación o parada y cada arista
        $e \in \aristas$ un segmento de ruta ponderado por tiempo de viaje
        $\peso{i}{j}$.
        \item \textit{Georreferenciación y validación espacial}: Asociación de
        coordenadas geográficas a cada nodo y verificación de la topología resultante
        mediante herramientas \gis{}.
    \end{itemize}
\end{description}

\subsection{Fase 2: Análisis de la topología actual con \textit{VFTModel}}
\label{subsec:fase2}

Esta fase tiene una ponderación del 30\% y cubre los \textbf{Objetivos
Específicos~2 y~3}.

\begin{description}
    \item[Objetivo de la Fase:] Caracterizar cuantitativamente la estructura y
    eficiencia de la red de transporte actual mediante el cálculo de indicadores
    de teoría de redes sobre el grafo construido en la Fase~1.

    \item[Actividades a Desarrollar:]~
    \begin{itemize}
        \item \textit{Cálculo del Índice de Ruta Directa ($DI$) y Tiempo de Viaje
        Promedio ($T$)}: Aplicación del motor \textit{VFTModel}
        \parencite{cabrera2026VFTModel} para obtener los indicadores base de
        eficiencia en la configuración topológica actual. El $DI$ mide la desviación
        de los trayectos reales respecto a la distancia euclidiana óptima; valores
        elevados indican ineficiencia geométrica en la red \parencite{owais2020}.
        \item \textit{Cálculo de la Centralidad de Intermediación ($B(v)$)}:
        Identificación de los nodos con mayor flujo de rutas mínimas, localizando
        los puntos neurálgicos cuya falla comprometería la conectividad global del
        sistema \parencite{batista2015, loterovélez2014}.
        \item \textit{Evaluación de la Robustez Geométrica ($\Delta E$)}: Cálculo
        del impacto sistémico ante la eliminación simulada de nodos o aristas de alta
        centralidad, cuantificando la degradación de la eficiencia global de la red
        \parencite{loterovélez2014}.
    \end{itemize}
\end{description}

\subsection{Fase 3: Diagnóstico de cobertura y simulación de escenarios}
\label{subsec:fase3}

Esta fase tiene una ponderación del 35\% y cubre los \textbf{Objetivos
Específicos~4 y~5}.

\begin{description}
    \item[Objetivo de la Fase:] A partir del diagnóstico estructural obtenido en
    la Fase~2, identificar las principales brechas de cobertura y eficiencia del
    sistema actual, y evaluar el impacto topológico de la incorporación de un
    corredor anillar en el Tercer Contorno metropolitano.

    \item[Actividades a Desarrollar:]~
    \begin{itemize}
        \item \textit{Diagnóstico de cobertura y eficiencia}: Análisis espacial de
        los indicadores calculados para identificar zonas con déficit de conectividad
        lateral, tiempos de traslado elevados y dependencia de nodos de alta
        intermediación \parencite{mishra2012, ugur2025}.
        \item \textit{Diseño del escenario propuesto}: Modelado del corredor circular
        sobre la sección sur del Anillo Periférico como un conjunto de nuevos nodos
        $\nodos'$ y aristas $\aristas'$ incorporados al grafo original, con pesos
        $\peso{i}{j}$ estimados a partir de la velocidad comercial de sistemas
        análogos de transporte masivo.
        \item \textit{Simulación comparativa}: Ejecución de \textit{VFTModel} sobre
        el grafo ampliado para calcular la variación neta en $DI$, $T$ y $B(v)$,
        cuantificando la reducción de tiempos de traslado y la redistribución de la
        centralidad de intermediación \parencite{cabrera2026VFTModel}.
    \end{itemize}
\end{description}

\subsection{Fase 4: Síntesis y lineamientos de integración}
\label{subsec:fase4}

Esta fase tiene una ponderación del 10\% y cubre el \textbf{Objetivo Específico~6}.

\begin{description}
    \item[Objetivo de la Fase:] Traducir los resultados cuantitativos de las fases
    anteriores en lineamientos concretos de integración operativa y tecnológica que
    orienten la toma de decisiones en materia de movilidad metropolitana.

    \item[Actividades a Desarrollar:]~
    \begin{itemize}
        \item \textit{Interpretación de resultados}: Lectura crítica de los
        indicadores comparativos para valorar la pertinencia del corredor anillar
        frente al modelo radial vigente, con énfasis en la equidad territorial y la
        reducción de la vulnerabilidad funcional del sistema.
        \item \textit{Formulación de lineamientos}: Desarrollo de recomendaciones
        para la integración modal, tarifaria y operativa del corredor circular,
        orientadas a reducir la dependencia del transporte concesionado en las
        alcaldías del sur y oriente de la \cdmx{} y la \zmvm{}
        \parencite{semovi2020diagnostico}.
    \end{itemize}
\end{description}

%% ── 2-MarcoTeorico/1-Antecedentes.tex
\subsection*{1 Antecedentes}
% !TeX root = ../tesis.tex

\section*{Antecedentes y Marco Teórico}
\label{ch:marco-teorico}

Al inicio del siglo XX, la capital mexicana contaba con escasos tranvías eléctricos
y unos pocos transportes de ruta, pero a medida que avanzó el siglo, la población
también creció, lo que obligó al gobierno federal a invertir en los transportes
colectivos. Este avance se vio frenado con el estallido de la Revolución Mexicana
en 1910, movimiento que no cesó hasta mediados de 1930. Para esos años ya se
habían implementado nuevas rutas de transporte colectivo, como la creación del
naciente tranvía y sus rutas que partían del centro de la Ciudad. Otro avance
notorio fue la promulgación de nuevos transportes de ruta a localidades alejadas,
siendo estos mayormente camiones. Sin embargo, el interés por la expansión de las
áreas urbanas no se dio hasta la época llamada ``El Milagro Mexicano'', la cual
abarca desde los inicios de 1940 hasta mediados de la década de 1970, según
\textcite{molina2018}.

Es a partir del Cardenismo que la reconstrucción de la economía, la infraestructura
y el tejido político se consumó, haciendo necesaria una inversión significativa sobre
las ciudades y su comunicación. En su informe entregado ante las cámaras federales
de 1935, el Ejecutivo reportó una recuperación notoria en las arcas públicas,
destinando buena parte del cálculo presupuestal a las necesidades sociales y de
comunicación:

\begin{quote}
``Es verdaderamente satisfactorio para el Encargado del Poder Ejecutivo de la
Nación dirigirse a sus conciudadanos con motivo del primero de año, enviándoles un
mensaje cordial y lleno de optimismo, pues las condiciones del país nos permiten
afirmar que algunos de los anhelos trascendentales del pueblo podrán realizarse
totalmente en el presente ciclo, e iniciarse la ejecución de otros varios que
complementen la base de un período nuevo de positivo progreso.''
\parencite[pp.~2--4]{biblioteca1935}
\end{quote}

Con este contexto, el urbanismo experimentó un acelerado crecimiento, el mismo
que se empalma con las necesidades de un país que transitaba de los caudillos y lo
rural hacia lo institucional y urbano.

% -----------------------------------------------------------------------
\section{La Hegemonía del Automóvil y la adopción de la automovilidad}
\label{sec:hegemonia-automovil}
% -----------------------------------------------------------------------

El problema de movilidad que actualmente afecta a la \cdmx{} y su periferia no es
un fenómeno coyuntural, sino el resultado de una adopción histórica y cultural de la
automovilidad como paradigma urbano. Este concepto, definido por
\textcite{franco2022} en su tesis sobre el periodo 1903--1933, establece que el
vehículo privado dejó de ser un simple medio de transporte para convertirse en el
eje estructurador de la política pública y el imaginario social.

Esta adopción se consolidó a mediados del siglo XX con la priorización política de
la infraestructura vial sobre otros modos de transporte. El crecimiento vehicular en
el Distrito Federal fue exponencial, pasando de 150,584 unidades en 1955 a
209,904 en 1958 \parencite{perlo2023}, una demanda que el Estado decidió gestionar
no mediante la diversificación modal, sino a través de la expansión de la red de
carreteras urbanas, sellando así la dependencia estructural \DIFdelbegin \DIFdel{al }\DIFdelend \DIFaddbegin \DIFadd{del }\DIFaddend automóvil que
caracteriza la problemática actual.

% -----------------------------------------------------------------------
\section{El Anillo Periférico: Un Proyecto Político Estructural}
\label{sec:anillo-periferico}
% -----------------------------------------------------------------------

La construcción del Anillo Periférico, junto con el Viaducto Miguel Alemán, debe
analizarse no como una simple necesidad de infraestructura, sino como la expresión
material y simbólica de la hegemonía de la automovilidad como política de Estado
durante la regencia de Ernesto P. Uruchurtu (1952--1966) \parencite{franco2022}.
La magnitud de estas obras subraya el giro drástico que la administración de Ruiz
Cortines imprimió a la agenda de gobierno para el Distrito Federal, requiriendo un
funcionario con la energía y la afinidad ideológica de Uruchurtu, cuya trayectoria
previa le había proporcionado un conocimiento directo de la función administrativa
del \textsc{ddf} \parencite[Tomo~1, p.~173]{perlo2023}.

% -----------------------------------------------------------------------
\subsection{Centralización Administrativa y la Creación de los XII Cuarteles}
\label{subsec:xii-cuarteles}
% -----------------------------------------------------------------------

Antes de la ejecución de la mega-obra vial, la regencia de Uruchurtu se enfocó en
la reestructuración administrativa y financiera del Departamento del Distrito Federal
(\textsc{ddf}). Esta etapa inicial (1952--1953) incluyó una reforma a la recaudación
de impuestos que generó un aumento sustancial del presupuesto para obra pública,
lo cual sentó las bases para el posterior desarrollo urbano.

Derivado de este crecimiento poblacional y la necesidad de optimizar la prestación
de servicios, el regente procedió a la división del \textsc{ddf} en doce Cuarteles.
Según la documentación, esta delimitación fue estrictamente administrativa y de
vigilancia, sin peso político-estatal directo, orientada a un control poblacional y
estratégico de la Ciudad de México. No obstante, el alcance de esta estrategia fue
rápidamente superado por la expansión física de la mancha urbana que forjó la
naciente Área Metropolitana de la Ciudad de México.

La importancia de esta división radica en su impacto urbanístico posterior: el diseño
del Anillo Periférico y el Circuito Interior representó la materialización geográfica
de los límites de estos Cuarteles \parencite[Tomo~1, p.~257]{perlo2023}. Esta
correlación entre el límite administrativo de control y la infraestructura vial más
ambiciosa de la época permite afirmar que la creación de ambos circuitos fungió como
una delimitación operativa y simbólica que, si bien fue superada en su función de
control, sigue siendo una referencia fundamental en la división económica y política
de la capital.

\begin{figure}[H]
    \centering
    \figinclude[width=0.82\textwidth]{Figures/Cap1/Uruchurtu1}
    \caption[Designación de los XII Cuarteles en la Ciudad de México]{Designación
    de los XII Cuarteles en la Ciudad de México y crecimiento de la mancha urbana
    (1940--1970).}
    \label{fig:xii-cuarteles}
    \fuentefigura{Perlo Cohen (2023, p.~257).}
\end{figure}

% -----------------------------------------------------------------------
\subsection{La Prioridad Presupuestal: El Sesgo Vial Hegemónico}
\label{subsec:sesgo-vial}
% -----------------------------------------------------------------------

La materialización de la modernidad urbana, bajo la óptica de Uruchurtu, se basó en
la implementación de grandes avenidas y vías rápidas orientadas exclusivamente al
automóvil. Entre las obras más destacadas se encuentran el Viaducto Miguel Alemán,
la Avenida Río Churubusco, el Circuito Interior y el Anillo Periférico. Este último,
que se extendía desde Cuatro Caminos hasta Cuemanco, se concibió como la
principal arteria de comunicación, representando una obra moderna y cosmopolita
destinada a servir al flujo vehicular masivo.

La decisión de priorizar estas obras no fue una casualidad de la planeación, sino un
sesgo presupuestal deliberado. \textcite[Tomo~1, p.~391]{perlo2023} documenta que
solo en 1958, el \textsc{ddf} destinó \$408 millones de pesos a obras en avenidas,
banquetas, pavimentos y pasos a desnivel. Esta inversión superó con creces la suma
destinada a servicios sociales esenciales como mercados, escuelas y centros
asistenciales, evidenciando que la eficiencia del flujo vehicular era el criterio rector
del desarrollo urbano por encima de la equidad social o la integración territorial.

% -----------------------------------------------------------------------
\subsection{El Déficit Estructural y el Rechazo al Metro}
\label{subsec:rechazo-metro}
% -----------------------------------------------------------------------

A pesar de la magnitud de la inversión vial, el \textsc{ddf} afrontaba una crisis de
movilidad imposible de ignorar. Para finales de la década de 1950, la mayoría de la
población carecía de automóvil propio, dependiendo de un sistema de autobuses de
rutas concesionadas caracterizado por ser precario e ineficiente. La respuesta del
gobierno se limitó a la ampliación de rutas y a la promoción de transportes eléctricos
(trolebuses y tranvías), una estrategia insuficiente ante el crecimiento desmedido
de la urbe.

El punto de inflexión que definió el futuro de la movilidad capitalina fue el rechazo
de Uruchurtu a la propuesta de un tren subterráneo ---el actual Sistema de Transporte
Colectivo Metro---, presentada por el ingeniero Bernardo Quintana Arrioja a principios
de la década de 1960. El regente desestimó la idea por considerarla costosa,
invasiva e innecesaria, argumentando la supuesta inviabilidad del proyecto debido al
suelo fangoso de la ciudad. Esta decisión tuvo un impacto trascendental, retrasando
la construcción del Metro hasta 1967 y consolidando el predominio del transporte
superficial y motorizado.

En síntesis, las obras de Uruchurtu carecieron de un programa articulado y basado
en un patrón urbanístico refinado, respondiendo principalmente a las necesidades
inmediatas que el regente logró identificar. El legado del Periférico es una
fragmentación funcional: internamente, delimitó el orden administrativo de los XII
Cuarteles; externamente, propició un crecimiento sin orden que ha dejado a los
territorios periféricos sin conectividad de transporte masivo hasta el día de hoy.

% -----------------------------------------------------------------------
\section{Marco Teórico: Dispersión, Fragmentación y Equidad}
\label{sec:marco-teorico-dispersion}
% -----------------------------------------------------------------------

La consecuencia directa del modelo vial radial fue la expansión desordenada de la
mancha urbana en un proceso denominado \termino{dispersión}. Sin embargo, para
los fines de esta tesis, el concepto más relevante es la \termino{fragmentación
urbana}, ya que denota una ruptura funcional y social que va más allá de la simple
extensión física del territorio.

\textcite{padilla2007} argumenta que la fragmentación se manifiesta en una dislocación
de la cohesión urbana y social, donde las nuevas áreas periféricas carecen de una
adecuada conexión funcional con el centro y entre sí. Esta dislocación se traduce en
una desigualdad de acceso que perpetúa las brechas socioespaciales. Por lo tanto,
el problema a resolver en el Anillo Periférico es la fragmentación territorial resultante
del modelo radial, que impide la movilidad equitativa para las alcaldías del sur.

% -----------------------------------------------------------------------
\subsection{Limitaciones del Modelo Radial}
\label{subsec:limitaciones-radial}
% -----------------------------------------------------------------------

La planificación del transporte público en la \cdmx{} siguió el esquema descrito por
\textcite{molinero2002}, donde las rutas radiales se priorizan para conectar áreas
periféricas con un Distrito Central de Negocios (\textsc{cbd}, por sus siglas en
inglés). Este esquema predominó durante la segunda mitad del siglo XX, cuando la
ciudad experimentó un crecimiento urbano desmedido. Sin embargo, en ciudades
con más de 300,000 habitantes este tipo de rutas comienza a ser ineficiente, ya que
concentra los movimientos en el centro y no considera las necesidades de
desplazamiento entre otras áreas urbanas \parencite{molinero2002}.

% -----------------------------------------------------------------------
\subsection{Transiciones de Fase en Redes Urbanas}
\label{subsec:transiciones-fase}
% -----------------------------------------------------------------------

A la fecha de realización de este estudio (2025), no existe en la \cdmx{} alguna ruta
circular que sirva de soporte para el transporte público masivo colectivo. A pesar de
que la Ciudad cuenta con una vialidad de trazo circular (el Anillo Periférico), \DIFdelbegin \DIFdel{ésta
}\DIFdelend \DIFaddbegin \DIFadd{esta
}\DIFaddend se ha orientado históricamente al automóvil privado. Esta ausencia representa una
\termino{brecha topológica} en la red: la inexistencia de conexiones transversales
impide que el sistema alcance una mayor eficiencia ante el crecimiento continuo de
la demanda.

% -----------------------------------------------------------------------
\subsection{Minimización de Transferencias en Redes Anillo-Radiales}
\label{subsec:minimizacion-transferencias}
% -----------------------------------------------------------------------

Desde la perspectiva de la teoría de redes, \textcite{owais2020} demuestra que la
incorporación de una geometría anillar a un sistema predominantemente radial reduce
el número de transbordos necesarios para conectar zonas periféricas entre sí. Esta
minimización de transferencias no solo mejora los tiempos de traslado, sino que
incrementa la percepción de calidad del servicio por parte del usuario, al reducir la
fricción del viaje multimodal que caracteriza actualmente la movilidad en las alcaldías
del sur de la \cdmx{}.

% -----------------------------------------------------------------------
\section{Teoría Latinoamericana}
\label{sec:teoria-latinoamericana}
% -----------------------------------------------------------------------

Para enmarcar el caso de la \cdmx{} en un contexto de crítica urbanística regional,
es fundamental recurrir a la teoría sobre la ciudad en América Latina. La obra
compilatoria de \textcite{padilla2007} proporciona el marco conceptual para entender
que la crisis de movilidad en la periferia es síntoma de un agotamiento del modelo
de planeación estatal y la adopción de enfoques que favorecieron el capital
inmobiliario y la privatización del espacio. La incapacidad de la planeación urbana
para gestionar la expansión territorial generada por los grandes proyectos viales,
como el Periférico, es el caldo de cultivo para la emergencia de la ciudad neoliberal.

La propuesta de un corredor circular es, en este contexto teórico, un esfuerzo por
revertir el impacto de la automovilidad hegemónica y la planificación fragmentada.
Busca pasar de un enfoque que prioriza la velocidad del flujo ---propio de la era
Uruchurtu--- a uno que prioriza la integración equitativa y la eficiencia operativa,
alineado con las demandas de un nuevo modelo de ciudad que ya no puede ser
definido únicamente por sus arterias motorizadas \parencite{franco2022}.

% -----------------------------------------------------------------------
\section{Casos Similares en Otros Países}
\label{sec:casos-similares}
% -----------------------------------------------------------------------

La implementación de corredores circulares de transporte masivo no es una idea
nueva a nivel global. Su efectividad ha sido documentada en diversas megaciudades
que enfrentaban problemas estructurales similares a los de la \cdmx{}: redes
radiales saturadas, periferias desconectadas y tiempos de traslado crecientes.

El caso más representativo es la \textbf{Gran Línea Circular del Metro de Moscú},
inaugurada en 2023 con 70 kilómetros de longitud. Esta línea fue concebida
explícitamente para descomprimir las líneas radiales del centro y articular la
movilidad entre zonas periféricas sin necesidad de pasar por las estaciones
centrales más congestionadas \parencite{efe2023}. Su puesta en operación
demostró cómo la adopción de geometrías anillares puede cambiar cualitativamente
la eficiencia de todo el sistema de transporte de una capital.

Un segundo referente es el \textbf{\textit{Suburban Rail Loop} de Melbourne},
Australia, actualmente en planeación y construcción. Este proyecto parte de un
diagnóstico análogo al de la \cdmx{}: las alcaldías del llamado \textit{Middle-Ring}
albergan hospitales, universidades y centros de empleo de alta demanda, pero
carecen de una conectividad lateral robusta que los articule entre sí
\parencite{suburbanrailloop2021}. El proyecto busca transformar una serie de
suburbios desconectados en una red policéntrica de precintos de innovación y
empleo, mediante la inserción de un corredor ferroviario de arco periférico.

Estos referentes internacionales validan la hipótesis central de esta investigación:
la inserción de una ruta anillar en la sección sur del Periférico no es una solución
experimental, sino una estrategia probada y replicable en el contexto específico de
la \cdmx{} y su \zmvm{}.

%% ── 3-Conceptos-Indicadores/3-1-Definiciones.tex
\subsection*{3 1 Definiciones}
% !TeX root = ../tesis.tex

\section*{Definiciones de Conceptos Fundamentales}
\label{ch:definiciones}

% =============================================================================
%  SECCIÓN 3.2 — Definición de Conceptos Fundamentales
% =============================================================================

\section{Definición de Conceptos Fundamentales}
\label{sec:conceptos-fundamentales}

El análisis de redes de transporte público desde un enfoque computacional y
matemático exige la precisión conceptual de los términos que articulan tanto el
diagnóstico del sistema como el diseño de intervenciones. En esta sección se
definen los cuatro conceptos estructurales del presente trabajo ---vulnerabilidad
urbana, planeación estratégica, eficiencia urbana e interconectividad entre
redes--- con base en la literatura especializada en teoría de grafos, redes complejas
y análisis de sistemas de movilidad. Estas definiciones no solo delimitan el alcance
teórico de la investigación, sino que también establecen los fundamentos formales
sobre los cuales se construirán los modelos computacionales, los algoritmos de
análisis y los indicadores de desempeño de la red de transporte público del área
metropolitana estudiada.

% -----------------------------------------------------------------------
\subsection{Vulnerabilidad Urbana}
\label{subsec:vulnerabilidad-urbana}
% -----------------------------------------------------------------------

\subsubsection{Concepto y Marco Teórico desde las Redes Complejas}
\label{subsubsec:vulnerabilidad-redes-complejas}

El estudio de la vulnerabilidad en sistemas de transporte urbano ha experimentado
una transición significativa: del análisis geográfico y estadístico tradicional hacia
un enfoque topológico fundamentado en la teoría de redes complejas. Este cambio
de paradigma es relevante para la presente investigación, pues permite representar
la red de transporte como un grafo matemático $\grafo = \{\nodos, \aristas\}$,
donde los nodos $\nodos$ corresponden a las estaciones y paradas, y las aristas
$\aristas$ a las rutas o conexiones entre ellas, de modo que las propiedades
estructurales globales del sistema quedan formalmente accesibles al análisis
algorítmico.

\textcite{loterovélez2014} señalan que una de las propiedades de mayor relevancia
práctica de las redes complejas es la capacidad del sistema para mantener algunas
de sus funciones cuando ocurren fallas, errores o ataques dirigidos a sus nodos o
vínculos; esta propiedad ha sido denominada de manera indistinta como
\termino{robustez}, \termino{resiliencia} o \termino{vulnerabilidad} por distintos
autores. En el contexto del transporte urbano, las fallas pueden corresponder a
cierres de estaciones por mantenimiento, incidentes operativos o desastres
naturales que impliquen, en términos de grafos, la eliminación de uno o varios nodos
o aristas de la red, afectando su conectividad y los flujos que la atraviesan.

No existe, sin embargo, un consenso único en la definición de vulnerabilidad. Según
el análisis de \textcite{loterovélez2014}, diversas corrientes han propuesto
conceptualizaciones complementarias: Holmgren (2006) la define como la
sensibilidad del sistema de infraestructura a las amenazas y perturbaciones que
puedan presentarse; Nagurney y Qiang (2011) la conciben como la cuantificación
del impacto generado por la remoción de un componente sobre el desempeño de la
red; mientras que Ouyang et al. (2009) la describen operativamente como la
disminución en la eficiencia de la red tras un ataque o falla. Para los propósitos del
presente trabajo, se adopta esta última perspectiva operativa, dado que permite
traducir el concepto de vulnerabilidad en métricas computables sobre el grafo.

\subsubsection{Fragilidad Estructural y Análisis Topológico}
\label{subsubsec:fragilidad-estructural}

La medición de la vulnerabilidad descansa en el análisis de las propiedades
topológicas del grafo que representa la red. \textcite{loterovélez2014} indican que
la forma más común de evaluar la robustez de una red es observar su comportamiento
después de la eliminación de uno o varios de sus elementos ---sean nodos, aristas o
una combinación de ambos---, donde dicha remoción puede ser aleatoria (simulando
una falla fortuita) o dirigida a los elementos más importantes (simulando un ataque).

Desde esta perspectiva, la vulnerabilidad estructural se manifiesta en el cambio de
las \termino{geodésicas} del grafo tras la perturbación. Las geodésicas, entendidas
como los caminos de menor longitud (\textit{shortest paths}) entre pares de nodos
origen-destino, constituyen el indicador central: si la eliminación de un nodo
incrementa significativamente la longitud media de dichos caminos o produce la
desconexión de componentes del grafo, ese nodo es clasificado como un punto de
fallo crítico. Albert, Jeong y Barabási (2000), citados por \textcite{loterovélez2014},
demostraron que las redes con distribución libre de escala ---caracterizadas por la
existencia de unos pocos nodos con un elevado número de conexiones, denominados
\termino{hubs}--- son robustas ante fallas aleatorias, pero altamente vulnerables ante
ataques dirigidos a dichos nodos centrales, patrón que resulta especialmente
relevante para las redes de transporte metropolitanas con nodos de alta centralidad.

Las medidas de centralidad constituyen, en consecuencia, los instrumentos
fundamentales para la identificación de nodos críticos. Según \textcite{loterovélez2014},
entre las más relevantes se encuentran: la \termino{centralidad de grado nodal}
(número de conexiones de un nodo), la \termino{centralidad de intermediación}
(proporción de caminos más cortos entre pares de nodos que pasan por el nodo en
cuestión), la \termino{centralidad de cercanía} (inverso de la suma de distancias del
nodo a los demás) y los \termino{coeficientes de cohesión} (conectividad entre los
vecinos de un nodo). En el modelo computacional que se implementa en este trabajo,
estos indicadores se calculan algorítmicamente sobre el grafo de la red para
identificar las estaciones cuya pérdida produciría el mayor impacto sistémico.

Más allá de la vulnerabilidad estructural, que considera exclusivamente la topología
del grafo, Ouyang et al. (2009), referenciados por \textcite{loterovélez2014},
distinguen una \termino{vulnerabilidad funcional} que incorpora los niveles reales de
funcionamiento de la red: frecuencias de servicio, tiempos de espera y capacidades
de carga. Esta distinción resulta crucial en el análisis de transporte público, donde la
pérdida de un nodo de alto flujo tiene consecuencias que trascienden la simple
alteración topológica y se traducen en redistribuciones de demanda sobre rutas
alternativas, potencialmente colapsando otros puntos de la red por sobrecarga.

Para los sistemas de transporte urbano, \textcite{loterovélez2014} afirman que una
aproximación plausible es el uso de herramientas conceptuales y analíticas de la
ciencia de la complejidad, entre ellas las redes complejas, que permiten mapear de
manera directa los elementos del sistema en nodos y vínculos. Esto fundamenta la
elección metodológica del presente trabajo: al modelar la red de transporte como un
grafo ponderado y dirigido, es posible simular computacionalmente escenarios de
falla progresiva y evaluar el comportamiento global del sistema bajo distintos niveles
de perturbación, obteniendo así un perfil de vulnerabilidad completo y cuantificado.


\begin{figure}[H]
    \centering
   \figinclude[width=0.90\textwidth]{Figures/Cap3/TopologiaNodos.png}
   \caption[Topologías de red y vulnerabilidad]{Cuatro topologías de red para
   50 nodos: regular, mundo pequeño, aleatoria y libre de escala.}
   \label{fig:topologias-red}
   \fuentefigura{\textcite{loterovélez2014}, p.~3.}
 \end{figure}

 % -----------------------------------------------------------------------
\subsection{Planeación Estratégica}
\label{subsec:planeacion-estrategica}
% -----------------------------------------------------------------------

\subsubsection{Caracterización Dinámica de la Oferta de Transporte}
\label{subsubsec:caracterizacion-dinamica}

La planeación estratégica del transporte público ha evolucionado de un proceso
esencialmente reactivo ---basado en la observación de demandas insatisfechas---
hacia un modelo predictivo y basado en datos que emplea el modelado computacional
de la red para anticipar los efectos de sus transformaciones. En este marco, el grafo
matemático que representa la red deja de ser un producto del análisis para
convertirse en el instrumento central de la planeación.

\textcite{fortin2016innovative} argumentan que, hasta la fecha de su investigación, no existía
una herramienta de análisis capaz de caracterizar sistemáticamente una red de
transporte y observar sus cambios a lo largo del tiempo de manera estructurada y
automatizada. Su propuesta aborda esta brecha mediante un método orientado a
grafos que importa datos en formato \gtfs{} (\textit{General Transit Feed
Specification}) hacia una base de datos de grafos, desde la cual se pueden calcular
algoritmos de caminos mínimos y derivar indicadores de desempeño de manera
automatizada. Este enfoque es el que adopta el presente trabajo: los datos \gtfs{}
de la red estudiada se procesan computacionalmente para construir un grafo dinámico
que refleja no solo la geometría de la red, sino también sus niveles de servicio,
frecuencias y capacidades operativas.

\textcite{moreno2023} refuerzan este enfoque al plantear que una metodología
robusta para la planeación requiere de una matriz origen-destino (\textsc{od}), una
red georreferenciada y los algoritmos necesarios para calcular la asignación de flujos
vehiculares. Señalan que el lenguaje Python y su ecosistema de librerías constituyen
una plataforma de reconocida eficacia para este propósito, dado que permite el
desarrollo modular del software de análisis, otorgando flexibilidad para actualizar los
modelos y extender su alcance según las necesidades de la investigación. Esta
premisa orienta directamente las decisiones de implementación tecnológica descritas
en los capítulos metodológicos posteriores.

\subsubsection{Toma de Decisiones Basada en Algoritmos y Comparación de Redes}
\label{subsubsec:comparacion-redes}

Una de las contribuciones más relevantes de la planeación estratégica basada en
grafos es la posibilidad de comparar cuantitativamente diferentes configuraciones de
red antes de ejecutar cambios físicos en la infraestructura.
\textcite{fortin2016innovative} proponen explícitamente la \termino{comparación de redes}
(\textit{Network Comparison}) como herramienta central de la planeación: al calcular
indicadores derivados de caminos mínimos sobre el grafo actual de la red y sobre un
grafo hipotético que incorpora modificaciones ---como la adición de nuevas líneas o
la reestructuración de rutas existentes---, es posible evaluar cuantitativamente el
impacto de cada escenario sobre el desempeño del sistema antes de tomar
decisiones de inversión.

Este principio resulta directamente aplicable al caso de estudio de la presente
investigación. Al modelar la red de transporte público como un grafo y simular sobre
él la incorporación del corredor anillar propuesto, los algoritmos de caminos mínimos
---como el algoritmo de \dijkstra{} o sus variantes para grafos ponderados--- permiten
cuantificar la reducción en tiempos de viaje, la mejora en la conectividad entre zonas
periféricas y la eliminación de cuellos de botella en los nodos de alta centralidad.
Así, la planeación abandona el terreno de la intuición para situarse en el de la
demostración algorítmica.

\textcite{moreno2023} complementan esta perspectiva al destacar que el desarrollo
modular del software de asignación permite adaptar el modelo a redes bimodales que
integren distintos modos de transporte ---como los sistemas carretera-ferrocarril o
superficie-metro--- y sus posibles intercambios a través de estaciones y puntos de
transferencia, lo que amplía significativamente el alcance analítico de la herramienta
más allá del análisis de una red monomodal.

\subsubsection{Enfoque en la Sostenibilidad y el Cambio Modal}
\label{subsubsec:sostenibilidad-cambio-modal}

La planeación estratégica del transporte público no puede disociarse de su dimensión
territorial y ambiental. \textcite{ugur2025} señalan que la integración de los
diferentes modos de transporte público ---lograda mediante la evaluación sistemática
de la conectividad entre líneas y la creación de oportunidades efectivas de
transferencia--- incrementa la cobertura espacial de la red, mejora la accesibilidad
urbana, promueve el uso del transporte colectivo y contribuye a la reducción de
emisiones, al tiempo que garantiza la equidad social en la movilidad y ofrece opciones
de viaje más económicas. Desde esta perspectiva, la planeación estratégica tiene
como objetivo último facilitar el cambio modal del vehículo privado hacia el transporte
masivo, para lo cual la optimización de la red ---modelada como grafo--- es condición
necesaria.

\textcite{moreno2023} enmarcan además este esfuerzo en el contexto de la soberanía
tecnológica del Estado, argumentando que el desarrollo de software propio de
asignación de tráfico, basado en herramientas de código abierto, reduce la
dependencia de licencias comerciales costosas y permite que las instituciones
públicas acumulen capacidades analíticas autónomas para la planeación del
transporte. Esta consideración es relevante para el presente trabajo, que adopta
deliberadamente una arquitectura tecnológica basada en Python y librerías de
acceso libre.

\begin{figure}[H]
   \centering
   \figinclude[width=0.85\textwidth]{Figures/Cap3/DiagramCompleteDataSetGTFS.png}
   \caption[Estructura relacional de archivos \gtfs{}]{Diagrama de la estructura
   relacional de los archivos que componen el estándar \gtfs{}.}
   \label{fig:gtfs-estructura}
   \fuentefigura{\textcite{fortin2016innovative}, p.~21.}
 \end{figure}

% -----------------------------------------------------------------------
\subsection{Eficiencia Urbana}
\label{subsec:eficiencia-urbana}
% -----------------------------------------------------------------------

\subsubsection{Eficiencia Topológica y Caminos Mínimos}
\label{subsubsec:eficiencia-topologica}

La eficiencia de una red de transporte se entiende, en su dimensión topológica, como
la capacidad del sistema para conectar cualquier par de puntos origen-destino a
través del camino de menor costo posible ---medido en términos de tiempo de viaje,
número de transbordos o distancia recorrida---. Esta propiedad es directamente
computable sobre el grafo de la red mediante algoritmos de caminos mínimos, lo que
la convierte en una de las métricas más operativas del análisis propuesto en este
trabajo.

\textcite{fortin2016innovative} plantean que los indicadores derivados de la teoría de grafos
---como el \termino{índice alfa} (grado de ciclicidad), el \termino{índice gamma}
(grado de conectividad) y el \termino{índice de superposición de líneas}---
caracterizan la eficiencia del diseño de la red de manera directamente relacionada
con su topología. Sin embargo, estos autores reconocen que tales indicadores no
consideran las características únicas del transporte público, como la frecuencia de
servicio, la posibilidad de transferencias entre líneas o la variabilidad del servicio a
lo largo del día; por ello proponen complementarlos con indicadores dinámicos
derivados del modelo de grafo expandido en el tiempo construido con datos \gtfs{}.
En este modelo, la eficiencia se mide a través del camino mínimo ponderado en la
dimensión temporal, que incorpora tiempos de espera y frecuencias reales de cada
línea.

Mishra et al. (2012), citados por \textcite{mishra2012}, introdujeron en este
contexto el concepto de \termino{índice de conectividad} de un nodo como una
medida de desempeño que va más allá de la centralidad tradicional. Este índice se
define como la suma de las potencias de conexión de todas las líneas que pasan por
un nodo, considerando características operativas como la velocidad, la capacidad, la
frecuencia, el tiempo de encabezamiento (\textit{headway}) y la distancia desde el
punto de origen. Así, la eficiencia topológica y la eficiencia operativa se integran en
una sola métrica nodal computable, que será adoptada como referencia en la
construcción de los indicadores del presente trabajo.

\subsubsection{Centralidad de Intermediación y Congestión de Rutas}
\label{subsubsec:centralidad-intermediacion}

La distribución de los flujos de tráfico a lo largo de la red determina, en gran medida,
su eficiencia operativa. Un sistema topológicamente bien diseñado puede resultar
ineficiente en la práctica si los flujos de demanda se concentran
desproporcionadamente sobre un subconjunto reducido de nodos o aristas, generando
congestión y tiempos de viaje superiores a los predichos por el modelo teórico.

\textcite{batista2015} proponen abordar este problema mediante una variación
ponderada de la \termino{centralidad de intermediación} (\textit{betweenness
centrality}), en la que el grafo de la red vial es ponderado no solo por la longitud de
las aristas, sino por el número de rutas generadas por la demanda real de tráfico,
representada mediante una matriz origen-destino (\textsc{od}). En su marco teórico,
un nodo o arista de alta centralidad de intermediación es aquel que aparece con mayor
frecuencia en los caminos mínimos entre los pares origen-destino del sistema; los
nodos con esta propiedad concentran los flujos de manera natural y son, por tanto,
los primeros en colapsar ante un incremento de demanda o una perturbación. La
eficiencia urbana aumenta cuando el sistema logra redistribuir estos flujos hacia rutas
alternativas, reduciendo la dependencia de los nodos centrales y homogeneizando
la carga sobre la red.

Esta perspectiva es fundamental para el análisis de la red metropolitana objeto del
presente estudio: la identificación computacional de los arcos y nodos con mayor
centralidad de intermediación ponderada por la demanda permite señalar aquellos
segmentos del sistema donde las inversiones en infraestructura o las modificaciones
de rutas producirían las mayores ganancias en eficiencia global.
 \subsubsection{Accesibilidad Espacial como Dimensión de la Eficiencia}
\label{subsubsec:accesibilidad-espacial}

La eficiencia urbana no se agota en la optimización interna de los flujos de la red;
incluye también la dimensión espacial de la \termino{accesibilidad}, es decir, la
capacidad del sistema de transporte para garantizar que cualquier usuario,
independientemente de su ubicación en el territorio, pueda alcanzar una cantidad
mínima de destinos dentro de umbrales de tiempo razonables.

\textcite{fortin2016innovative} señalan que los indicadores de accesibilidad más comunes se
calculan definiendo un radio de cobertura (por ejemplo, 500 metros como distancia
caminable aceptable) alrededor de las paradas de la red, y evaluando la proporción
de población o de empleos que quedan dentro de dicha cobertura. Sin embargo, estos
autores reconocen que este enfoque estático no incorpora la demanda real de los
usuarios ---representada por sus pares origen-destino específicos--- ni la frecuencia
del servicio, lo que limita su capacidad para reflejar la experiencia real de
accesibilidad. Proponen por ello indicadores dinámicos que, construidos sobre el
grafo ponderado en el tiempo, permiten calcular la accesibilidad real desde cualquier
nodo de la red hacia todos los demás, considerando los horarios y frecuencias del
servicio. Esta \termino{accesibilidad dinámica}, calculada algorítmicamente sobre el
modelo de grafo, constituye una medida integral de la eficiencia urbana que se
incorporará al análisis computacional de la presente investigación.

\begin{figure}[H]
   \centering
   \figinclude[width=0.85\textwidth]{Figures/Cap3/TablaIndicatorsAdaptedGraphTeory.png}
   \DIFdelbeginFL %DIFDELCMD < \caption[Indices de Grafos \gtfs{}]{%%%
\DIFdelendFL \DIFaddbeginFL \caption[Índices de Grafos \gtfs{}]{\DIFaddendFL Tabla de indicadores de grafos \gtfs{}.}
   \label{fig:indices-grafos}
   \fuentefigura{\textcite{fortin2016innovative}, p.~23.}
 \end{figure}

% -----------------------------------------------------------------------
\subsection{Interconectividad entre Redes}
\label{subsec:interconectividad}
% -----------------------------------------------------------------------

\subsubsection{Integración Multimodal y Conectividad a Nivel de Nodo}
\label{subsubsec:integracion-multimodal}

La interconectividad entre redes es el grado en que distintos modos y líneas de
transporte público ---metro, autobús, trolebús, \brt{}, ciclovías, entre otros---
operan de manera coordinada como un sistema único y fluido, minimizando las
fricciones de la transferencia y maximizando la cobertura territorial efectiva del
sistema. En la representación por grafos, la interconectividad se modela
introduciendo \termino{aristas de transferencia} que conectan nodos de diferentes
modos o líneas dentro de un mismo grafo multimodal, con pesos que representan el
tiempo de caminata entre andenes, el tiempo de espera esperado y, en su caso, el
costo tarifario adicional del transbordo.

\textcite{ugur2025} afirman que la integración juega un papel crucial para garantizar
que los diferentes modos de transporte público operen de manera armoniosa,
facilitando las transferencias y permitiendo a los pasajeros desplazarse con facilidad
entre destinos. En un sistema bien integrado, la fortaleza de la conexión entre líneas
no solo amplía el área accesible para el usuario, sino que también reduce las
dificultades experimentadas durante el viaje, lo que influye positivamente en la
demanda de transporte público y contribuye a la sustitución del vehículo privado.
Su investigación sobre la línea M7 del metro de Estambul demuestra que la
evaluación sistemática de la conectividad en cada estación ---mediante un
\termino{Índice de Conectividad Simplificado} (\textsc{sci}, por sus siglas en
inglés) que incorpora variables de velocidad, capacidad de pasajeros, frecuencia,
horas de operación y distancia al origen--- revela patrones de integración desigual
que orientan directamente las decisiones de inversión y mejora operativa.

\subsubsection{Medición de la Conectividad de Parada y Penalización por Transbordo}
\label{subsubsec:penalizacion-transbordo}

La evaluación cuantitativa de la interconectividad requiere un índice que capture la
calidad de la conexión en cada nodo del grafo multimodal, más allá del simple conteo
de líneas que lo sirven. \textcite{mishra2012} proponen un marco de análisis basado
en la teoría de grafos que evalúa la conectividad a múltiples niveles del sistema de
transporte: el nivel de nodo (parada individual), el nivel de línea (ruta) y el nivel de
centro de transferencia (\textit{transfer center}). En este marco, el índice de
conectividad de un nodo se define como el producto de las potencias de conexión de
todas las líneas que lo atraviesan, donde la potencia de conexión de cada línea
depende de sus características operativas ---velocidad, distancia, frecuencia, tiempo
de encabezamiento, capacidad, aceleración y desaceleración--- así como de la
direccionalidad de la ruta.

Este enfoque supera las limitaciones de la centralidad de grado simple, que
únicamente cuenta el número de conexiones directas de un nodo sin considerar la
calidad ni la frecuencia del servicio ofrecido. \textcite{mishra2012} subrayan que
medir la conectividad puede resultar una tarea compleja y esquiva, pues es necesario
considerar tarifas, horarios, capacidad, frecuencia y otras características del sistema
en su conjunto; por ello, la evaluación de la conectividad del transporte requiere un
enfoque sistemático que integre los múltiples parámetros involucrados en la
prestación del servicio real. La diferencia entre un sistema integrado y uno
fragmentado se traduce directamente en la magnitud de la \termino{penalización por
transbordo} percibida por el usuario: cuando la conectividad en los nodos de
transferencia es alta, esta penalización se reduce, haciendo el sistema más atractivo
respecto al automóvil privado.

\subsubsection{El Grafo Multimodal como Marco Formal de Evaluación}
\label{subsubsec:grafo-multimodal}

La representación matemática de la interconectividad entre redes exige la
construcción de un \termino{grafo multimodal} en el que coexistan, sobre una misma
estructura topológica, los nodos y aristas de los distintos modos de transporte,
vinculados mediante aristas artificiales de transferencia debidamente ponderadas.
\textcite{mishra2012} señalan que, a diferencia de un enlace en una red vial ---que
es un segmento físico que conecta un nodo con otro---, un enlace en una red de
transporte público multimodal forma parte de una línea que sirve una secuencia de
paradas; dado que diferentes líneas pueden servir una misma parada, pueden existir
múltiples aristas entre los mismos nodos del grafo, lo que exige representaciones
formales más complejas que las de una red monomodal.

En el modelo computacional implementado en este trabajo, la red de transporte
metropolitana se representa como un grafo dirigido y ponderado
$\grafo = \{\nodos,\, \aristas,\, W\}$, donde $\nodos$ comprende el conjunto de todas
las paradas de todos los modos, $\aristas$ incluye tanto las aristas de trayecto
(intra-línea) como las aristas de transferencia (inter-modo), y $W$ es la función de
peso que asigna a cada arista un costo compuesto ---tiempo de viaje, tiempo de
espera, tiempo de caminata--- calculado a partir de los datos \gtfs{} de la red. Sobre
este grafo unificado, los algoritmos de caminos mínimos permiten evaluar la
interconectividad del sistema de manera integral: una transferencia deficiente se
traduce en un peso elevado en la arista de transbordo correspondiente, lo que el
algoritmo reflejará automáticamente como una ruta subóptima, señalando así los
nodos de transferencia que requieren intervención prioritaria.

% Sugerencia: insertar aquí el mapa de resultados SCI de Uğur y Duran (2025)
% sobre la línea M7 de Estambul, o el diagrama de grafo multimodal de Mishra
% et al. (s.f.) para la región Washington-Baltimore, una vez disponibles.
\begin{figure}[H]
  \centering
  \figinclude[width=0.88\textwidth]{Figures/Cap3/M7Estambul.png}
  \caption[Índice de Conectividad Simplificado en la línea M7 de Estambul]{
  Distribución del \textsc{sci} a lo largo de las 17 estaciones de la línea M7
  del metro de Estambul (buena, media y débil conectividad).}
  \label{fig:sci-estambul}
  \fuentefigura{\textcite{ugur2025}, p.~XX.}
 \end{figure}

%% ── 3-Conceptos-Indicadores/3-2-Grafos-Matematicas.tex
\subsection*{3 2 Grafos Matematicas}
% !TeX root = ../tesis.tex
%%%%%%%%%%%%%%%%%%%%%%%%%%%%%%%%%%%%%%%%%%%%%%%%%%%%%%%%%%%%%%%%%%%
% SECCIÓN: FUNDAMENTOS DE LA TEORÍA DE GRAFOS E IMPLEMENTACIÓN
%%%%%%%%%%%%%%%%%%%%%%%%%%%%%%%%%%%%%%%%%%%%%%%%%%%%%%%%%%%%%%%%%%%
\section{Grafos Matemáticos y su Implementación Computacional}
\label{sec:grafos-matematicos}

\subsection{Origen y Fundamentos: La Abstracción de Euler}
\label{subsec:euler}
    El estudio formal de la teoría de grafos se remonta a 1736, cuando Leonhard Euler publicó su resolución al
    problema de los siete puentes de la ciudad de Königsberg. El núcleo de ese aporte residió en la capacidad de
    abstraer la geografía física de la ciudad ---dividida por el río Pregel en cuatro zonas terrestres--- hacia un
    esquema puramente relacional: las zonas se representaron como puntos (vértices) y los puentes como líneas
    (aristas), de modo que el problema se resolviera no mediante la geometría física, sino a través de la
    configuración topológica del sistema, tal como se aprecia en la Figura~\ref{fig:königsberg}
    \parencite{garcia2005, braicovich2009}.

\begin{figure}[H]
    \centering
    \figinclude[width=0.70\textwidth]{Figures/Cap3/MapaKonisnberg.png}
    \caption{Mapa histórico de Königsberg y su representación como grafo}
    \label{fig:königsberg}
    \fuentefigura{Figura extraída de \textcite{garcia2005}.}
\end{figure}

\subsection{Grafos de Euler y la Condición de Transitabilidad}
\label{subsec:grafos-euler}
    Derivado del trabajo fundacional de Euler, se establece que un grafo es euleriano si admite un circuito que
    recorra todas sus aristas exactamente una vez, regresando al vértice de origen. La condición necesaria y
    suficiente para que un grafo conexo admita este recorrido es que todos sus vértices posean un grado par; si el
    grafo posee exactamente dos vértices de grado impar, admite en cambio un camino de Euler abierto, que debe
    comenzar en uno de esos vértices y terminar en el otro. Ambas configuraciones se ilustran comparativamente en
    la Figura~\ref{fig:diagramas-aristas} \parencite{garcia2005, braicovich2009}. Esta distinción resulta
    fundamental en problemas de transporte y logística como el del ``cartero chino'', donde se busca minimizar el
    costo de un recorrido que cubra toda una infraestructura vial \parencite{allauca2023}.

\begin{figure}[H]
    \centering
    \begin{subfigure}[b]{0.47\textwidth}
        \centering
        \figinclude[width=\textwidth]{Figures/Cap3/Diagrama1Aristas.png}
        \caption{Representación esquemática de un grafo}
        \label{fig:DiagramaAristas1}
    \end{subfigure}
    \hfill
    \begin{subfigure}[b]{0.47\textwidth}
        \centering
        \figinclude[width=\textwidth]{Figures/Cap3/Diagrama2Aristas.png}
        \caption{Nodos, aristas y condiciones de grado}
        \label{fig:DiagramaAristas2}
    \end{subfigure}
    \caption{Elementos constitutivos de un grafo y condiciones de transitabilidad euleriana}
    \label{fig:diagramas-aristas}
    \fuentefigura{Figuras extraídas de \textcite{garcia2005}.}
\end{figure}

\subsection{Grafos Computacionales: Eficiencia mediante Mapas de Adyacencia}
\label{subsec:grafos-computacionales}
    En el ámbito de las ciencias de la computación, existe una distinción operativa entre el grafo como ente
    matemático abstracto y su implementación en sistemas de datos. \textcite{valiente2022} sostiene que, si bien
    las representaciones tradicionales como las matrices de adyacencia ofrecen un acceso rápido, suelen ser
    ineficientes en el uso de memoria para grafos dispersos. Como alternativa superior, se destacan los
    \textbf{Mapas de Adyacencia (\termino{Adjacency Maps})}, estructuras que particionan nodos y aristas mediante
    diccionarios indexados y permiten verificar la adyacencia entre dos nodos en un tiempo esperado constante
    $O(1)$, optimizando así el rendimiento de los algoritmos en redes de gran escala, tal como se esquematiza en
    la Figura~\ref{fig:PresentacionCaminos}.

\begin{figure}[H]
    \centering
    \figinclude[width=0.75\textwidth]{Figures/Cap3/RepresentacionCaminos.png}
    \caption{Representación matemática de caminos sobre una matriz de adyacencia}
    \label{fig:PresentacionCaminos}
    \fuentefigura{Figura extraída de \textcite{braicovich2009}.}
\end{figure}

\subsection{Optimización de Redes: El Cálculo de Caminos Mínimos}
\label{subsec:caminos-minimos}
    La aplicación de la teoría de grafos en la optimización logística permite resolver desafíos complejos de
    enrutamiento. \textcite{allauca2023} expone que las redes de transporte pueden modelarse como grafos dirigidos
    valorados, donde los nodos representan estaciones y las aristas son rutas ponderadas por variables como tiempo
    y costo. Bajo este enfoque, se destacan algoritmos fundamentales como el de \dijkstra{} para hallar el camino
    más corto desde un nodo origen, y el de \floyd{} para calcular distancias mínimas entre todos los pares de
    nodos, tal como se ilustra en la Figura~\ref{fig:CaminosMinimo}. Estudios de caso demuestran que la aplicación
    de estas técnicas permite reducciones estimadas del 8\% en tiempos de entrega y del 5\% en costos operativos
    \parencite{allauca2023}.

\begin{figure}[H]
    \centering
    \figinclude[width=0.75\textwidth]{Figures/Cap3/CaminosMinimos.png}
    \caption{Representación esquemática de caminos mínimos en un problema teórico}
    \label{fig:CaminosMinimo}
    \fuentefigura{Figura extraída de \textcite{braicovich2009}.}
\end{figure}

\subsection{Integración de Grafos con Sistemas \gis{} y Datos \gtfs{}}
\label{subsec:grafos-sig-gtfs}
    La convergencia entre la teoría de grafos y los Sistemas de Información Geográfica (\gis{}) permite un
    análisis avanzado de la conectividad urbana. \textcite{fortin2016innovative} y \textcite{adas2013} proponen
    el uso del estándar \textbf{Google Transit Data (\gtfs{})} como fuente primaria para construir grafos
    dinámicos que integran la topología de la red con su ubicación georreferenciada. Esta convergencia permite
    generar indicadores clave de desempeño, como el \textbf{Detour Factor ($DI$)}, que cuantifica la eficiencia
    geométrica del trazado comparando la distancia real en la red con\DIFdelbegin \DIFdel{tra}\DIFdelend  la distancia euclidiana; asimismo, el
    uso de modelos de grafos expandidos en el tiempo facilita la evaluación de la accesibilidad dinámica
    considerando frecuencias reales y penalizaciones por transferencia \parencite{fortin2016innovative, adas2013}.

\subsection{Implementaciones Tecnológicas y Algoritmos de Búsqueda}
\label{subsec:implementaciones-grafos}
    En el entorno industrial moderno, las teorías de búsqueda en grafos se implementan a través de motores
    analíticos especializados que procesan datos masivos. \textcite{wang2025} destaca que algoritmos como la
    Búsqueda en Anchura (\termino{BFS}), la Búsqueda en Profundidad (\termino{DFS}) y el algoritmo \termino{A*}
    constituyen el núcleo de las soluciones de navegación actuales. Un ejemplo de esta tendencia es la plataforma
    \textbf{PuppyGraph}, que ofrece una arquitectura de consultas analíticas directamente sobre almacenes de datos
    relacionales, sin procesos de extracción intermedios, y utiliza algoritmos distribuidos para gestionar grafos
    de millones de nodos en aplicaciones de seguridad y logística.

%% ── 3-Conceptos-Indicadores/3-3-Indicadores-Formula.tex
\subsection*{3 3 Indicadores Formula}
% !TeX root = ../tesis.tex

% =============================================================================
%  SECCIÓN 3.3 — Indicadores de Análisis: Definición, Formulación y Aplicación
%  Archivo: 3-Conceptos-Indicadores/3-3-Indicadores-Formula.tex
% =============================================================================

\section{Indicadores de Análisis: Definición, Formulación y Aplicación al
         Caso de Estudio}
\label{sec:indicadores-formulas}

El análisis computacional de la red de transporte público de la \cdmx{} y su zona
metropolitana se articula a través de ocho indicadores cuantitativos organizados en
cuatro fases de construcción progresiva. Esta jerarquía no es arbitraria: cada
indicador depende de los insumos producidos por los anteriores, de modo que la
arquitectura metodológica constituye en sí misma un modelo de grafo que crece en
complejidad desde el análisis espacial puro hasta las pruebas de estrés topológico.
A continuación se presenta\DIFdelbegin \DIFdel{n}\DIFdelend  cada uno de los indicadores con su fundamento
matemático, la fuente académica que lo sustenta, y un ejemplo práctico calibrado
con datos del área de estudio.


% =============================================================================
\subsection*{Grupo 1: Indicadores Espaciales y de Cobertura}
% =============================================================================

% -----------------------------------------------------------------------
\subsection{Nivel de Cobertura de la Red de Transporte \texorpdfstring{$C$}{C}}
\label{subsec:indicador-cobertura}
% -----------------------------------------------------------------------

\textbf{Fuente principal:} \textcite{fortin2016innovative}

\subsubsection{Definición y Fundamentación}

El nivel de cobertura es el indicador de entrada del modelo: el más independiente
de todos, pues su cálculo se realiza exclusivamente mediante análisis espacial y no
requiere que el grafo de la red esté completamente interconectado. Mide la
proporción del territorio urbano que queda dentro de una distancia caminable de
algún nodo de la red de transporte.

\textcite{fortin2016innovative} explican en su sección de evaluación de accesibilidad que los
indicadores de cobertura más comunes se calculan trazando áreas de influencia
---típicamente un \textit{buffer} de 500 metros o 400 metros como radio de caminata
promedio aceptable--- alrededor de cada parada de la red. A partir de estas áreas,
es posible identificar la proporción de individuos o de empleos que quedan dentro
de cierta distancia de la infraestructura de transporte, lo que permite evaluar el nivel
de accesibilidad en el territorio y señalar con precisión los lugares donde deben
ocurrir las mejoras.

\subsubsection{Fórmula}

\begin{equation}
    C = \frac{\text{Área Cubierta por Buffers}}{\text{Área Total de la Zona de Estudio}}
    \label{eq:cobertura}
\end{equation}

Donde:
\begin{itemize}
    \item $C$ = Nivel de Cobertura de la Red (adimensional, rango $[0, 1]$).
    \item \textbf{Área Cubierta} = Unión de los polígonos de cobertura (isócronas o
    \textit{buffers} circulares) generados alrededor de cada parada de la red,
    medida en km².
    \item \textbf{Área Total} = Superficie del polígono urbano delimitado de la zona
    de estudio, en km².
\end{itemize}

\DIFdelbegin \DIFdel{Insertar figura del indicador C cuando esté disponible
}\DIFdelend %DIF >  [Insertar figura del indicador C cuando esté disponible]
\begin{figure}[H]
    \centering
    \figinclude[width=0.80\textwidth]{Figures/Cap3/CoberturaTransporte.png}
    \caption[Diagrama del indicador de Cobertura $C$]{Esquema del cálculo del
    indicador $C$: polígono de la zona de estudio (Área Total), nodos de transporte
    y círculos de radio 400\,m (Área Cubierta).}
    \label{fig:indicador-C}
    \fuentefigura{Elaboración propia con base en \textcite{fortin2016innovative}.}
\end{figure}


\subsubsection{Consideraciones de Implementación}

Para una representación más precisa que la del \textit{buffer} circular estático, se
implementará el concepto de \termino{isócronas}: áreas de influencia construidas a
partir de la red vial real de OpenStreetMap, que delimitan todos los puntos
alcanzables caminando desde una parada en exactamente 15 minutos a la velocidad
promedio peatonal (1.5 km/h en zonas densas). Las isócronas producen polígonos
irregulares que respetan barreras físicas como vías rápidas, cuerpos de agua o
pendientes pronunciadas, ofreciendo una medida de cobertura sustancialmente más
realista que un círculo uniforme.

El resultado de las isócronas se superpone al censo poblacional del INEGI para
calcular no solo el área cubierta, sino la \termino{población cubierta}: la fracción de
habitantes de la zona metropolitana que reside dentro del área de influencia de
alguna parada de la red.

\subsubsection{Ejemplo Práctico: \cdmx{} y Zona Metropolitana}

\textbf{Datos de entrada:} Coordenadas georreferenciadas de todas las paradas de
la red (Metro, Metrobús, RTP, Trolebús, Cablebús, Mexibús) en formato \geojson{};
polígono del área urbana de la \zmvm{} $\approx 7{,}866$\,km²; capa de manzanas
del INEGI con datos de población del Censo 2020.


\begin{table}[H]
    \centering
    \caption{Estimaciones del indicador $C$ por escenario de red.}
    \label{tab:cobertura-escenarios}
    \begin{tabularx}{\textwidth}{X r r}
        \toprule
        \textbf{Escenario} & \textbf{Área Cubierta} & \textbf{Cobertura $C$} \\
        \midrule
        Red masiva actual (solo Metro + Metrobús + Cablebús)
            & $\approx 450$\,km² & $C \approx 0.06$ \\
        Red completa actual (incluyendo RTP y concesionados)
            & $\approx 2{,}800$\,km² & $C \approx 0.36$ \\
        Red con propuesta del Anillo Periférico (+60 estaciones)
            & $\approx 3{,}150$\,km² & $C \approx 0.40$ \\
        \bottomrule
    \end{tabularx}
\end{table}

\textbf{Implementación en Python (Teórica):}

\begin{lstlisting}[style=estiloPython, caption={Cálculo del indicador $C$ con GeoPandas.}, label={lst:cobertura}]
# GeoPandas + Shapely
gdf_stops = gpd.read_file("paradas_red.geojson")
buffers = gdf_stops.geometry.buffer(800)       # Buffer 800 m
covered_area = buffers.unary_union             # Union de poligonos
city_polygon = gpd.read_file("zmvm.geojson").geometry[0]
C = covered_area.intersection(city_polygon).area / city_polygon.area
\end{lstlisting}

% -----------------------------------------------------------------------
\subsection{Nivel de Alimentación Capilar --- Fuerza Nodal \texorpdfstring{($C_i$)}{(Ci)}}
\label{subsec:alimentacion-capilar}
% -----------------------------------------------------------------------

\textbf{Fuente principal:} \textcite{lee2008}

\subsubsection{Definición y Fundamentación}

Una vez que las estaciones están georreferenciadas y las aristas de la red están
conectadas, el siguiente paso es identificar el rol funcional de cada nodo en la
jerarquía del sistema. \textcite{lee2008}, en su análisis estadístico de la red de
metro de Seúl, establecen que cada estación puede caracterizarse por dos
propiedades complementarias: su \termino{grado} ($k$), definido como el número de
conexiones (aristas) que entran o salen del nodo, y su \termino{fuerza} ($s$),
definida como la suma de los pesos (flujos de pasajeros) de todos los enlaces
conectados a esa estación.

Los autores demuestran que la distribución de la fuerza de los nodos en la red de
metro de Seúl sigue un comportamiento log-normal con un pico característico, lo que
indica que la mayoría de las estaciones atienden un volumen similar de pasajeros,
pero un pequeño número de nodos ---los grandes concentradores o
\textit{hubs}--- acumula flujos desproporcionadamente elevados. Este patrón es
directamente relevante para la red de la \cdmx{}, donde los macro-CETRAMs
(Pantitlán, Indios Verdes, Cuatro Caminos, Taxqueña) operan como \textit{hubs}
nodales que concentran la articulación entre las redes capilares de superficie y las
arterias masivas.

\subsubsection{Fórmula}

\begin{equation}
    C_i = \sum_{j \in S} A_{ij} \cdot \peso{i}{j}
    \label{eq:fuerza-nodal}
\end{equation}

Donde:
\begin{itemize}
    \item $C_i$ = Nivel de Alimentación Capilar (Fuerza Nodal) de la estación
    masiva $i$.
    \item $S$ = Subconjunto de rutas pertenecientes exclusivamente a la red de
    superficie (RTP, microbús, trolebús, concesionados).
    \item $A_{ij}$ = Elemento de la matriz de adyacencia: 1 si la ruta de superficie
    $j$ conecta con la estación masiva $i$, 0 en caso contrario.
    \item $\peso{i}{j}$ = Peso de la conexión: capacidad de pasajeros por hora o flujo
    real medido de la ruta $j$ en la estación $i$.
\end{itemize}

 
\textcite{lee2008} establecen que el peso $\peso{i}{j}$ de un enlace representa el
flujo de pasajeros que transita por él, y que la fuerza $s_i$ de una estación
corresponde al número de pasajeros que llegan y parten de ella. En el modelo
propuesto, $\peso{i}{j}$ se inicializa con la capacidad operativa declarada de cada
ruta y se refinará en fases posteriores con el Coeficiente de Fricción Vial.

% Figura de Fuerza Nodal 
\begin{figure}[H]
    \centering
    \figinclude[width=0.70\textwidth]{Figures/Cap3/NiveldeAlimentacionCapilar.png}
    \caption[Diagrama del indicador de Alimentación Capilar $C_i$]{Esquema de
    fuerza nodal: rutas de superficie $F_1 \ldots F_n$ convergiendo en una
    estación masiva de alta conectividad de entrada.}
    \label{fig:indicador-Ci}
    \fuentefigura{Elaboración propia con base en \textcite{lee2008}.}
 \end{figure}


\subsubsection{Ejemplo Práctico: \cdmx{} y Zona Metropolitana}

\begin{table}[H]
    \centering
    \caption{Estimaciones del indicador $C_i$ en nodos representativos.}
    \label{tab:capilar-nodos}
    \begin{tabularx}{\textwidth}{X r r}
        \toprule
        \textbf{Estación masiva} &
        \textbf{$k_{in}$ (rutas capilares)} &
        \textbf{$C_i$ estimado (pax/h)} \\
        \midrule
        Pantitlán (Hub CETRAM)          & $\approx 120$ rutas & Alto \\
        Indios Verdes (CETRAM)          & $\approx 90$ rutas  & Alto \\
        Nodo Norte propuesto (Periférico) & $\approx 40$ rutas (est.) & Medio-alto \\
        \bottomrule
    \end{tabularx}
\end{table}

\textbf{Implementación en Python (Teórica):}

\begin{lstlisting}[style=estiloPython, caption={Cálculo de grado capilar por nodo con NetworkX.}, label={lst:fuerza-nodal}]
# NetworkX - Grado del nodo filtrado por tipo de arista
feeder_degree = {
    node: sum(1 for _, _, d in G.edges(node, data=True)
              if d.get('tipo') == 'superficie')
    for node in G.nodes()
}
\end{lstlisting}


% =============================================================================
\subsection*{Grupo 2: Indicadores de Eficiencia (Tiempo y Forma)}
% =============================================================================

% -----------------------------------------------------------------------
\subsection{Índice de Ruta Directa --- \texorpdfstring{\textit{Detour Factor} ($DI$)}{Detour Factor (DI)}}
\label{subsec:ruta-directa}
% -----------------------------------------------------------------------

\textbf{Fuente principal:} \textcite{fortin2016innovative}

\subsubsection{Definición y Fundamentación}

El Índice de Ruta Directa cuantifica qué tan eficiente es la topología de la red en
términos geométricos puros: mide cuánto se alarga un viaje al estar obligado a
seguir la infraestructura real de la red en lugar de desplazarse en línea recta entre
origen y destino. \textcite{fortin2016innovative} proponen este indicador como parte de su
método orientado a grafos para caracterizar sistemáticamente redes de transporte a
partir de datos \gtfs{}, donde el análisis de la topología implica calcular los caminos
mínimos entre pares de nodos y derivar indicadores de desempeño.

Es importante precisar una limitación inherente: el $DI$ no incorpora los pesos
ponderados de las aristas (tiempos reales, frecuencias, factores de congestión),
sino únicamente la ubicación geométrica de los nodos. Es, por tanto, un indicador
de la \termino{eficiencia estructural} de la red ---independiente de sus condiciones
operativas--- y resulta especialmente útil para evaluar el trazado de nuevas rutas o
comparar configuraciones alternativas antes de incorporar la dinámica de tráfico.

\subsubsection{Fórmula}

\begin{equation}
    DI = \frac{d_{\text{red}}(i,\, j)}{d_{\text{euclidiana}}(i,\, j)}
    \label{eq:detour-factor}
\end{equation}

Donde:
\begin{itemize}
    \item $DI$ = Índice de Ruta Directa (adimensional, siempre $DI \geq 1.0$).
    \item $d_{\text{red}}(i, j)$ = Distancia del camino mínimo entre los nodos $i$
    y $j$ medida a lo largo de las aristas del grafo, en km.
    \item $d_{\text{euclidiana}}(i, j)$ = Distancia en línea recta (distancia
    \textit{haversine} para coordenadas geográficas) entre los nodos $i$ y $j$,
    en km.
\end{itemize}

Un valor $DI = 1.0$ representa eficiencia geométrica perfecta (el trayecto en red
coincide con la línea recta). Valores $DI \gg 1.0$ indican que la topología fuerza
a los usuarios a rodeos significativos respecto al desplazamiento ideal.

% Figura deligura del Detour Factor (Distancia Euclidiana)
\begin{figure}[H]
    \centering
    \figinclude[width=0.65\textwidth]{Figures/Cap3/DistanciaEuclidiana.png}
    \caption[Diagrama del \textit{Detour Factor} $DI$]{Comparación entre la
    distancia en red (arco curvado) y la distancia euclidiana (línea punteada)
    entre origen O y destino D.}
    \label{fig:indicador-DI}
    \fuentefigura{Elaboración propia con base en \textcite{fortin2016innovative}.}
\end{figure}

\subsubsection{Ejemplo Práctico: \cdmx{} y Zona Metropolitana}

\begin{table}[H]
    \centering
    \caption{Comparativa del indicador $DI$ antes y después del Anillo Periférico.}
    \label{tab:detour-factor}
    \begin{tabularx}{\textwidth}{X r r r r r}
        \toprule
        \textbf{Par O--D} &
        $d_{\text{eucl.}}$ &
        $d_{\text{red}}$ actual &
        $DI$ actual &
        $d_{\text{red}}$ con Anillo &
        $DI$ con Anillo \\
        \midrule
        Xochimilco $\to$ Santa Fe    & 15\,km & 28\,km & 1.87 & 18\,km & 1.20 \\
        Perisur $\to$ Naucalpan      & 18\,km & 42\,km & 2.33 & 26\,km & 1.44 \\
        Iztapalapa $\to$ Tlalnepantla & 22\,km & 55\,km & 2.50 & 30\,km & 1.36 \\
        \bottomrule
    \end{tabularx}
\end{table}

\textbf{Implementación en Python (Teórica):}

\begin{lstlisting}[style=estiloPython, caption={Cálculo del \textit{Detour Factor} $DI$.}, label={lst:detour}]
from geopy.distance import geodesic
import networkx as nx
 
d_euclidiana = geodesic(coords_origen, coords_destino).km
ruta_minima = nx.shortest_path(G, origen, destino, weight='distancia_km')
d_red = sum(G[u][v]['distancia_km'] for u, v in zip(ruta_minima, ruta_minima[1:]))
DI = d_red / d_euclidiana
\end{lstlisting}

% -----------------------------------------------------------------------
\subsection{Tiempo de Viaje Promedio \texorpdfstring{($T$)}{(T)}}
\label{subsec:tiempo-viaje}
% -----------------------------------------------------------------------

\textbf{Fuente principal:} \textcite{aldous2019optimal}

\subsubsection{Definición y Fundamentación}

El Tiempo de Viaje Promedio es la métrica de desempeño global del sistema y el
criterio de optimización central del modelo. \textcite{aldous2019optimal} plantean
que el criterio de optimalidad en el diseño de redes de transporte consiste en
minimizar el tiempo promedio de viaje entre pares de puntos elegidos
independientemente según su distribución en el territorio.

De manera particularmente relevante para la hipótesis central de este trabajo,
\textcite{aldous2019optimal} demuestran analíticamente que, a medida que la
longitud total de la red aumenta, ocurre una transición brusca en un valor crítico
$L_c$ en el que aparece un bucle o anillo en la topología óptima de la red. Esto
implica que, para redes de transporte a escala metropolitana, la incorporación de un
corredor perimetral ---como el Anillo Periférico propuesto--- no solo es viable, sino
que constituye la forma \textit{óptima esperada} de expansión de la red una vez
superado cierto umbral de longitud total. Este resultado teórico provee el sustento
matemático directo para la hipótesis de investigación.

\subsubsection{Fórmula}

\begin{equation}
    T = \frac{1}{N \cdot (N-1)} \sum_{i \neq j} t_{ij}
    \label{eq:tiempo-promedio}
\end{equation}

Donde:
\begin{itemize}
    \item $T$ = Tiempo de Viaje Promedio global de la red, en minutos.
    \item $N$ = Número total de nodos (estaciones y paradas) del grafo.
    \item $t_{ij}$ = Tiempo mínimo de viaje desde el nodo $i$ hasta el nodo $j$,
    calculado mediante el algoritmo de \dijkstra{} sobre el grafo ponderado,
    en minutos. Incluye tiempo en vehículo + tiempo de caminata + tiempo de
    espera en andén + penalización por transbordo.
\end{itemize}

El indicador $T$ se calculará en dos escenarios: la red metropolitana actual y la red
con la incorporación del Anillo Periférico. La diferencia
$\Delta T = T_{\text{actual}} - T_{\text{propuesto}}$ constituye la medida
cuantitativa principal de la viabilidad de la propuesta.

% Figura del indicador T (Tiempo de Viaje Promedio)
\begin{figure}[H]
    \centering
    \figinclude[width=0.72\textwidth]{Figures/Cap3/TiempoViajePromedio.png}
    \caption[Diagrama del indicador de Tiempo de Viaje Promedio $T$]{Red con
    nodos $N$, aristas grises (enlaces conectores) y caminos mínimos
    $t_{ij}$ (caminos más cortos) resaltados.}
    \label{fig:indicador-T}
    \fuentefigura{Elaboración propia con base en \textcite{aldous2019optimal}.}
\end{figure}


\subsubsection{Ejemplo Práctico: \cdmx{} y Zona Metropolitana}

\begin{table}[H]
    \centering
    \caption{Estimaciones del indicador $T$ por escenario.}
    \label{tab:tiempo-promedio}
    \begin{tabularx}{\textwidth}{X r}
        \toprule
        \textbf{Escenario} & \textbf{$T$ estimado} \\
        \midrule
        Red radial actual (sin anillo) & $\approx 85$ minutos \\
        Red con propuesta del Anillo Periférico & $\approx 65$ minutos \\
        Reducción $\Delta T$ & $\approx 20$ min ($-24\%$) \\
        \bottomrule
    \end{tabularx}
\end{table}

\textbf{Implementación en Python:}

\begin{lstlisting}[style=estiloPython, caption={Cálculo del Tiempo de Viaje Promedio $T$ con NetworkX.}, label={lst:tiempo}]
import networkx as nx
# Para redes grandes: usar muestra representativa de pares O-D
T = nx.average_shortest_path_length(G, weight='tiempo_minutos')
\end{lstlisting}


% -----------------------------------------------------------------------
\subsection{Penalización por Transferencia \texorpdfstring{($W_{s,h}$)}{(Ws,h)}}
\label{subsec:penalizacion-transferencia}
% -----------------------------------------------------------------------

\textbf{Fuentes principales:} \textcite{allauca2023}; \textcite{aldous2019optimal}

\subsubsection{Definición y Fundamentación}

La penalización por transferencia es el costo temporal adicional \DIFaddbegin \DIFadd{en el }\DIFaddend que un usuario
incurre al cambiar de una línea o modo de transporte a otro durante su viaje. Su
correcta modelación es fundamental para que el grafo multimodal refleje la
experiencia real del usuario: un algoritmo de \dijkstra{} que no penalice los
transbordos subestimará sistemáticamente los tiempos de viaje en rutas que
requieren múltiples cambios de línea, induciendo al modelo a preferir rutas con
transbordos frecuentes que en la práctica resultan más lentas o incómodas.

\textcite{allauca2023} justifica la incorporación de pesos compuestos en las aristas
del grafo dirigido valorado como mecanismo para modelar correctamente el costo
total de cada ruta. La teoría de grafos provee algoritmos exactos de caminos
mínimos ---específicamente el algoritmo de \dijkstra{} y el algoritmo de
\floyd{}--- para determinar la ruta óptima entre cada par de nodos una vez que los
pesos de las aristas reflejan todos los componentes del costo de viaje.

\subsubsection{Fórmulas}

Se presentan dos variantes complementarias:

\medskip
\noindent\textbf{Variante 1 --- Penalización Estática (simplificada):}

\begin{equation}
    \text{Peso total} = T_{\text{viaje}} + \left(N_{\text{transbordos}}
    \times P_{\text{constante}}\right)
    \label{eq:penalizacion-estatica}
\end{equation}

Donde $P_{\text{constante}}$ es un castigo fijo de tiempo por cada transbordo
(valor de referencia: 12 minutos).

\medskip
\noindent\textbf{Variante 2 --- Penalización Dinámica (por arista de transferencia):}

\begin{equation}
    W_{s,h} = T_{\text{walk}}(s,\, h) + T_{\text{wait}}(s,\, h)
    \label{eq:penalizacion-dinamica}
\end{equation}

Donde:
\begin{itemize}
    \item $W_{s,h}$ = Peso total de la arista artificial de transferencia entre la
    parada $s$ y el andén de la línea $h$, en minutos.
    \item $T_{\text{walk}}(s, h)$ = Tiempo de caminata entre la parada de descenso
    y el andén de abordaje, en minutos.
    \item $T_{\text{wait}}(s, h)$ = Tiempo de espera promedio en el andén de la
    línea $h$, calculado como la mitad del intervalo de paso (\textit{headway})
    promedio de la línea.
\end{itemize}


% Figura de Penalización por Transferencia
\begin{figure}[H]
    \centering
    \figinclude[width=0.68\textwidth]{Figures/Cap3/PenalizacionTransferencia.png}
    \caption[Diagrama del indicador de Penalización por Transferencia $W_{s,h}$]{
    Nodo de transferencia con la descomposición del peso total en sus componentes
    de tiempo de caminata y espera.}
    \label{fig:indicador-W}
    \fuentefigura{Elaboración propia con base en \textcite{allauca2023}.}
\end{figure}

 
\subsubsection{Nota Metodológica sobre Fuentes de Datos}

El uso de horarios \gtfs{} estáticos para estimar $T_{\text{wait}}$ introduce un
sesgo conocido: los datos estandarizados representan únicamente horarios
planificados, no tiempos reales de operación. \textcite{fortin2016innovative} reconocen esta
limitación y señalan que los horarios inscritos pueden ser incorrectos debido a la
congestión o a errores de codificación. Para mitigar este sesgo, la variante
dinámica de la penalización se calibrará mediante peticiones automatizadas a la
\textit{Google Maps Directions API} (modo \texttt{transit}, con variación del
parámetro \texttt{departure\_time}), que utiliza datos de \textit{crowdsourcing}
(GPS de teléfonos móviles) para devolver tiempos de espera y caminata ajustados
a condiciones históricas reales en distintos cortes horarios.

\subsubsection{Ejemplo Práctico: \cdmx{} y Zona Metropolitana}

\textbf{Par de transferencia:} Pantitlán Línea 1 Metro $\to$ Pantitlán Línea 9 Metro.

\begin{table}[H]
    \centering
    \caption{Estimaciones del indicador $W_{s,h}$ en la transferencia Pantitlán L1--L9.}
    \label{tab:penalizacion}
    \begin{tabularx}{\textwidth}{X r r}
        \toprule
        \textbf{Componente} & \textbf{Valor estático} &
        \textbf{Valor dinámico (hora pico 8:00\,h)} \\
        \midrule
        $T_{\text{walk}}$ (caminata entre andenes) & 3 min & 5 min \\
        $T_{\text{wait}}$ (espera en andén L9) & 4.5 min & 7 min \\
        $W_{s,h}$ \textbf{total} & \textbf{7.5 min} & \textbf{12 min} \\
        \bottomrule
    \end{tabularx}
\end{table}

% -----------------------------------------------------------------------
\subsection{Coeficiente de Fricción Vial \texorpdfstring{($CF_h$)}{(CFh)}}
\label{subsec:friccion-vial}
% -----------------------------------------------------------------------

\textbf{Fuente principal:} \textcite{fortin2016innovative}

\subsubsection{Definición y Fundamentación}

El Coeficiente de Fricción Vial es un multiplicador dinámico que corrige la
sobreestimación de eficiencia que ocurre cuando se asigna a las rutas de transporte
de superficie un tiempo de viaje ideal, sin considerar el efecto de la congestión
vehicular. \textcite{fortin2016innovative} reconocen que los datos estandarizados del \gtfs{}
representan únicamente horarios planificados, y que utilizar tiempos ajustados a las
condiciones reales de tráfico devuelve estimaciones sustancialmente más precisas al
ejecutar el algoritmo de \dijkstra{}.

Este indicador distingue dos categorías de modos con comportamientos radicalmente
distintos respecto a la congestión: los \termino{modos confinados} (Metro, Cablebús,
Suburbano, Tren Ligero), que circulan en infraestructura segregada y mantienen un
$CF = 1.0$ constante, y los \termino{modos de tráfico mixto} (RTP, Trolebús,
Metrobús en tramos compartidos, microbuses concesionados), cuyo tiempo de viaje
real se incrementa en función de la hora del día y las condiciones de la vialidad.

\subsubsection{Fórmula}

\begin{equation}
    CF_h = \frac{T_{\text{real},h}}{T_{\text{ideal}}}
    \label{eq:friccion-vial}
\end{equation}

\begin{equation}
    T_{\text{real},h} = T_{\text{ideal}} \times CF_h
    \label{eq:tiempo-real}
\end{equation}

Donde:
\begin{itemize}
    \item $CF_h$ = Coeficiente de Fricción Vial en la franja horaria $h$
    (adimensional, $CF_h \geq 1.0$).
    \item $T_{\text{real},h}$ = Tiempo de viaje observado en el tramo a la hora $h$,
    en minutos.
    \item $T_{\text{ideal}}$ = Tiempo de viaje en condiciones de flujo libre
    (referencia: 05:00\,h, sin congestión), en minutos.
\end{itemize}

\subsubsection{Ejemplo Práctico: \cdmx{} y Zona Metropolitana}

\textbf{Tramo de referencia:} Ruta RTP, Insurgentes Norte (5\,km, tiempo en flujo
libre: 10 minutos).

\begin{table}[H]
    \centering
    \caption{Valores del coeficiente $CF_h$ por franja horaria.}
    \label{tab:friccion-vial}
    \begin{tabularx}{\textwidth}{X r r}
        \toprule
        \textbf{Franja horaria} & \textbf{$T_{\text{real}}$ observado} &
        \textbf{$CF_h$} \\
        \midrule
        05:00\,h (flujo libre)          & 10 min   & 1.00 (referencia) \\
        08:00\,h (pico máximo)          & 18.5 min & \textbf{1.85} \\
        15:00\,h (pico vespertino)      & 16 min   & 1.60 \\
        19:00\,h (pico nocturno)        & 17 min   & 1.70 \\
        Metro (cualquier hora)          & 10 min   & \textbf{1.00} (confinado) \\
        \bottomrule
    \end{tabularx}
\end{table}

% =============================================================================
\subsection*{Grupo 3: Topología y Rigor Matemático}
% =============================================================================

% -----------------------------------------------------------------------
\subsection{Centralidad de Intermediación \texorpdfstring{($B(v)$)}{(B(v))}}
\label{subsec:centralidad-intermediacion}
% -----------------------------------------------------------------------

\textbf{Fuentes principales:} \textcite{lee2008}; \textcite{allauca2023}

\subsubsection{Definición y Fundamentación}

La Centralidad de Intermediación cuantifica la importancia de cada nodo de la red
en función de cuántas rutas óptimas de la ciudad lo utilizan como nodo de paso.
Un nodo con alta centralidad de intermediación es aquel que aparece con mayor
frecuencia en los caminos mínimos entre los distintos pares origen-destino del
sistema: su remoción o saturación produce el mayor impacto sobre los tiempos de
viaje globales.

\textcite{allauca2023} señala que la representación de la red de transporte como un
grafo permite realizar un análisis de centralidad para evaluar la importancia de los
nodos (estaciones), lo que sirve fundamentalmente para identificar los nodos críticos
y las rutas más importantes del sistema. \textcite{lee2008}, en su caso de estudio
del metro de Seúl, aplican este principio empíricamente, demostrando que en un
sistema metropolitano la distribución del grado nodal sigue una ley de potencias,
confirmando la naturaleza de red libre de escala del sistema y su consecuente
vulnerabilidad a ataques dirigidos.

\subsubsection{Fórmula}

\begin{equation}
    B(v) = \sum_{s \neq v \neq t} \frac{\sigma_{st}(v)}{\sigma_{st}}
    \label{eq:betweenness}
\end{equation}

Donde:
\begin{itemize}
    \item $B(v)$ = Centralidad de Intermediación de la estación $v$ (adimensional,
    normalizada en $[0, 1]$).
    \item $\sigma_{st}$ = Número total de caminos mínimos desde la estación de
    origen $s$ hasta el destino $t$.
    \item $\sigma_{st}(v)$ = Número de esos caminos mínimos que pasan
    forzosamente por la estación $v$.
\end{itemize}

La versión ponderada del indicador, $B(v, w)$, utiliza el tiempo de viaje real como
métrica de costo en la búsqueda de caminos mínimos, de modo que los nodos
contabilizados son aquellos que aparecen en las rutas más rápidas, no
necesariamente las más cortas en distancia.


% Figura de Betweenness cuando esté disponible
\begin{figure}[H]
    \centering
    \figinclude[width=0.62\textwidth]{Figures/Cap3/CentralidadItnermediacion.png}
    \caption[Diagrama de Centralidad de Intermediación $B(v)$]{Nodo de
    transferencia con alta centralidad: concentra el mayor número de caminos
    mínimos del sistema.}
    \label{fig:indicador-Bv}
    \fuentefigura{Elaboración propia con base en \textcite{lee2008} y
    \textcite{allauca2023}.}
\end{figure}


\subsubsection{Ejemplo Práctico: \cdmx{} y Zona Metropolitana}

\begin{table}[H]
    \centering
    \caption{Variación del indicador $B(v)$ antes y después del Anillo Periférico.}
    \label{tab:betweenness}
    \begin{tabularx}{\textwidth}{X r r r}
        \toprule
        \textbf{Nodo} & \textbf{$B(v)$ Red Actual} &
        \textbf{$B(v)$ con Anillo} & \textbf{Variación} \\
        \midrule
        Pantitlán               & $\approx 0.40$ & $\approx 0.22$ & $-45\%$ \\
        Pino Suárez             & $\approx 0.38$ & $\approx 0.20$ & $-47\%$ \\
        Tacubaya                & $\approx 0.35$ & $\approx 0.19$ & $-46\%$ \\
        Estaciones Anillo (prom.) & ---          & $\approx 0.18$ & $+\infty$ (nuevos) \\
        \bottomrule
    \end{tabularx}
\end{table}

La caída en la centralidad de intermediación de los nodos del centro de la ciudad,
combinada con el surgimiento de nuevos nodos de alta centralidad en el Anillo
Periférico, constituye la demostración algorítmica de que la propuesta descongestiona
el núcleo urbano y redistribuye los flujos hacia la periferia.

\textbf{Implementación en Python:}

\begin{lstlisting}[style=estiloPython, caption={Cálculo de Centralidad de Intermediación con NetworkX.}, label={lst:betweenness}]
import networkx as nx
centrality = nx.betweenness_centrality(G, weight='tiempo_minutos', normalized=True)
# Ordenar por centralidad descendente
top_nodes = sorted(centrality.items(), key=lambda x: x[1], reverse=True)[:10]
\end{lstlisting}


% -----------------------------------------------------------------------
\subsection{Robustez Geométrica --- Vulnerabilidad \texorpdfstring{($\Delta E$)}{(Delta E)}}
\label{subsec:robustez-geometrica}
% -----------------------------------------------------------------------

\textbf{Fuente principal:} \textcite{loterovélez2014}

\subsubsection{Definición y Fundamentación}

La Robustez Geométrica es el indicador de mayor dependencia jerárquica del
modelo: para calcularlo, se requiere conocer previamente qué nodo destruir ---dado
por el Indicador de Centralidad de Intermediación
(\autoref{subsec:centralidad-intermediacion})--- y medir el impacto de esa destrucción
sobre el desempeño global ---dado por el Indicador de Tiempo de Viaje Promedio
(\autoref{subsec:tiempo-viaje})\DIFaddbegin \DIFadd{---}\DIFaddend .

\textcite{loterovélez2014} definen la vulnerabilidad operativamente como la
disminución de la eficiencia de la red después de un ataque o falla. La metodología
estándar para medir la robustez consiste en cuantificar el impacto de la remoción de
un componente de la red, ya sea de forma aleatoria ---simulando una falla fortuita
como el cierre de una estación por inundación o mantenimiento--- o de forma dirigida
hacia los elementos más importantes, simulando un ataque a un nodo crítico. Los
autores señalan que las redes con distribución libre de escala son robustas ante
fallas aleatorias, pero altamente vulnerables ante ataques dirigidos a sus nodos más
conectados, resultado que es la base teórica del experimento de fase final del
presente trabajo.

\subsubsection{Fórmula}

\begin{equation}
    \Delta E = \frac{E_0 - E_f}{E_0}
    \label{eq:robustez}
\end{equation}

Donde:
\begin{itemize}
    \item $\Delta E$ = Caída relativa de eficiencia de la red tras la remoción de un
    nodo o conjunto de nodos (adimensional, rango $[0, 1]$).
    \item $E_0$ = Eficiencia global de la red en el estado completo, expresada como
    el inverso del Tiempo de Viaje Promedio:
    $E_0 = 1 / T_{\text{completa}}$.
    \item $E_f$ = Eficiencia global de la red en el estado fallido (tras la remoción
    del nodo o nodos críticos):
    $E_f = 1 / T_{\text{fallida}}$.
\end{itemize}

Un valor $\Delta E \approx 0$ indica que la red es robusta: la pérdida del nodo no
afecta significativamente el desempeño global (existen rutas alternativas eficientes).
Un valor $\Delta E \to 1$ indica que la red es frágil: la pérdida de ese nodo produce
una degradación severa del servicio.


% figura de Robustez Geométrica
\begin{figure}[H]
    \centering
    \figinclude[width=0.80\textwidth]{Figures/Cap3/RobustezGeometrica.png}
    \caption[Diagrama del indicador de Robustez Geométrica $\Delta E$]{Estado
    completo con ruta directa ($E_0$) versus estado con fallo y rutas
    alternativas de desvío ($E_f$).}
    \label{fig:indicador-deltaE}
    \fuentefigura{Elaboración propia con base en \textcite{loterovélez2014}.}
\end{figure}


\subsubsection{Ejemplo Práctico: \cdmx{} y Zona Metropolitana}

\textbf{Experimento:} Remoción del nodo ``Pantitlán'' (nodo de mayor $B(v)$ en
la red actual).

\begin{table}[H]
    \centering
    \caption{Impacto del indicador $\Delta E$ ante la remoción de Pantitlán.}
    \label{tab:robustez}
    \begin{tabularx}{\textwidth}{X r r r}
        \toprule
        \textbf{Condición} & \textbf{$T$ promedio} &
        \textbf{$E = 1/T$} & \textbf{$\Delta E$} \\
        \midrule
        Red actual completa ($E_0$)          & 85 min & 0.01176 & --- \\
        Red actual sin Pantitlán ($E_f$)     & 133 min & 0.00752 & $\mathbf{36\%}$ \\
        Red con Anillo completo ($E_0'$)     & 65 min & 0.01538 & --- \\
        Red con Anillo sin Pantitlán ($E_f'$) & 78 min & 0.01282 & $\mathbf{17\%}$ \\
        \bottomrule
    \end{tabularx}
\end{table}

La diferencia entre $\Delta E = 36\%$ (red radial sin anillo) y $\Delta E = 17\%$
(red con anillo periférico) demuestra algorítmicamente que la propuesta reduce la
vulnerabilidad del sistema ante la pérdida del nodo más crítico en casi un 53\%,
dotando a la red metropolitana de redundancia estructural.

\textbf{Implementación en Python (Teórica):}

\begin{lstlisting}[style=estiloPython, caption={Cálculo de Robustez Geométrica $\Delta E$ con NetworkX.}, label={lst:robustez}]
import networkx as nx, copy
 
def eficiencia_global(G):
    T = nx.average_shortest_path_length(G, weight='tiempo_minutos')
    return 1 / T
 
E0 = eficiencia_global(G)
 
# Remover nodo critico
G_fallido = copy.deepcopy(G)
G_fallido.remove_node('Pantitlan')
Ef = eficiencia_global(G_fallido)
 
delta_E = (E0 - Ef) / E0
print(f"Caida de eficiencia: {delta_E:.2%}")
\end{lstlisting}

 
% =============================================================================
\subsection{Resumen de Indicadores y Fases de Construcción}
\label{subsec:resumen-indicadores}
% =============================================================================

La Tabla~\ref{tab:resumen-indicadores} sintetiza los ocho indicadores, su símbolo,
la fase metodológica a la que pertenecen y sus dependencias jerárquicas.

\begin{table}[H]
    \centering
    \caption{Resumen de los ocho indicadores de análisis y sus fases de
             construcción.}
    \label{tab:resumen-indicadores}
    \begin{tabularx}{\textwidth}{c l c c X}
        \toprule
        \textbf{\#} & \textbf{Indicador} & \textbf{Símbolo} &
        \textbf{Fase} & \textbf{Dependencias} \\
        \midrule
        1 & Nivel de Cobertura de la Red     & $C$          & 1 & Ninguna (espacial puro) \\
        2 & Nivel de Alimentación Capilar    & $C_i$        & 1 & Indicador 1 \\
        3 & Índice de Ruta Directa           & $DI$         & 1 & Indicadores 1--2 \\
        4 & Tiempo de Viaje Promedio         & $T$          & 3 & Indicadores 1--2 + Fases 1 y 2 completas \\
        5 & Penalización por Transferencia   & $W_{s,h}$    & 2 & Indicadores 1--3 \\
        6 & Coeficiente de Fricción Vial     & $CF_h$       & 2 & Indicadores 1--3 \\
        7 & Centralidad de Intermediación    & $B(v)$       & 3 & Indicadores 1--6 \\
        8 & Robustez Geométrica              & $\Delta E$   & 4 & Todos los anteriores \\
        \bottomrule
    \end{tabularx}
\end{table}

%% ── 4-Capas-Codigo/4-1-Apimetro.tex
\subsection*{4 1 Apimetro}
% !TeX root = ../tesis.tex
%%%%%%%%%%%%%%%%%%%%%%%%%%%%%%%%%%%%%%%%%%%%%%%%%%%%%%%%%%%%%%%%%%%
% SECCIÓN: CAPAS DE CÓDIGO
%%%%%%%%%%%%%%%%%%%%%%%%%%%%%%%%%%%%%%%%%%%%%%%%%%%%%%%%%%%%%%%%%%%
\section*{Capas de Código}
\label{ch:Capas-codigo}
% =============================================================================
%  SECCIÓN 4.1 — Filosofía de Arquitectura
% =============================================================================

\section{Filosofía de la Arquitectura: Hacia la democratización de la Tecnología de planeación}
\label{sec:filosofia-arquitectura}
La arquitectura computacional de esta investigación se fundamenta en la transición hacia herramientas de análisis de acceso abierto, 
rompiendo con la dependencia histórica de ecosistemas de software cerrados y onerosos para la administración pública. 
Históricamente, la planeación de redes de transporte ha estado supeditada al acceso a bases de datos comerciales, como NAVTEQ Map Data 
\parencite{navteq2012}, o al uso de plataformas SIG especializadas como TransCAD \parencite{caliper2010}, cuyos costos de licenciamiento y opacidad 
algorítmica imponen barreras de entrada significativas para los organismos locales.
Frente a este modelo, la democratización tecnológica se apoya en proyectos como OpenStreetMap (OSM) \parencite{osm2012}, que proporcionan capas 
de infraestructura vial libre, permitiendo una construcción transparente de grafos ruteables. Este enfoque se alinea con los preceptos de la \textbf{Free Software Foundation (FSF)}, 
la cual establece que la libertad del software radica en la capacidad de los usuarios para usar, estudiar, modificar y redistribuir las herramientas según las necesidades de la comunidad.
En términos de planeación urbana, la adopción de esta filosofía permite alcanzar una soberanía tecnológica, facilitando que las instituciones públicas acumulen capacidades analíticas autónomas 
y reduzcan la brecha de información en las periferias urbanas. Consecuentemente, el desarrollo de motores propios como \textbf{Apimetro} y \textbf{VFTModel}, basados en lenguajes de programación modernos como Go y Python, 
garantiza que el modelado y la optimización de la red multimodal de la ZMVM sea un proceso reproducible, transparente y accesible para el interés público.
\DIFaddbegin 

\DIFadd{La arquitectura de la investigación se articula en tres capas tecnológicas
interdependientes, cuya interacción se resume en la
Figura~\ref{fig:macro-arquitectura-stack}: }\textbf{\DIFadd{Apimetro}} \DIFadd{como capa de
ingesta y datos; }\textbf{\DIFadd{VFTModel}} \DIFadd{como motor analítico; y
}\textbf{\DIFadd{Transport-gis-zmvm-mjg}} \DIFadd{como interfaz cartográfica de presentación.
Los capítulos siguientes desarrollan cada capa en detalle.
}

\begin{figure}[H]
    \centering
    \figplaceholder{Figures/TiKz_Libraries/Diagramas/macro-arquitectura-stack}
    \caption[Macro-arquitectura del \textit{stack} tecnológico de la investigación]{
    \DIFaddFL{Las tres capas del }\textit{\DIFaddFL{stack}} \DIFaddFL{tecnológico de la investigación. La Capa~1
    (Apimetro) provee los datos geoespaciales vía }\api{} \DIFaddFL{REST; la Capa~2
    (VFTModel) construye el grafo topológico y calcula los indicadores; la
    Capa~3 (Transport-gis) produce la cartografía de presentación mediante QGIS.
    Las flechas indican el flujo de datos entre capas.}}
    \label{fig:macro-arquitectura-stack}
    \fuentefigura{Elaboración propia.}
\end{figure}

\DIFaddend \section{Apimetro: El motor de Datos metropolitano (Ingesta y ETL)}
\label{sec:apimetro-intro}

El proyecto \texttt{Apimetro} tiene su génesis entre los años 2023 y 2024, derivado de una necesidad cartográfica y de modelación espacial para el universo narrativo de \textbf{"La Saga de la Catacumba"}, obra del autor de la presente investigación. 
El primer volumen de esta serie, titulado \textbf{La historia de Quince}, explora una distopía urbana denominada \DIFdelbegin \DIFdel{"Ciudad Capital"}\DIFdelend \DIFaddbegin \DIFadd{«Ciudad Capital»}\DIFaddend , la cual funciona como un homónimo especulativo de la Ciudad de México donde la habitabilidad se rige por la interacción crítica 
de tres variables: transporte, educación y tiempo. En dicha narrativa, el \DIFdelbegin \DIFdel{"tiempo" }\DIFdelend \DIFaddbegin \DIFadd{«tiempo» }\DIFaddend no es solo una dimensión física, sino un recurso de equidad social condicionado por la eficiencia de la infraestructura de movilidad. Para la construcción de este entorno literario, 
se requirió el desarrollo de un plan urbanístico que resolviera la fricción espacial de los protagonistas, quienes residían en polos opuestos de la metrópoli: el oriente y el sur profundo. En la realidad geográfica de la Ciudad de México, estas zonas corresponden a 
demarcaciones con un rezago histórico en transporte masivo, tales como Milpa Alta, Xochimilco, Tláhuac y Tlalpan. En estos sectores, la oferta de movilidad recae primordialmente en la \textbf{Red de Transporte de Pasajeros (RTP)}, 
cuya baja frecuencia y velocidades comerciales limitadas generaban, a nivel narrativo, \DIFdelbegin \DIFdel{"tiempos muertos" }\DIFdelend \DIFaddbegin \DIFadd{«tiempos muertos» }\DIFaddend que dificultaban la fluidez de la trama. Esta problemática orilló a una convergencia multidisciplinaria que dio origen a cuatro proyectos paralelos: 
\textit{Apimetro}, \textit{La Catacumba}, \textit{VFTModel} y \textit{Diarios}.

La primera aproximación técnica para sistematizar la geografía metropolitana fue el desarrollo de \textit{\DIFdelbegin \DIFdel{"}\DIFdelend \DIFaddbegin \DIFadd{«}\DIFaddend TrainMap\DIFdelbegin \DIFdel{"}\DIFdelend \DIFaddbegin \DIFadd{»}\DIFaddend }, una herramienta de visualización que, aunque funcional 
para propósitos creativos, carecía del rigor matemático necesario para una auditoría urbana formal. A partir del año 2025, con el ingreso a la \textbf{Maestría en Urbanismo de la UNAM}, la investigación dio un giro hacia la formalización científica, 
integrando estándares como el \textit{General Transit Feed Specification} (GTFS) y sistemas de información geográfica (SIG).

\texttt{Apimetro} recibió una actualización estructural significativa a mediados de marzo de 2026, orientándose hacia un modelo de 
\textbf{Software as a Service (SaaS)} especializado en el manejo de tipos de datos \texttt{GeoJSON} para la representación cartográfica y el análisis ruteable. El núcleo operativo del sistema se define mediante un proceso \textbf{ETL}\footnote{En esta investigación se adopta la 
denominación \textit{Extraction, Treatment and Liberation} para enfatizar el carácter transformativo y de acceso abierto del proceso; el acrónimo estándar de la industria es \DIFaddbegin \DIFadd{«}\DIFaddend \textit{Extract, Transform, Load}\DIFaddbegin \DIFadd{»}\DIFaddend .},
diseñado para transformar la rigidez de los archivos \gtfs{} en geometrías flexibles. Si bien el estándar \gtfs{} posee una estructura relacional robusta, su formato original dificulta el análisis 
directo en hojas de cálculo o capas \texttt{SHP} convencionales, por lo que la conversión a \texttt{GeoJSON} se determinó como la solución \DIFdelbegin \DIFdel{más }\DIFdelend óptima para la visualización y el cálculo de indicadores topológicos. La persistencia y estructuración de la información 
se ejecutan dentro de una base de datos relacional \textbf{PostgreSQL 15}, potenciada por la extensión espacial \textbf{PostGIS 3.4}. Esta elección tecnológica permite almacenar datos geográficos de manera eficiente y proporciona operaciones de consulta avanzadas, 
como \texttt{ST\_AsGeoJSON} y \texttt{ST\_DWithin}, fundamentales para la ingesta y el ruteo asíncrono. Para la implementación de esta fase, se utilizó el lenguaje de programación \textbf{Python}, aprovechando la robustez de librerías especializadas en el manejo de 
grandes volúmenes de datos como \texttt{Pandas} y \texttt{PySpark}. Gracias a esta arquitectura, la ingesta de datos se consume de manera autosuficiente y óptima, resultando en una base de datos de apenas 38 MB en estado de desarrollo que es fácilmente replicable en 
diversos entornos operativos, incluyendo Linux, Windows, MacOS y plataformas web. Esta portabilidad facilita la auditoría de la red mediante gestores de consultas como \texttt{Postman} o su despliegue local a través de tecnologías de contenedores como \textbf{Docker} y 
\textbf{Docker Compose}. En última instancia, \texttt{Apimetro} garantiza la soberanía tecnológica de la investigación, permitiendo que la replicación del esquema de datos se realice de forma transparente mediante el archivo \texttt{seed.sql} disponible en su repositorio oficial.

\begin{figure}[htbp]
  \centering
  \figinclude[width=0.45\textwidth]{Figures/Cap3/APImetroIcon.png}
  \caption[Ícono Personal de Apimetro]{Ícono Creado para el proyecto Apimetro}
  \label{fig:apimetro-icon}
  \fuentefigura{Elaboración propia.}
\end{figure}

\subsection{Modelo de Datos y Esquema Relacional}
\label{subsec:modelo-datos}

\subsubsection{Esquema Entidad-Relación E-R}
\label{subsubsec:EsquemaER}

El diseño del esquema relacional busca simplificar el estándar complejo del \gtfs{} (que incluye tablas como \textit{trips}, \textit{stop\_times} y \textit{shapes}) en tres entidades primordiales optimizadas para el consumo de la \api{} y el cálculo de indicadores topológicos:

\begin{figure}[H]
    \centering
    \figinclude[width=0.85\textwidth]{Figures/Cap4/DiagramaER.png}
    \caption[Esquema Relacional Simplificado de Apimetro]{Representación de las entidades principales (Estaciones, Líneas y Polígonos) y sus relaciones para el cálculo de indicadores topológicos.}
    \label{fig:diagrama-er}
    \fuentefigura{Elaboración propia.}
\end{figure}


Como se observa en la Figura \ref{fig:diagrama-er}, el modelo se articula bajo los siguientes componentes:

\begin{itemize} 
    \item \textbf{Estaciones:} Nodo puntual que contiene atributos de conectividad y ubicación. 
    \item \textbf{Líneas:} Aristas de tipo \texttt{MultiLineString} que almacenan métricas operativas (velocidad, capacidad, frecuencia). 
    \item \textbf{Polígonos:} Áreas administrativas para filtrado por alcaldía o municipio. 
    \item \textbf{Relaciones:} Las líneas se vinculan con las estaciones a través de un identificador de sistema y ramal, permitiendo la reconstrucción del grafo metropolitano.
\end{itemize}


\subsubsection{Definición de Tablas Principales}
\label{subsubsec:DefinicionSQL}

El esquema relacional se implementa en tres tablas: \texttt{estaciones} (nodos puntuales con atributos de conectividad), \texttt{lineas} (aristas \texttt{MultiLineString} con métricas operativas de velocidad, frecuencia y capacidad) y \texttt{poligonos} (áreas administrativas para filtrado por alcaldía o municipio). Todas las columnas de geometría utilizan el tipo \texttt{GEOMETRY} con proyección WGS~84 (SRID~4326) e índices espaciales \texttt{GIST} para optimizar las consultas de ruteo. La definición completa del esquema \textsc{sql} se presenta en el Anexo~\ref{anx:apimetro-sql}.

La arquitectura de estas tablas responde a dos necesidades fundamentales identificadas en la investigación:
\begin{itemize}
    \item \textbf{Optimización de Consultas Geográficas:} El uso de columnas de tipo \texttt{GEOMETRY} y la creación de índices
    \texttt{GIST} permiten realizar operaciones espaciales masivas en milisegundos, tales como \texttt{ST\_AsGeoJSON} para la
    exportación de geometrías, \texttt{ST\_DWithin} para el cálculo de áreas de influencia o cobertura, y \texttt{ST\_Length} para
    determinar las distancias reales sobre la red vial.
    \item \textbf{Cálculo de Indicadores Operativos:} A diferencia de un \texttt{SIG}
    convencional, la tabla \texttt{lineas} integra atributos operativos críticos como la \textit{capacidad} del vehículo y la
    \textit{frecuencia} de paso. Esto permite que el motor \texttt{VFTModel} calcule automáticamente indicadores de desempeño como la
    \textbf{Fuerza Nodal ($C_i$)} y el \textbf{Tiempo de Viaje ($T$)}, ponderando el grafo metropolitano no solo por su fricción espacial,
    sino por la oferta real de transporte y las penalizaciones por transferencia.
\end{itemize}

Finalmente, es importante señalar que esta infraestructura está alineada con los preceptos de soberanía tecnológica y reproducibilidad
científica, siendo totalmente instalable en cualquier entorno (Docker, Linux, Windows o MacOS) mediante el archivo de inicialización
\texttt{init.sql} (o \texttt{seed.sql}) disponible en el repositorio oficial de \texttt{Apimetro} en GitHub.


\subsection{Tecnologías Backend}
\label{subsec:backend}
La arquitectura del motor de datos de la \zmvm se ha implementado utilizando el lenguaje de programación \textbf{Go 1.25}. 
La elección de esta tecnología no es fortuita; responde a la necesidad de procesar un grafo metropolitano con miles de nodos y aristas bajo condiciones 
de alta concurrencia y baja latencia. Go, al ser un lenguaje compilado, ofrece un rendimiento cercano a C++, pero con una gestión de memoria (\textit{Garbage Collection}) 
eficiente, lo cual es ideal para el ruteo asíncrono y el despacho de geometrías masivas en formato \texttt{GeoJSON}.
Para la capa de servicios web, se seleccionó el framework \textbf{Gin v1.9}. Gin se distingue por su enfoque minimalista y su velocidad en el ruteo de peticiones HTTP mediante 
un árbol de prefijos, lo que permite que los endpoints de \texttt{Apimetro} respondan en milisegundos incluso bajo cargas pesadas. El sistema se organiza bajo el patrón de diseño 
\textbf{MVC (Modelo-Vista-Controlador)}, permitiendo una separación clara entre la lógica de negocio, el dominio matemático y la persistencia de datos.

\subsubsection{Configuración y Entorno de Ejecución}
\label{subsubsec:configuracionejecucion}

Para garantizar la portabilidad y soberanía tecnológica, el backend se configura mediante variables de entorno que permiten su despliegue en contenedores \textbf{Docker}. Las variables críticas cubren la conexión al motor espacial (\texttt{DB\_HOST}, \texttt{DATABASE\_URL}) y el modo de ejecución del framework (\texttt{GIN\_MODE}); su especificación completa se presenta en el Cuadro~\ref{tab:env-vars} del Anexo~\ref{anx:apimetro-backend}.

\subsubsection{Estructura de Endpoints y Lógica de Respuesta}
\label{subsubsec:estructuraEndpoints}
La capa de rutas de \texttt{Apimetro} expone la interfaz pública del sistema, actuando como el punto de entrada para el consumo de datos georreferenciados. 
Todos los endpoints de mapas han sido diseñados para devolver objetos de tipo \texttt{FeatureCollection} bajo el estándar \textbf{GeoJSON RFC 7946} \parencite{cabrera2025apimetro}. 
Esta arquitectura permite que el sistema opere bajo un modelo de \textbf{SaaS (Software as a Service)}, donde clientes analíticos como \textit{VFTModel} o visores 
cartográficos consumen datos ya normalizados y ruteables. Esto elimina la necesidad de procesar los archivos de texto plano de los feeds \gtfs{} originales en cada consulta, 
optimizando significativamente el rendimiento computacional al calcular indicadores como el \textbf{Tiempo de Viaje Promedio (T)} y el \textbf{Índice de Ruta Directa (DI)} \parencite{cabrera2026VFTModel}.

Para ejemplificar la versatilidad del backend desarrollado en \textbf{Go}, se analizan las dos estructuras principales del sistema: la gestión de nodos puntuales (estaciones) y la recuperación de geometrías 
complejas (líneas y ramales).
\paragraph{Control de Estaciones (Nodos):} El controlador \texttt{estaciones.go} gestiona la lógica de filtrado de las paradas del sistema metropolitano capturando parámetros dinámicos de la \textsc{url} (sistema, alcaldía, jerarquía de transporte) para segmentar la red y permitir auditorías locales precisas. El código fuente completo se presenta en el Anexo~\ref{anx:apimetro-backend}.

\paragraph{Optimización de Indicadores:}
La separación de estas lógicas en controladores especializados permite que una petición de ruteo para el \textbf{Tercer Contorno} sea resuelta mediante índices espaciales \texttt{GIST} en milisegundos \parencite{cabrera2025apimetro}.

Esta estructura garantiza que la planeación urbana no dependa de la manipulación manual de datos pesados, permitiendo que la simulación de escenarios de transporte sea un proceso dinámico, 
reproducible y basado en la soberanía tecnológica de herramientas de código abierto \parencite{cabrera2025apimetro}.


\subsection{Salidas Estructurales (GeoJSON)}
\label{subsec:estructura-geojson}

La elección del formato de intercambio de datos geoespaciales constituye una decisión crítica para la operatividad de los sistemas de transporte modernos. En el desarrollo del motor \texttt{Apimetro}, se ha priorizado
 el uso del estándar \textbf{GeoJSON (RFC 7946)} sobre el formato tradicional \textit{ESRI Shapefile} (SHP). Mientras que el formato SHP representa un estándar de legado que requiere una colección de múltiples archivos 
independientes (\textit{.shp, .dbf, .shx}) para su funcionamiento y presenta limitaciones de codificación en entornos web, el GeoJSON es un formato basado íntegramente en texto plano (\textit{JSON}) diseñado 
específicamente para la transferencia ágil de datos geográficos en aplicaciones móviles y sistemas distribuidos \parencite{cabrera2025apimetro}.

Una de las ventajas determinantes de este formato es su \textbf{compatibilidad nativa} con las librerías de mapas de alto rendimiento más utilizadas en la industria tecnológica, tales como \textit{Leaflet}, 
\textit{Mapbox} y \textit{MapLibre} \parencite{cabrera2025apimetro}. Esta característica permite que los datos de movilidad de la \zmvm se integren directamente en visores interactivos sin necesidad de 
realizar transformaciones pesadas en el lado del cliente o recurrir a servidores de mapas intermediarios. Al ser un estándar basado en texto, facilita la transparencia en la auditoría de los datos y 
asegura la soberanía tecnológica de la investigación al utilizar insumos oficiales de acceso abierto \parencite{adip_gtfs_2024}.

Asimismo, la integración de GeoJSON con bases de datos espaciales como \textbf{PostGIS} permite que la generación de mapas sea un proceso dinámico y asíncrono. Mediante funciones nativas de conversión, el sistema 
transforma la topología de la red almacenada en tablas de estaciones y líneas en geometrías ruteables de forma inmediata \parencite{fortin2016innovative}. 
Este enfoque supera la rigidez de los archivos vectoriales estáticos, permitiendo que cada objeto geográfico transporte, dentro de su estructura de \textit{properties}, no solo su ubicación espacial\DIFaddbegin \DIFadd{, }\DIFaddend sino
también sus métricas operativas críticas —como velocidad promedio, frecuencia y capacidad del vehículo—. Esta consolidación de información en un solo objeto estructurado es fundamental para el procesamiento 
analítico del motor \textit{VFTModel} al evaluar la eficiencia de la red \parencite{cabrera2026VFTModel}.

\subsection{Estructura de JSON Descriptivo (No Geográfico)}
\label{subsec:estructura-geojson-descriptivo}

Complementariamente a las salidas espaciales, el motor \texttt{Apimetro} provee una capa de \textbf{respuestas descriptivas} en formato JSON estándar. Esta estructura está diseñada específicamente para gestionar atributos 
alfanuméricos, metadatos históricos y parámetros de estilo que no requieren de una representación geométrica inmediata para su procesamiento \parencite{cabrera2025apimetro}. 
La utilidad de esta salida radica en su eficiencia para alimentar componentes de interfaz de usuario, tales como fichas técnicas de estaciones, tablas de atributos operativos o la configuración de la simbología visual 
en visores cartográficos sin incurrir en el costo computacional de transferir coordenadas geográficas masivas.

A diferencia del estándar GeoJSON, orientado al renderizado de mapas, el JSON descriptivo permite realizar consultas rápidas sobre la jerarquía administrativa y la cronología del sistema de transporte~\parencite{cabrera2025apimetro}. Estos registros vinculan la información operativa almacenada en la tabla \texttt{historico\_operacion} con las identidades comerciales de cada línea. Un ejemplo representativo de esta salida se presenta en el Anexo~\ref{anx:apimetro-salidas}.

Entre los campos de mayor relevancia para la planeación se encuentran \texttt{color\_esp} —que permite la estandarización visual de las rutas bajo los lineamientos oficiales de identidad institucional de la \zmvm{}— y \texttt{clasificacion} —que distingue entre infraestructura existente, futura o eliminada—, ambos fundamentales para que el motor \textit{VFTModel} realice comparaciones estadísticas entre la red actual y los escenarios de expansión propuestos~\parencite{cabrera2026VFTModel, adip_gtfs_2024}. Esta separación de intereses entre los datos descriptivos y los geográficos asegura que el flujo de información sea escalable y se adapte a las necesidades específicas de la planeación urbana estratégica y la auditoría de redes complejas~\parencite{fortin2016innovative}.

\subsection{Estructura de GeoJSON (Geográfico)}
\label{subsec:estructura-geojson-geografico}

La salida geográfica de \texttt{Apimetro} constituye la columna vertebral del sistema, transformando la persistencia relacional de la base de datos en objetos espaciales listos para el análisis y la visualización \parencite{cabrera2025apimetro}. 
A diferencia de los formatos tabulares, la estructura GeoJSON permite una jerarquía anidada que encapsula la topología de la red metropolitana en un solo archivo de texto plano. A continuación, se describen los cuatro componentes fundamentales 
que integran cada respuesta del motor:

\paragraph{1. FeatureCollection:} Es el contenedor de nivel superior que agrupa un conjunto de entidades geográficas. En el contexto de esta investigación, un \texttt{FeatureCollection} puede representar la totalidad de las estaciones de un sistema 
(como el Metrobús) o el trazo completo de un ramal operativo.

\paragraph{2. Feature:} Representa una unidad individual dentro de la red, ya sea un nodo puntual (estación) o una arista (segmento de línea). Cada \texttt{Feature} actúa como el vínculo entre la geometría física y sus atributos descriptivos.

\paragraph{3. Geometry:} Define el componente espacial de la entidad utilizando coordenadas geográficas bajo el elipsoide \textbf{WGS 84 (SRID 4326)}. \texttt{Apimetro} emplea primordialmente dos tipos de geometrías: 
\texttt{Point} para la geolocalización exacta de estaciones y paradas, y \texttt{MultiLineString} para representar los trazos de las rutas, permitiendo capturar la continuidad de la infraestructura vial incluso en segmentos fragmentados.

\paragraph{4. Properties:} Es el objeto donde se almacena la metadata enriquecida y las métricas operativas. En esta sección, el motor integra datos como la \texttt{alcaldia\_municipio}, el \texttt{sistema}, y variables críticas para el motor \textit{VFTModel} 
como la \texttt{velocidad\_promedio\_kmh} y la \texttt{frecuencia\_minutos}.

Un ejemplo representativo de la estructura de una estación de transporte masivo generada por la \api{} se presenta en el Anexo~\ref{anx:apimetro-salidas}.

La propiedad \texttt{coordinates} proporciona la base para el cálculo de la \textbf{Distancia Euclidiana} necesaria para obtener el \textbf{Índice de Ruta Directa ($DI$)}. Al mismo tiempo, las \texttt{properties} permiten realizar filtrados dinámicos en el lado del cliente, facilitando que aplicaciones móviles o visores web presenten únicamente la información relevante para la planeación del corredor anillar propuesto \parencite{cabrera2025apimetro, adip_gtfs_2024}. Esta arquitectura asegura que la información geográfica no sea solo una representación visual, sino un insumo matemático auditable y reproducible \parencite{fortin2016innovative}.
\subsection{El Futuro de Apimetro}
\label{subsec:futuro-apimetro}

\texttt{Apimetro} se define no solo como una herramienta técnica, sino como un proyecto de \textbf{software libre} comprometido con la democratización de los datos de movilidad en la \zmvm. Bajo la premisa de soberanía tecnológica, 
el proyecto se encuentra en una fase de expansión activa, proyectando un crecimiento significativo en su cobertura y capacidades analíticas en los próximos meses \parencite{cabrera2025apimetro}. 
Para facilitar la transparencia, el estudio y la redistribución del conocimiento generado, el ecosistema de \texttt{Apimetro} se articula a través de cuatro plataformas principales:

\paragraph{1. Página Oficial (apimetro.dev):} Es el punto de entrada principal para usuarios y desarrolladores. En este portal se aloja la documentación interactiva (Swagger UI), el estado de salud del servidor de producción y los 
accesos rápidos a los recursos de la API masiva.

\begin{figure}[htbp]
  \centering
  \figinclude[width=0.85\textwidth]{Figures/Cap3/PantallaInicialApimetro.png}
  \caption{Interfaz principal del portal \url{apimetro.dev}.}
  \label{fig:screenshot-web}
\end{figure}

\paragraph{2. Repositorio de Código en GitHub:} Constituye el núcleo técnico del proyecto. Aquí se encuentra el código fuente en \textbf{Go}, los motores de ingesta \textbf{ETL} y las configuraciones de despliegue mediante \textbf{Docker}. 
El repositorio está bajo una licencia \textit{Business Source License 1.1}, la cual garantiza su uso gratuito para fines académicos, de investigación y del sector público, con una transición automática a la licencia \textit{Apache 2.0} 
en abril de 2029 \parencite{cabrera2025apimetro}.

\begin{figure}[htbp]
  \centering
  \figinclude[width=0.85\textwidth]{Figures/Cap3/PantallaInicialApimetroGithub.png}
  \caption{Estructura del c{\'o}digo fuente y control de versiones en GitHub.}
  \label{fig:screenshot-repo}
\end{figure}

\paragraph{3. Comunidad de Desarrollo (Discussions):} Espacio dedicado a la colaboración abierta y la gobernanza del proyecto. A través de este canal, la comunidad puede proponer nuevas funcionalidades, reportar incidencias en los trazos 
de transporte o discutir la integración de nuevos sistemas de la Zona Metropolitana, fomentando una construcción colectiva de la infraestructura de datos.

\begin{figure}[htbp]
  \centering
  \figinclude[width=0.85\textwidth]{Figures/Cap3/PantallaInicialComunidadApimetro.png}
  \caption{Foro de discusi{\'o}n y colaboraci{\'o}n de la comunidad Apimetro.}
  \label{fig:screenshot-discussions}
\end{figure}

\paragraph{4. Wiki de Apimetro:} Centro de documentación técnica profunda donde se detallan las guías de implementación, el diccionario de datos de las tablas \texttt{SQL} y los procedimientos para replicar el motor de datos de forma local 
utilizando el archivo \texttt{seed.sql}.

\begin{figure}[htbp]
  \centering
  \figinclude[width=0.85\textwidth]{Figures/Cap3/PantallawikiApimetro.png}
  \caption{Repositorio de documentaci{\'o}n t{\'e}cnica y gu{\'i}as de implementaci{\'o}n.}
  \label{fig:screenshot-wiki}
\end{figure}

Esta infraestructura digital asegura que la investigación no termine con la entrega de la tesis, sino que perdure como un \textbf{bien público tecnológico} capaz de evolucionar y adaptarse a las dinámicas cambiantes de la movilidad 
urbana en México \parencite{fortin2016innovative}.

%% ── 4-Capas-Codigo/4-2-VFTModel.tex
\subsection*{4 2 VFTModel}
% !TeX root = ../tesis.tex
%%%%%%%%%%%%%%%%%%%%%%%%%%%%%%%%%%%%%%%%%%%%%%%%%%%%%%%%%%%%%%%%%%%
% SECCIÓN 4.2 — VFTModel: Vanishing Fig-Tree Model
%%%%%%%%%%%%%%%%%%%%%%%%%%%%%%%%%%%%%%%%%%%%%%%%%%%%%%%%%%%%%%%%%%%

% =============================================================================
%  4.2 — VISIÓN GENERAL
% =============================================================================
\section{VFTModel: El Motor Analítico de la Red}
\label{sec:vftmodel-intro}

\textbf{VFTModel} (\termino{Vanishing Fig-Tree Model}) es el segundo componente
de la arquitectura de código de esta investigación. Si \textbf{Apimetro} cumple
la función de fuente de datos —exponiendo la geometría de la red de transporte
mediante una \api{} \textsc{rest}—, VFTModel es el \textbf{motor analítico}
que consume esos datos, los transforma en una estructura matemática de grafos y
ejecuta los algoritmos que producen los indicadores topológicos presentados en
el Capítulo~\ref{ch:resultados-vft}.

\begin{figure}[htbp]
  \centering
  \figinclude[width=0.45\textwidth]{Figures/Cap4/VFTLogo1.png}
  \caption[Ícono Personal de VFTModel]{Ícono Creado para el proyecto VFTModel}
  \label{fig:VFTModel-icon}
  \fuentefigura{Elaboración propia.}
\end{figure}


\subsection{Contexto y Motivación}
\label{subsec:vft-contexto}

La \zmvm{} alberga una de las redes de transporte público más heterogéneas del
mundo: desde el Metro, que opera en vías exclusivas subterráneas a velocidades
comerciales de hasta 36~km/h, hasta los Corredores Concesionados que comparten
carril con el automóvil particular y absorben la saturación vial de la ciudad.
Esta diversidad opera simultáneamente como fortaleza y como desafío analítico.
Contabilizar kilómetros de línea o número de estaciones no es suficiente para
determinar si la red es eficiente: un sistema puede tener muchas estaciones y
aun así obligar a sus usuarios a recorrer trayectos tortuosos, o concentrar
toda la conectividad en grandes nodos de transbordo mientras las zonas periféricas
quedan sin alimentación capilar adecuada.

VFTModel nace como respuesta a esta necesidad: un motor analítico reproducible,
de código abierto y matemáticamente fundamentado, capaz de cuantificar la
eficiencia estructural de la red mediante indicadores topológicos derivados de
la teoría de grafos \parencite{newman2010networks}. El modelo transforma los
datos geográficos del sistema metropolitano en una estructura matemática de nodos
y aristas sobre la cual aplica algoritmos de análisis de redes, produciendo
métricas concretas y comparables que permiten evaluar la conectividad del sistema
e identificar nodos vulnerables.

\subsection{El Nombre: Vanishing Fig-Tree}
\label{subsec:vft-nombre}

La denominación \termino{Vanishing Fig-Tree} —árbol de higuera en desaparición \DIFdelbegin \DIFdel{— o
—}\DIFdelend \DIFaddbegin \DIFadd{o
}\DIFaddend Modelo de la Higuera Evanescente— es una metáfora del comportamiento topológico
de una red de transporte fragmentada.
\\
\termino{Vanishing Fig-Tree Model} se basa en dos referencias clave en diferentes ámbitos\DIFdelbegin \DIFdel{, 
mientras uno es literario }\DIFdelend \DIFaddbegin \DIFadd{:
una es literaria }\DIFaddend y describe la desesperanza de una situación crítica ante una decisión
tan trascendente\DIFdelbegin \DIFdel{. El otro }\DIFdelend \DIFaddbegin \DIFadd{, mientras que la otra }\DIFaddend describe de manera lírica el llamado \termino{Punto de No Retorno}:

\citaclave{
  \begin{center}
    ``From the tip of every branch, like a fat purple fig, a wonderful future beckoned and 
    winked. One fig was a husband and a happy home and children, and another fig was a famous 
    poet and another fig was a brilliant professor, and another fig was Europe and Africa and 
    South America, and another fig was Constantine and Socrates and Attila and a pack of other 
    lovers with queer names and offbeat professions, and another fig was an Olympic lady crew 
    champion, and beyond and above these figs were many more figs I couldn't quite make out.
    \\ \vspace{1ex}
    I saw myself sitting in the crotch of this fig tree, starving to death, just because 
    I couldn't make up my mind which of the figs I would choose. I wanted each and every 
    one of them, but choosing one meant losing all the rest, and, as I sat there, unable 
    to decide, the figs began to wrinkle and go black, and, one by one, they plopped 
    to the ground at my feet.'' 

    \vspace{1.5ex}
    \hrule
    \vspace{1.5ex}

    ``De la punta de cada rama, como si de un grueso higo morado se tratara, pendía un 
    maravilloso futuro, señalado y rutilante. Un higo era un marido y un hogar feliz e 
    hijos y otro higo era un famoso poeta, y otro higo era un brillante profesor, y otro 
    higo era Europa y África y Sudamérica y otro higo era Constantino y Sócrates y Atila 
    y un montón de otros amantes con nombres raros y profesionales poco usuales, y otro 
    higo era una campeona de equipo olímpico de atletismo, y más allá y por encima de 
    aquellos higos había muchos más higos que no podía identificar claramente.
    \\ \vspace{1ex}
    Me vi a mí misma sentada en la bifurcación de ese árbol de higos, muriéndome de 
    hambre sólo porque no podía decidir cuál de los higos escoger. Quería todos y cada 
    uno de ellos, pero elegir uno significaba perder el resto, y, mientras yo estaba allí 
    sentada, incapaz de decidirme, los higos empezaron a arrugarse y a tornarse negros y, 
    uno por uno, cayeron al suelo, a mis pies.'' \\
    \vspace{1ex}
    \parencite{plath2019campana}
  \end{center}
}

\citaclave{
  \begin{center}
    ``My life ain't no holiday \\
    I've been through the point of no return \\
    I've seen what a man can do \\
    I've seen all the hate of a woman too.'' 

    \vspace{1.5ex}
    \hrule
    \vspace{1.5ex}

    ``Mi vida no es ninguna vacación \\
    He cruzado el punto de no retorno \\
    He visto lo que un hombre puede hacer \\
    También he visto todo el odio de una mujer.'' \\
    \vspace{1ex}
    \parencite{neworder_vanishingpoint_1989}
  \end{center}
}

Una higuera no muere cuando pierde sus ramas más delgadas. El tronco y las ramas
principales permanecen, pero su capacidad de expansión y de alimentación colapsa
silenciosamente: las raíces capilares dejan de nutrir las zonas más alejadas, y
el árbol entra en una decadencia progresiva que no es visible desde el exterior
hasta que es demasiado tarde. De manera análoga, una red de transporte puede
mantener sus corredores troncales operativos con plena visibilidad pública,
mientras sus rutas de alimentación capilar desaparecen o se deterioran sin que
ese deterioro aparezca en los indicadores convencionales de pasajeros
transportados o kilómetros de línea. Las zonas que pierden su conexión capilar
quedan aisladas funcionalmente del resto del sistema, aunque sigan figurando en
los mapas oficiales.

El modelo propone medir precisamente esta fragilidad estructural mediante
indicadores que capturen tanto la conectividad interna de los nodos —la
\termino{Fuerza Capilar}— como la eficiencia de los recorridos reales frente a
la distancia teórica mínima —el \termino{Factor de Desviación}—, y el impacto
que tendría la pérdida de un nodo crítico sobre el sistema en su conjunto.

% =============================================================================
%  4.2.2 — ARQUITECTURA
% =============================================================================
\section{Arquitectura del Sistema}
\label{sec:vft-arquitectura}

\subsection{Tres Capas con Responsabilidades Delimitadas}
\label{subsec:vft-capas}

\begin{figure}[htbp]
    \centering
    \DIFdelbeginFL %DIFDELCMD < \figinclude[width=0.45\textwidth]{Figures/Cap4/EsquemaVFTArquitecture.png}
%DIFDELCMD <   \caption[Aquitectura esquemática VFTModel]{%%%
\DIFdelFL{Esquema creado para el }\DIFdelendFL \DIFaddbeginFL \figplaceholder{Figures/TiKz_Libraries/Diagramas/vft-arquitectura-ports-adapters}
    \caption[Arquitectura \textit{Ports and Adapters} de VFTModel]{\DIFaddFL{Diagrama de la
    arquitectura }\termino{Ports and Adapters} \DIFaddFL{de }\DIFaddendFL VFTModel\DIFaddbeginFL \DIFaddFL{. El Dominio central
    (}\texttt{\DIFaddFL{src/core/}}\DIFaddFL{) contiene los algoritmos topológicos y permanece aislado
    de las tecnologías externas. Los tres adaptadores gestionan la comunicación
    con actores externos: API REST, herramientas de visualización y capa de
    persistencia. Todo el sistema corre dentro de un contenedor Docker.}\DIFaddendFL }
    \label{fig:VFTModel-arquitectura}
    \DIFdelbeginFL %DIFDELCMD < \fuentefigura{Elaboración propia con Gemini.}
%DIFDELCMD < %%%
\DIFdelendFL \DIFaddbeginFL \fuentefigura{Elaboración propia.}
\DIFaddendFL \end{figure}

VFTModel está organizado bajo el patrón de \termino{arquitectura hexagonal}
(\termino{Ports and Adapters}). La idea central es que el núcleo matemático del
sistema —los algoritmos que calculan los indicadores topológicos— no debe depender
de cómo lleguen los datos ni de cómo se presenten los resultados. Esta
independencia se logra separando el código en tres capas concéntricas:

\begin{description}
    \item[\textbf{Capa de \api{}} (\texttt{src/api/})]
    Recibe las peticiones \textsc{http}, valida los datos de entrada mediante
    contratos de tipo estrictos y devuelve respuestas \textsc{json}. Es la única
    capa visible para quien consume el sistema desde el exterior.

    \item[\textbf{Capa de Dominio} (\texttt{src/core/})]
    Contiene los algoritmos topológicos, los modelos matemáticos y los servicios
    de construcción del grafo. Aquí reside toda la inteligencia analítica del
    modelo. No depende de cómo lleguen los datos ni de cómo se muestren.

    \item[\textbf{Capa de Infraestructura} (\texttt{src/infrastructure/})]
    Gestiona los clientes \textsc{http} para comunicarse con Apimetro, el acceso
    a archivos \geojson{} locales de respaldo y las colas de procesamiento
    asíncrono. Se encarga de obtener los datos, no de analizarlos.
\end{description}

La Figura~\ref{fig:vft-arquitectura} ilustra la jerarquía de las tres capas y
el sentido de las dependencias entre ellas.

\begin{figure}[H]
    \centering
    \DIFdelbeginFL %DIFDELCMD < \begin{tikzpicture}[
%DIFDELCMD <         capa/.style = {
%DIFDELCMD <             rectangle, draw=unamAzul, rounded corners=5pt,
%DIFDELCMD <             text width=11cm, align=left,
%DIFDELCMD <             minimum height=1.6cm, inner sep=10pt
%DIFDELCMD <         },
%DIFDELCMD <         flecha/.style = {Stealth-Stealth, thick, color=unamAzul!55}
%DIFDELCMD <     ]
%DIFDELCMD <         \node[capa, fill=unamAzul!5] (api) {
%DIFDELCMD <             \textbf{Capa de \api{}} \hfill \texttt{src/api/}\\
%DIFDELCMD <             {\small Peticiones \textsc{http} · Validación Pydantic · Respuestas \textsc{json}}
%DIFDELCMD <         };
%DIFDELCMD <         \node[capa, fill=unamAzul!12, below=0.45cm of api] (core) {
%DIFDELCMD <             \textbf{Capa de Dominio} \hfill \texttt{src/core/}\\
%DIFDELCMD <             {\small Algoritmos topológicos · Modelos matemáticos · Constructor del grafo}
%DIFDELCMD <         };
%DIFDELCMD <         \node[capa, fill=unamAzul!22, below=0.45cm of core] (infra) {
%DIFDELCMD <             \textbf{Capa de Infraestructura} \hfill \texttt{src/infrastructure/}\\
%DIFDELCMD <             {\small Clientes \textsc{http} · Archivos \geojson{} locales · Colas asíncronas}
%DIFDELCMD <         };
%DIFDELCMD <         \draw[flecha] (api)  -- (core);
%DIFDELCMD <         \draw[flecha] (core) -- (infra);
%DIFDELCMD <     \end{tikzpicture}
%DIFDELCMD <     %%%
\DIFdelendFL \DIFaddbeginFL \figplaceholder{Figures/TiKz_Libraries/Diagramas/vft-arquitectura-hexagonal}
    \DIFaddendFL \caption[Arquitectura hexagonal de VFTModel]{Las tres capas de la arquitectura
    hexagonal de VFTModel. Las flechas bidireccionales indican dependencias
    controladas: cada capa solicita servicios a la inferior; la inferior nunca
    dicta qué calcular a las superiores.}
    \label{fig:vft-arquitectura}
    \fuentefigura{Elaboración propia.}
\end{figure}

El beneficio práctico de esta separación es la reproducibilidad: si el formato
de datos de Apimetro cambia, o si los resultados se requieren en un visor
cartográfico distinto, solo se modifica la capa correspondiente sin alterar el
núcleo matemático.

\subsection{Flujo de Procesamiento}
\label{subsec:vft-flujo}

El procesamiento de una solicitud de análisis sigue una cadena de seis
transformaciones sucesivas, desde la ingesta de datos geográficos crudos hasta
la entrega de métricas finales, como se esquematiza en la
Figura~\ref{fig:vft-flujo}.

\begin{figure}[H]
    \centering
    \DIFdelbeginFL %DIFDELCMD < \resizebox{\linewidth}{!}{%
%DIFDELCMD <     \begin{tikzpicture}[
%DIFDELCMD <         node distance = 0.5cm and 0.35cm,
%DIFDELCMD <         caja/.style = {
%DIFDELCMD <             rectangle, draw=unamAzul, fill=unamAzul!7,
%DIFDELCMD <             text width=1.95cm, align=center,
%DIFDELCMD <             minimum height=1.3cm, rounded corners=4pt,
%DIFDELCMD <             font=\small
%DIFDELCMD <         },
%DIFDELCMD <         flecha/.style = {-Stealth, thick, color=unamAzul!70},
%DIFDELCMD <         etiq/.style = {
%DIFDELCMD <             font=\scriptsize\itshape, color=gray,
%DIFDELCMD <             align=center, text width=1.95cm
%DIFDELCMD <         }
%DIFDELCMD <     ]
%DIFDELCMD <         \node[caja] (apy) {\textbf{Apimetro}\\{\tiny Go Backend}};
%DIFDELCMD <         \node[caja, right=of apy]  (val) {\textbf{Validación}\\{\tiny Pydantic}};
%DIFDELCMD <         \node[caja, right=of val]  (grf) {\textbf{Constructor}\\{\tiny del Grafo}};
%DIFDELCMD <         \node[caja, right=of grf]  (imp) {\textbf{Impedancia}\\{\tiny Motor VFT}};
%DIFDELCMD <         \node[caja, right=of imp]  (ana) {\textbf{Analizadores}\\{\tiny Topológicos}};
%DIFDELCMD <         \node[caja, right=of ana]  (res) {\textbf{Resultados}\\{\tiny JSON / CSV}};
%DIFDELCMD < 

%DIFDELCMD <         \node[etiq, below=0.18cm of apy]  {Estaciones\\y Líneas\\\geojson{}};
%DIFDELCMD <         \node[etiq, below=0.18cm of val]  {Contratos\\de datos};
%DIFDELCMD <         \node[etiq, below=0.18cm of grf]  {Grafo\\NetworkX};
%DIFDELCMD <         \node[etiq, below=0.18cm of imp]  {Pesos en\\minutos};
%DIFDELCMD <         \node[etiq, below=0.18cm of ana]  {8 Indicadores};
%DIFDELCMD <         \node[etiq, below=0.18cm of res]  {{\upshape\api{}} REST};
%DIFDELCMD < 

%DIFDELCMD <         \draw[flecha] (apy) -- (val);
%DIFDELCMD <         \draw[flecha] (val) -- (grf);
%DIFDELCMD <         \draw[flecha] (grf) -- (imp);
%DIFDELCMD <         \draw[flecha] (imp) -- (ana);
%DIFDELCMD <         \draw[flecha] (ana) -- (res);
%DIFDELCMD <     \end{tikzpicture}}%%%
\DIFdelendFL \DIFaddbeginFL \resizebox{\linewidth}{!}{%
    \figplaceholder{Figures/TiKz_Libraries/Diagramas/vft-flujo-procesamiento}}\DIFaddendFL %
    \caption[Flujo de procesamiento de VFTModel]{Cadena de transformación de
    VFTModel: desde los datos geográficos crudos de Apimetro hasta los
    indicadores topológicos publicados vía \api{} \textsc{rest}.}
    \label{fig:vft-flujo}
    \fuentefigura{Elaboración propia.}
\end{figure}

% =============================================================================
%  4.2.3 — ENTRADAS Y VALIDACIÓN
% =============================================================================
\section{Entradas y Validación de Datos}
\label{sec:vft-entradas}

\subsection{Formato de Ingesta: GeoJSON}
\label{subsec:vft-geojson}

VFTModel se alimenta de los datos geográficos producidos por Apimetro en formato
\geojson{} \parencite{butler2016geojson}. Cada petición devuelve una
\texttt{FeatureCollection} con dos tipos de elementos:

\begin{itemize}
    \item \textbf{Estaciones} (campo \texttt{tipo\_entidad} = \texttt{estacion}): representadas
    como puntos geográficos (\texttt{Point}), con coordenadas GPS y atributos
    operativos —nombre, sistema, jerarquía de transporte, alcaldía y si es un
    nodo \textsc{cetram}—.
    \item \textbf{Rutas} (campo \texttt{tipo\_entidad}\allowbreak{} = \texttt{ruta}): representadas como
    trazos lineales (\texttt{MultiLineString}), con las métricas que el motor de
    impedancia requiere —velocidad operacional, frecuencia de servicio, tipo de
    derecho de vía y sentido de circulación—.
\end{itemize}

Antes de procesar cualquier elemento, VFTModel valida que todos los campos
requeridos existan y sean válidos, como se describe en la siguiente sección.

\subsection{Validación Estricta con Pydantic}
\label{subsec:vft-pydantic}

Cada \texttt{Feature} recibido es sometido a una validación estricta mediante
\textbf{Pydantic}, una biblioteca de Python que verifica que los datos cumplan
exactamente los contratos de tipo definidos en el modelo. Si algún campo es
inválido —por ejemplo, un sistema de transporte no reconocido o un valor de
velocidad negativo—, el sistema lo rechaza antes de que contamine los cálculos.
Este filtro está implementado en \texttt{src/api/schemas/schemas.py}.

El contrato verifica ocho campos críticos: \texttt{sistema} (Enum de 12 valores), \texttt{tipo\_entidad} (\texttt{estacion} o \texttt{ruta}), \texttt{jerarquia\_transporte}, \texttt{derecho\_de\_via}, \texttt{velocidad\_promedio\_kmh}, \texttt{frecuencia\_minutos}, \texttt{sentido} y \texttt{es\_cetram}. La especificación completa del contrato de datos se presenta en el Cuadro~\ref{tab:vft-campos-validados} del Anexo~\ref{anx:vft-validacion}.

\subsection{Catálogo de Sistemas de Transporte}
\label{subsec:vft-sistemas}

VFTModel reconoce exactamente doce sistemas de transporte metropolitano. Cuando los datos de velocidad o frecuencia no llegan desde Apimetro, el modelo aplica valores de \termino{fallback} calibrados a la operación real de cada sistema. Los sistemas estructurales de mayor peso en el análisis son el STC Metro (36 km/h, vía exclusiva), Metrobús (16.3 km/h, vía confinada), Tren Ligero (22 km/h, vía exclusiva), \textsc{rtp} (14 km/h, vía mixta) y Corredores Concesionados (11 km/h, vía mixta). El catálogo completo con claves de código y parámetros de \termino{fallback} se presenta en el Cuadro~\ref{tab:vft-sistemas} del Anexo~\ref{anx:vft-catalogo}.

El tipo de derecho de vía es el parámetro que determina el coeficiente de
fricción aplicado a cada arista del grafo, como se describe en la
Sección~\ref{subsec:vft-cf}.

% =============================================================================
%  4.2.4 — CONSTRUCCIÓN DEL GRAFO
% =============================================================================
\section{Construcción del Grafo Topológico}
\label{sec:vft-grafo}

\subsection{De la Geometría GeoJSON a NetworkX}
\label{subsec:vft-geom-networkx}

El módulo central del sistema es la clase \texttt{\textbf{VFTGraphBuilder}}\allowbreak{}
(\path{src/core/services/graph_builder.py}), responsable de transformar las
geometrías \geojson{} en un \textbf{dígrafo\DIFdelbegin \DIFdel{dirigido}\DIFdelend } representado mediante la
biblioteca \textbf{NetworkX} de Python \parencite{hagberg2008networkx}. La
construcción opera en dos fases internas:

\begin{description}
    \item[\textbf{Fase 1 — Registro de estaciones (nodos):}]
    Cada estación del \geojson{} se convierte en un nodo del grafo,
    identificado por sus coordenadas geográficas (longitud, latitud) y
    enriquecido con sus atributos operativos: nombre, sistema de transporte,
    jerarquía, alcaldía y condición de nodo \textsc{cetram}.

    \item[\textbf{Fase 2 — Trazado de rutas (aristas):}]
    Para cada ruta del \geojson{}, el constructor recorre el
    \texttt{MultiLineString} segmento a segmento, acumulando distancia y
    detectando qué estaciones se encuentran sobre el trazo mediante una
    estructura de búsqueda espacial (\termino{KDTree}) con una tolerancia de
    50~m. Por cada par de estaciones consecutivas detectadas, crea una arista
    con la distancia real acumulada a lo largo del trazo —no la distancia en
    línea recta entre las estaciones—. Esta distinción es crítica para sistemas
    como \textsc{rtp} y Corredores Concesionados, cuyos trazos urbanos pueden
    ser hasta un 40~\% más largos que la cuerda directa.
\end{description}
\subsection{Snapping Lógico Geográfico}
\label{subsec:vft-snapping}

En el modo de integración realista (\texttt{REALISTIC\_INTEGRATION}), el
constructor agrega aristas adicionales de \textbf{transbordo peatonal} entre
estaciones de distintos sistemas que se encuentran dentro de una distancia
configurable. Este proceso —denominado \termino{Snapping Lógico Geográfico}—
representa el comportamiento real del pasajero, que puede caminar de una
estación a otra aunque no estén formalmente conectadas en los datos oficiales.

% --- BLOQUE DE IMÁGENES DEL PROCESO DE SNAPPING ---
\begin{figure}[H]
    \centering
    \begin{subfigure}[b]{0.32\textwidth}
        \centering
        \figinclude[width=\textwidth]{Figures/Cap4/SnappingProcess1.png}
        \caption{Nodos aislados}
        \label{fig:snap1}
    \end{subfigure}
    \hfill
    \begin{subfigure}[b]{0.32\textwidth}
        \centering
        \figinclude[width=\textwidth]{Figures/Cap4/SnappingProcess2.png}
        \caption{Evaluación de umbral}
        \label{fig:snap2}
    \end{subfigure}
    \hfill
    \begin{subfigure}[b]{0.32\textwidth}
        \centering
        \figinclude[width=\textwidth]{Figures/Cap4/SnappingProcess3.png}
        \caption{Arista de transbordo}
        \label{fig:snap3}
    \end{subfigure}
    \caption[Secuencia del Snapping Lógico Geográfico]
    {Representación visual del proceso de \termino{Snapping Lógico Geográfico}. Las estaciones 
    cercanas de distintos sistemas se conectan mediante aristas de transbordo si cumplen con 
    la tolerancia de distancia establecida.}
    \label{fig:proceso-snapping}
    \fuentefigura{Elaboración propia (VFTModel, 2026).}
\end{figure}
% --------------------------------------------------

La tolerancia de transbordo se selecciona a partir de seis umbrales estadísticos
derivados del análisis de las distancias reales entre estaciones de la red
metropolitana, definidos bajo la constante \texttt{STATISTICAL\_THRESHOLDS} en
\texttt{src/core/services/graph\_builder.py}.

\begin{table}[H]
    \centering
    \caption[Umbrales estadísticos de transbordo peatonal en VFTModel]{Umbrales
    estadísticos de tolerancia para el \termino{Snapping Lógico Geográfico}.
    El umbral \texttt{Q1} (85~m) es el recomendado para integración conservadora;
    \texttt{Q2} (180~m) integra la mayoría de los \textsc{cetram}s formales.}
    \label{tab:vft-umbrales}
    \begin{tabularx}{\textwidth}{|l|r|X|}
        \hline
        \textbf{Clave} & \textbf{Distancia (m)} & \textbf{Interpretación} \\
        \hline
        \texttt{MIN}  &  15 & Mismo recinto arquitectónico \\
        \texttt{Q1}   &  85 & Transbordo rápido y directo (integración conservadora) \\
        \texttt{Q2}   & 180 & Transbordo promedio real; integra la mayoría de \textsc{cetram}s \\
        \texttt{MEAN} & 245 & Promedio matemático (sesgado por valores atípicos) \\
        \texttt{Q3}   & 420 & Límite máximo de caminata forzada aceptable \\
        \texttt{MAX}  & 880 & Caso atípico; posible falso transbordo \\
        \hline
    \end{tabularx}
    \fuentefigura{Elaboración propia con base en \texttt{graph\_builder.py}
    (VFTModel, 2026).}
\end{table}

El grafo de la \zmvm{} en configuración \texttt{REALISTIC\_\allowbreak{}INTEGRATION Q1}
resultante contiene \textbf{11{,}115 nodos topológicos} y
\textbf{45{,}679 aristas} —25{,}197 de tránsito y 20{,}482 de transbordo
peatonal generadas algorítmicamente—, tras eliminar 1{,}562 aristas que
superaron el umbral de falsa conexión de 5~km.

Los registros de construcción completos de la corrida del 2026-04-18, que documentan cada fase del proceso, se presentan en el Anexo~\ref{anx:vft-codigo}.

\subsection{Modos de Operación}
\label{subsec:vft-modos}

\texttt{VFTGraphBuilder} ofrece dos modos de construcción del grafo, que
determinan si se generan o no las aristas de transbordo peatonal:

\begin{description}
    \item[\texttt{STRICT\_TOPOLOGY}:]
    Construye el grafo matemático puro: únicamente las aristas de tránsito
    definidas explícitamente por las rutas del \geojson{}. No agrega conexiones
    entre sistemas distintos. Útil para analizar la topología interna de cada
    línea en aislamiento o para verificar que el grafo refleja exactamente los
    trazos geográficos.

    \item[\texttt{REALISTIC\_INTEGRATION}:]
    Además de las aristas de tránsito, genera las aristas de transbordo peatonal
    entre estaciones próximas de distintos sistemas. El costo de cada arista de
    transbordo incluye el tiempo de caminata (a 5~km/h, límite superior del
    \parencite{sedatu2019manual}) más la espera promedio del siguiente servicio
    ($F_v / 2$). Este modo es el utilizado en todos los cálculos de indicadores
    presentados en el Capítulo~\ref{ch:resultados-vft}.
\end{description}

% =============================================================================
%  4.2.5 — MOTOR DE IMPEDANCIA
% =============================================================================
\section{Motor de Impedancia: La Ecuación VFT}
\label{sec:vft-impedancia}

Una vez construido el grafo topológico, el módulo \texttt{VFTImpedanceModel}
(\path{src/core/models/impedance.py}) recorre cada arista y le asigna un
\textbf{peso temporal en minutos}: el tiempo estimado de recorrido que los
algoritmos de caminos mínimos utilizarán para determinar la ruta más eficiente
entre dos estaciones.

\subsection{Distancia Geodésica: Fórmula de Haversine}
\label{subsec:vft-haversine}

Todos los cálculos de distancia del modelo utilizan la \textbf{fórmula de
Haversine}, que obtiene la longitud del arco más corto sobre la superficie de
una esfera entre dos puntos de latitud y longitud conocidos. A diferencia de
la distancia euclidiana plana, Haversine opera sobre coordenadas geográficas
reales y produce metros con un error menor al 0.3~\% para distancias urbanas
dentro de la \zmvm{}.

La fórmula se calcula en dos pasos:

\begin{equation}
    a = \sin^{2}\!\left(\frac{\Delta\varphi}{2}\right)
        + \cos\varphi_1 \cdot \cos\varphi_2
          \cdot \sin^{2}\!\left(\frac{\Delta\lambda}{2}\right)
    \label{eq:haversine-a}
\end{equation}

\begin{equation}
    d = 2R \cdot \arctan2\!\left(\sqrt{a},\;\sqrt{1-a}\right)
    \label{eq:haversine-d}
\end{equation}

\begin{itemize}
    \item $\varphi_1, \varphi_2$: latitudes de los dos puntos en radianes.
    \item $\Delta\varphi = \varphi_2 - \varphi_1$: diferencia de latitudes.
    \item $\Delta\lambda = \lambda_2 - \lambda_1$: diferencia de longitudes.
    \item $R = 6{,}371{,}000$~m: radio medio de la Tierra.
    \item $d$: distancia geodésica resultante en metros.
\end{itemize}

La implementación como método estático de \texttt{VFTImpedanceModel} es reutilizada por prácticamente todos los módulos del sistema. El código fuente completo se presenta en el Anexo~\ref{anx:vft-codigo}.

\subsection{Coeficiente de Fricción Vial}
\label{subsec:vft-cf}

El \textbf{Coeficiente de Fricción Vial} ($CF$) traduce el tipo de derecho de
vía de cada segmento en un penalizador multiplicativo sobre el tiempo de
traslado. Se obtiene mediante:

\begin{equation}
    CF = 1 + \alpha \cdot \beta_{\text{ZMVM}}
    \label{eq:vft-cf}
\end{equation}

\begin{itemize}
    \item $\alpha \in \{0.0, 0.2, 0.5, 1.0\}$: factor de exposición al tráfico
    según el tipo de derecho de vía, con referencia normativa en
    \parencite{sedatu2019manual}.
    \item $\beta_{\text{ZMVM}} = 0.759$: índice de saturación vial de la
    \zmvm{} (\termino{TomTom Traffic Index}).
\end{itemize}

La clase \texttt{\textbf{FrictionCalculator}} implementa este cálculo de forma aislada, aplicando el principio de vía mixta ($CF = 1.759$) como caso conservador por defecto. La implementación Python completa se presenta en el Anexo~\ref{anx:vft-codigo}.

\subsection{Ecuación VFT: Peso Temporal de Cada Arista}
\label{subsec:vft-ecuacion}

Con la distancia geodésica y el coeficiente de fricción disponibles, el método
\path{apply_impedance()} asigna el peso temporal a cada arista del grafo:

\begin{equation}
    w_{ij} = \frac{d_{ij}}{v_{ij}} \cdot CF_{ij}
    \label{eq:vft-peso}
\end{equation}

\begin{itemize}
    \item $w_{ij}$: peso de la arista en minutos (tiempo de viaje estimado entre
    estaciones $i$ y $j$).
    \item $d_{ij}$: distancia real del trazo entre $i$ y $j$ en metros, acumulada
    vértice a vértice sobre el \texttt{MultiLineString} durante la construcción
    del grafo.
    \item $v_{ij}$: velocidad comercial del sistema en metros por minuto:
    $v_{ij} = v_{\text{km/h}} \times (1000/60)$.
    \item $CF_{ij}$: Coeficiente de Fricción Vial de la arista
    (Ecuación~\ref{eq:vft-cf}).
\end{itemize}

La implementación Python del método \path{apply_impedance()} se presenta en el Anexo~\ref{anx:vft-codigo}.

Las aristas de transbordo peatonal quedan excluidas de este proceso porque
ya recibieron su peso durante la construcción del grafo:
$w_{\text{transbordo}} = t_{\text{caminata}} + F_v / 2$ (véase la
Sección~\ref{subsec:res-penalizacion-transferencia} del Capítulo~\ref{ch:resultados-vft}).
El atributo \texttt{weight} de cada arista —en minutos— es la entrada de todos
los algoritmos de caminos mínimos ejecutados en los indicadores de Fase~3.

% =============================================================================
%  4.2.6 — INDICADORES
% =============================================================================
\section{Indicadores Topológicos Implementados}
\label{sec:vft-indicadores}

VFTModel organiza sus indicadores en cuatro fases de cálculo progresivo,
cuya dependencia es estricta: las Fases~III y IV requieren un grafo completamente
ponderado —producto de las Fases~I y II— para ejecutar los algoritmos de caminos
mínimos. El Cuadro~\ref{tab:vft-indicadores-estado} presenta el mapa completo
de indicadores con su estado de implementación a la fecha de cierre del documento.

\begin{table}[H]
    \centering
    \caption[Indicadores topológicos de VFTModel y estado de implementación]{
    Indicadores topológicos de VFTModel, sus símbolos y estado de implementación.
    Los resultados de cada indicador se presentan en el
    Capítulo~\ref{ch:resultados-vft}.}
    \label{tab:vft-indicadores-estado}
    \begin{tabularx}{\textwidth}{|c|c|X|l|}
        \hline
        \textbf{Fase} & \textbf{\#} & \textbf{Indicador} & \textbf{Estado} \\
        \hline
        \multirow{4}{*}{\textbf{I}}
          & 1 & Validación de Impedancia                  & Implementado \\
          & 2 & Cobertura Espacial ($C$)                  & Implementado \\
          & 3 & Fuerza Capilar ($FC$)                     & Implementado \\
          & 4 & Factor de Desviación ($DF$)               & Implementado \\
        \hline
        \multirow{3}{*}{\textbf{II}}
          & 5 & Coeficiente de Fricción Vial ($CF$)       & Implementado \\
          & 6 & Penalización por Transferencia ($T_b$)    & Parcial      \\
          & 7 & Node Score ($\text{VFT\_score}$)          & Planificado  \\
        \hline
        \multirow{2}{*}{\textbf{III}}
          & 8 & Tiempo de Viaje Promedio ($T$)            & Pendiente    \\
          & 9 & Centralidad de Intermediación ($B(v)$)   & Pendiente    \\
        \hline
        \textbf{IV}
          & 10 & Vulnerabilidad Estructural ($\Delta E$) & Pendiente    \\
        \hline
    \end{tabularx}
    \fuentefigura{Elaboración propia (VFTModel, 2026-04-28).}
\end{table}

Los indicadores de Fase~I son implementados en módulos independientes dentro de
\path{src/core/algorithms/topologicalIndicators/}: la clase
\texttt{SpatialCoverageAnalyzer} para la cobertura espacial, la clase
\texttt{CapillaryStrengthAnalyzer} para la fuerza capilar, y el submódulo
\texttt{detaurFactor/} —con sus módulos \texttt{engine.py} y
\texttt{orchestator.py}— para el factor de desviación.

% =============================================================================
%  4.2.7 — SALIDAS
% =============================================================================
\section{Salidas: API REST con FastAPI}
\label{sec:vft-salidas}

Los resultados de VFTModel se entregan a través de una \textbf{\api{} \textsc{rest}}
construida con \textbf{FastAPI}, expuesta localmente en \texttt{localhost:8000}
con documentación interactiva disponible en el endpoint \texttt{/docs}.

La \api{} expone cinco endpoints bajo \texttt{/api/v1/network/}: \texttt{build-auto} (POST, construye el grafo en caché), \texttt{spatial-coverage} (GET, cobertura por alcaldía), \texttt{topological/capillary-strength} (GET, ranking de fuerza capilar), \texttt{topological/geo-capillary} (GET, macro-hubs) y \texttt{topological/detour-factor} (GET, factor de desviación por par O-D). La especificación completa de los endpoints se presenta en el Cuadro~\ref{tab:vft-endpoints} del Anexo~\ref{anx:vft-api}.

Cada respuesta \textsc{json} puede contener tres tipos de datos:

\begin{itemize}
    \item \textbf{Tablas de métricas:} listas de nodos con su puntuación en el
    indicador calculado, ordenadas de mayor a menor relevancia.
    \item \textbf{Geometrías de ruta:} la trayectoria real sobre la red y la
    línea recta teórica entre origen y destino, ambas en formato \geojson{},
    listas para ser renderizadas en visores cartográficos.
    \item \textbf{Resúmenes estadísticos:} promedio, máximo y distribución del
    indicador sobre el conjunto de pares O-D analizados en el muestreo.
\end{itemize}

%% ── 4-Capas-Codigo/4-3-Visualizacion.tex
\subsection*{4 3 Visualizacion}
% !TeX root = ../tesis.tex
%%%%%%%%%%%%%%%%%%%%%%%%%%%%%%%%%%%%%%%%%%%%%%%%%%%%%%%%%%%%%%%%%%%
% SECCIÓN 4.3 — Transport-gis: Motor de Visualización
%%%%%%%%%%%%%%%%%%%%%%%%%%%%%%%%%%%%%%%%%%%%%%%%%%%%%%%%%%%%%%%%%%%

% =============================================================================
%  4.3 — VISIÓN GENERAL
% =============================================================================
\section{Transport-gis: El motor de Visualización}
\label{sec:transport-gis-intro}

\textbf{Transport-gis-zmvm-mjg} (\termino{Transport GIS — Zona Metropolitana del Valle de México, M.J. Garibelt}) 
es el tercer y último componente de la arquitectura de código de esta investigación. Mientras \textbf{Apimetro} 
opera como el motor de datos unificado —responsable de la ingesta, estandarización y exposición de la red de transporte 
como servicio— y \textbf{VFTModel} constituye el núcleo analítico que calcula los indicadores de desempeño de la red, 
\textbf{Transport-gis} representa la capa de visualización: el sistema que traduce los resultados del modelo en 
representaciones cartográficas interpretables para el análisis urbanístico.

Al igual que los componentes precedentes, la denominación del sistema proviene de la obra literaria ficcional 
\textit{Saga Diarios, tomo~2}, de autoría del investigador. En dicha novela, la urbanista ficcional 
\textit{Mildred Julietta Cordero Garibelt} —referida por su acrónimo \textit{M.J.G}— es la responsable de la cartografía 
y la planeación espacial del sistema de movilidad de \textit{Ciudad Capital}, ucronía de la Ciudad de México en la década de 1960. 
La trama sitúa a este personaje en el centro de eventos políticos relevantes de la historia mexicana del siglo~\textsc{xx} 
\parencite{lopezmateos2023comoelcantodelcisne}, cuya intervención en las decisiones sobre infraestructura y movilidad urbana articula 
el argumento central de la novela. El paralelismo con la presente investigación es deliberado: así como \textit{M.J.G} produce 
los mapas que hacen legible el sistema de movilidad de \textit{Ciudad Capital}, \textbf{Transport-gis} genera las representaciones 
cartográficas que permiten interpretar los resultados del \textbf{VFTModel} sobre la red metropolitana real de la \zmvm.

\subsection{Transport-gis: Su arquitectura y funcionamiento}
\label{subsec:transport-gis-arquitectura-funcionamiento}

\textbf{Transport-gis} se organiza en cuatro módulos funcionales interdependientes: un módulo de \textbf{análisis} (\texttt{analysis/}), que gestiona la consulta a los servicios 
externos y el procesamiento de indicadores; un módulo de \textbf{datos} (\texttt{data/processed/}), que almacena las capas geoespaciales en formato \termino{GeoPackage} (\texttt{.gpkg}) 
bajo control de Git~\textsc{lfs}; un módulo \textbf{cartográfico} (\texttt{maps/}), que agrupa los proyectos QGIS, las plantillas tipográficas y los archivos de exportación; y un 
módulo de \textbf{orquestación} (\texttt{pipeline/}), que coordina la ejecución secuencial del proceso completo. La Figura~\ref{fig:tg-pipeline} ilustra el flujo de datos 
entre estos módulos y los componentes externos de la arquitectura de la investigación.

% =============================================================================
\subsubsection{Flujo de producción de datos}
\label{subsubsec:tg-flujo-datos}
% =============================================================================

El proceso de producción cartográfica parte de dos fuentes primarias externas: los \textit{feeds} \gtfs{} de la \textsc{semovi} (2024), que contienen los horarios, 
rutas y paradas de los sistemas de transporte de la \cdmx{}; y la cartografía base del \textsc{inegi} (2023), que provee los límites político-administrativos y la 
infraestructura vial de la \zmvm{}. Estas fuentes alimentan a \textbf{Apimetro} —que las transforma en una \api{} \textsc{rest} consumible con la topología del 
grafo de la red— y a \textbf{VFTModel}, que las procesa para calcular los indicadores de eficiencia nodal.

Dentro del repositorio de \textbf{Transport-gis}, el script \texttt{apimetro\_client.py} consulta la \api{} de \textbf{Apimetro} para recuperar la topología del 
grafo y los atributos de cada modo de transporte, mientras que \texttt{vft\_runner.py} ejecuta las rutinas del \textbf{VFTModel} para producir los indicadores por nodo en 
formato \texttt{.csv}. El orquestador \texttt{pipeline/run\_pipeline.py} gestiona la secuencia completa de transformaciones hasta exportar las capas procesadas al 
directorio \texttt{data/processed/}. Los fragmentos de código de estos componentes se detallarán en un anexo técnico.\pendiente{Crear Anexo — Documentación Técnica Transport-gis}

\begin{figure}[htbp]
\centering
\DIFdelbeginFL %DIFDELCMD < \begin{tikzpicture}[
%DIFDELCMD <   srcbox/.style = {rectangle, rounded corners=3pt, draw=gray!55, fill=gray!8,
%DIFDELCMD <                    minimum width=3.0cm, minimum height=0.90cm,
%DIFDELCMD <                    align=center, font=\small},
%DIFDELCMD <   extbox/.style = {rectangle, rounded corners=3pt, draw=unamAzul, fill=unamAzul!10,
%DIFDELCMD <                    minimum width=3.2cm, minimum height=0.90cm,
%DIFDELCMD <                    align=center, font=\small},
%DIFDELCMD <   intbox/.style = {rectangle, rounded corners=3pt, draw=unamAzul!60, fill=unamGrisClaro,
%DIFDELCMD <                    minimum width=6.4cm, minimum height=1.05cm,
%DIFDELCMD <                    align=center, font=\small},
%DIFDELCMD <   outbox/.style = {rectangle, rounded corners=3pt, draw=unamOro, fill=unamOro!15,
%DIFDELCMD <                    minimum width=3.2cm, minimum height=0.90cm,
%DIFDELCMD <                    align=center, font=\small},
%DIFDELCMD <   ->, >=Stealth, thick
%DIFDELCMD < ]
%DIFDELCMD < %% — Fuentes externas —
%DIFDELCMD < \node[srcbox] (gtfs)  at (0.0, 9.2) {\gtfs{} SEMOVI~2024};
%DIFDELCMD < \node[srcbox] (inegi) at (5.4, 9.2) {INEGI~2023 (cartografía)};
%DIFDELCMD < 

%DIFDELCMD < %% — Componentes externos —
%DIFDELCMD < \node[extbox] (api) at (0.0, 7.5) {\textbf{Apimetro}\\{\footnotesize API REST (Go)}};
%DIFDELCMD < \node[extbox] (vft) at (5.4, 7.5) {\textbf{VFTModel}\\{\footnotesize (Python / NetworkX)}};
%DIFDELCMD < 

%DIFDELCMD < %% — Transport-gis: análisis —
%DIFDELCMD < \node[intbox] (scripts) at (2.7, 5.5)
%DIFDELCMD <   {{\footnotesize\texttt{apimetro\_client.py} $\cdot$ \texttt{vft\_runner.py}}\\
%DIFDELCMD <    {\footnotesize\texttt{pipeline/run\_pipeline.py}}};
%DIFDELCMD < 

%DIFDELCMD < %% — Transport-gis: datos —
%DIFDELCMD < \node[intbox] (gpkg) at (2.7, 4.0)
%DIFDELCMD <   {{\small\texttt{data/processed/}}\\
%DIFDELCMD <    {\footnotesize 10 capas \texttt{.gpkg} (Git~LFS)}};
%DIFDELCMD < 

%DIFDELCMD < %% — Transport-gis: QGIS —
%DIFDELCMD < \node[intbox] (qgis) at (2.7, 2.5)
%DIFDELCMD <   {{\small QGIS~3.x — \texttt{maps/projects/*.qgz}}\\
%DIFDELCMD <    {\footnotesize\texttt{maps/templates/*.qpt} (paletas UNAM)}};
%DIFDELCMD < 

%DIFDELCMD < %% — Transport-gis: exports —
%DIFDELCMD < \node[intbox] (exports) at (2.7, 1.0)
%DIFDELCMD <   {{\small\texttt{maps/exports/}}\\
%DIFDELCMD <    {\footnotesize\texttt{*.pdf} $\cdot$ \texttt{*.png}}};
%DIFDELCMD < 

%DIFDELCMD < %% — Destino final —
%DIFDELCMD < \node[outbox] (tesis) at (2.7, -0.5)
%DIFDELCMD <   {\textbf{Tesis LaTeX}\\{\footnotesize\texttt{Figures/}}};
%DIFDELCMD < 

%DIFDELCMD < %% — Flechas: fuentes → externos —
%DIFDELCMD < \draw[color=gray!60] (gtfs)  -- (api);
%DIFDELCMD < \draw[color=gray!60] (inegi) -- (vft);
%DIFDELCMD < 

%DIFDELCMD < %% — Flechas: externos → scripts (convergentes, con codo exterior) —
%DIFDELCMD < \draw[color=unamAzul!70] (api.south) -- ++(0,-0.55) -| (scripts.north west);
%DIFDELCMD < \draw[color=unamAzul!70] (vft.south) -- ++(0,-0.55) -| (scripts.north east);
%DIFDELCMD < 

%DIFDELCMD < %% — Flujo interno —
%DIFDELCMD < \draw[color=unamAzul!70] (scripts) -- (gpkg);
%DIFDELCMD < \draw[color=unamAzul!70] (gpkg)    -- (qgis);
%DIFDELCMD < \draw[color=unamAzul!70] (qgis)    -- (exports);
%DIFDELCMD < 

%DIFDELCMD < %% — Exports → LaTeX —
%DIFDELCMD < \draw[color=unamOro] (exports) -- (tesis);
%DIFDELCMD < 

%DIFDELCMD < %% — Recuadro del repositorio Transport-gis (en fondo) —
%DIFDELCMD < \begin{scope}[on background layer]
%DIFDELCMD <   \node[draw=unamAzul!40, dashed, rounded corners=6pt, fill=unamAzul!3,
%DIFDELCMD <         fit={(scripts)(gpkg)(qgis)(exports)},
%DIFDELCMD <         inner sep=14pt] (tgbox) {};
%DIFDELCMD < \end{scope}
%DIFDELCMD < %% — Label DENTRO del recuadro, esquina sup-izq —
%DIFDELCMD < \node[font=\small\bfseries\color{unamAzul}, anchor=north west]
%DIFDELCMD <   at ($(tgbox.north west) + (4pt,-3pt)$) {Transport-gis-zmvm-mjg};
%DIFDELCMD < 

%DIFDELCMD < \end{tikzpicture}
%DIFDELCMD < %%%
\DIFdelendFL \DIFaddbeginFL \figplaceholder{Figures/TiKz_Libraries/Diagramas/tg-pipeline-cartografico}
\DIFaddendFL \caption[Pipeline de producción cartográfica de Transport-gis]%
  {Flujo de datos del sistema \textbf{Transport-gis-zmvm-mjg}. 
  Los componentes externos \textbf{Apimetro} y \textbf{VFTModel} proveen la información primaria; 
  dentro del repositorio, los \textit{scripts} de análisis consolidan los datos en capas \texttt{.gpkg} 
  que QGIS transforma en los mapas de presentación exportados a la tesis.}
\label{fig:tg-pipeline}
\end{figure}

% =============================================================================
\subsubsection{Capas del sistema de información geográfica}
\label{subsubsec:tg-capas}
% =============================================================================

El módulo de datos mantiene diez capas \termino{GeoPackage} en \texttt{data/processed/}, 
todas proyectadas en \textsc{epsg}:32614 (UTM zona~14N), referencia cartográfica adecuada 
para la extensión de la \zmvm{}. Las capas se organizan en tres categorías temáticas, ilustradas en la 
Figura~\ref{fig:tg-capas}: (i)~\textbf{contexto territorial}, con los límites político-administrativos de la \cdmx{}, del 
Estado de México, el polígono de estudio y las variables de medio físico natural; (ii)~\textbf{red e infraestructura}, 
con las geometrías de los sistemas de transporte, la infraestructura vial y el trazado del Periférico; y 
(iii)~\textbf{indicadores del modelo VFT}, con las capas derivadas del cálculo de cobertura espacial, 
fuerza capilar y factor de desviación descritos en el Capítulo~\ref{ch:resultados-vft}.

\begin{figure}[htbp]
\centering
\DIFdelbeginFL %DIFDELCMD < \begin{tikzpicture}[
%DIFDELCMD <   hdr1/.style = {rectangle, rounded corners=3pt,
%DIFDELCMD <                  draw=gray!60, fill=gray!55,
%DIFDELCMD <                  minimum width=4.2cm, minimum height=0.75cm,
%DIFDELCMD <                  align=center, font=\small\bfseries\color{white}},
%DIFDELCMD <   lyr1/.style = {rectangle, rounded corners=2pt,
%DIFDELCMD <                  draw=gray!45, fill=gray!8,
%DIFDELCMD <                  minimum width=4.2cm, minimum height=0.65cm,
%DIFDELCMD <                  align=center, font=\footnotesize},
%DIFDELCMD <   hdr2/.style = {rectangle, rounded corners=3pt,
%DIFDELCMD <                  draw=unamAzul, fill=unamAzul!70,
%DIFDELCMD <                  minimum width=4.2cm, minimum height=0.75cm,
%DIFDELCMD <                  align=center, font=\small\bfseries\color{white}},
%DIFDELCMD <   lyr2/.style = {rectangle, rounded corners=2pt,
%DIFDELCMD <                  draw=unamAzul!50, fill=unamAzul!8,
%DIFDELCMD <                  minimum width=4.2cm, minimum height=0.65cm,
%DIFDELCMD <                  align=center, font=\footnotesize},
%DIFDELCMD <   hdr3/.style = {rectangle, rounded corners=3pt,
%DIFDELCMD <                  draw=unamOro!80, fill=unamOro!75,
%DIFDELCMD <                  minimum width=4.2cm, minimum height=0.75cm,
%DIFDELCMD <                  align=center, font=\small\bfseries\color{white}},
%DIFDELCMD <   lyr3/.style = {rectangle, rounded corners=2pt,
%DIFDELCMD <                  draw=unamOro!60, fill=unamOro!10,
%DIFDELCMD <                  minimum width=4.2cm, minimum height=0.65cm,
%DIFDELCMD <                  align=center, font=\footnotesize},
%DIFDELCMD <   node distance = 0.12cm
%DIFDELCMD < ]
%DIFDELCMD < %% — Columna 1: Contexto territorial —
%DIFDELCMD < \node[hdr1] (h1) at (0,0) {Contexto territorial};
%DIFDELCMD < \node[lyr1, below=of h1] (c1a) {\texttt{LimitesPoliticos.gpkg}};
%DIFDELCMD < \node[lyr1, below=of c1a] (c1b) {\texttt{PoligonosEdoMex.gpkg}};
%DIFDELCMD < \node[lyr1, below=of c1b] (c1c) {\texttt{PoligonoEstudio1.gpkg}};
%DIFDELCMD < \node[lyr1, below=of c1c] (c1d) {\texttt{MedioFisicoNatural.gpkg}};
%DIFDELCMD < 

%DIFDELCMD < %% — Columna 2: Red e infraestructura —
%DIFDELCMD < \node[hdr2] (h2) at (4.6,0) {Red e infraestructura};
%DIFDELCMD < \node[lyr2, below=of h2] (c2a) {\texttt{Transporte.gpkg}};
%DIFDELCMD < \node[lyr2, below=of c2a] (c2b) {\texttt{InfraEstructura.gpkg}};
%DIFDELCMD < \node[lyr2, below=of c2b] (c2c) {\texttt{Vialidades.gpkg}};
%DIFDELCMD < 

%DIFDELCMD < %% — Columna 3: Indicadores del modelo —
%DIFDELCMD < \node[hdr3] (h3) at (9.2,0) {Indicadores del modelo};
%DIFDELCMD < \node[lyr3, below=of h3] (c3a) {\texttt{Cobertura800m.gpkg}};
%DIFDELCMD < \node[lyr3, below=of c3a] (c3b) {\texttt{FuerzaCapilar.gpkg}};
%DIFDELCMD < \node[lyr3, below=of c3b] (c3c) {\texttt{FactorDesviacion.gpkg}};
%DIFDELCMD < \end{tikzpicture}
%DIFDELCMD < %%%
\DIFdelendFL \DIFaddbeginFL \figplaceholder{Figures/TiKz_Libraries/Diagramas/tg-capas-geopackage}
\DIFaddendFL \caption[Organización de las capas GeoPackage de Transport-gis]%
  {Las diez capas \termino{GeoPackage} del directorio \texttt{data/processed/}, organizadas por 
  categoría temática. Todas están proyectadas en \textsc{epsg}:32614 (UTM zona~14N) y versionadas con Git~\textsc{lfs}.}
\label{fig:tg-capas}
\end{figure}

% =============================================================================
\subsubsection{Proyectos cartográficos y plantillas institucionales}
\label{subsubsec:tg-proyectos}
% =============================================================================

La Figura~\ref{fig:tg-arbol} muestra la estructura de directorios completa de \textbf{Transport-gis}. El repositorio se articula
en cinco carpetas de primer nivel con responsabilidades bien delimitadas: \texttt{analysis/} centraliza los \textit{scripts} de
procesamiento y los cuadernos de exploración; \texttt{data/processed/} almacena las diez capas \textit{GeoPackage} versionadas en
Git~\textsc{lfs}; \texttt{maps/} agrupa los proyectos QGIS~(\texttt{.qgz}), las plantillas visuales~(\texttt{.qpt}) con las paletas
institucionales de la \textsc{unam} y el directorio de exportaciones; \texttt{pipeline/} contiene el orquestador del proceso; y
\texttt{setup/} distribuye los activos tipográficos para garantizar la reproducibilidad visual en distintos sistemas operativos.
Los mapas terminados se exportan en \texttt{.pdf} (alta resolución para la tesis) y \texttt{.png} (presentaciones del coloquio),
cerrando el ciclo descrito en la Figura~\ref{fig:tg-pipeline}.

\begin{figure}[htbp]
\centering
\DIFdelbeginFL %DIFDELCMD < \iffalse
\iffalse
\begin{forest}
%DIFDELCMD <   %%%
\DIFdelFL{for tree = }%DIFDELCMD < {
%DIFDELCMD <     %%%
\DIFdelFL{font   = }%DIFDELCMD < \footnotesize\ttfamily%%%
\DIFdelFL{,
    grow'  = 0,
    child anchor  = west,
    parent anchor = east,
    anchor = west,
    draw   = gray!45,
    fill   = gray!7,
    rounded corners = 2pt,
    inner xsep = 6pt,
    inner ysep = 3pt,
    l sep  = 20pt,
    s sep  = 5pt,
    edge path = }%DIFDELCMD < {
%DIFDELCMD <       \noexpand\path [%%%
\DIFdelFL{draw=gray!50, thin, rounded corners=2pt}%DIFDELCMD < ]
%DIFDELCMD <         %%%
\DIFdelFL{(!u.parent anchor) -- +(8pt,0) |- (.child anchor)
        }%DIFDELCMD < \forestoption{edge label}%%%
\DIFdelFL{;
    }%DIFDELCMD < }%%%
\DIFdelFL{,
  }%DIFDELCMD < }
%DIFDELCMD <   [%%%
\DIFdelFL{Transport-gis-zmvm-mjg/, fill=unamAzul!12, draw=unamAzul!70,
                             font=}%DIFDELCMD < \small\bfseries\ttfamily
%DIFDELCMD <     [%%%
\DIFdelFL{analysis/
      }%DIFDELCMD < [%%%
\DIFdelFL{notebooks/}%DIFDELCMD < ]
%DIFDELCMD <       [%%%
\DIFdelFL{scripts/
        }%DIFDELCMD < [%%%
\DIFdelFL{apimetro\_client.py,
          draw=none, fill=none, font=}%DIFDELCMD < \footnotesize\ttfamily\color{unamAzul}]
%DIFDELCMD <         [%%%
\DIFdelFL{vft\_runner.py,
          draw=none, fill=none, font=}%DIFDELCMD < \footnotesize\ttfamily\color{unamAzul}]
%DIFDELCMD <       ]
%DIFDELCMD <     ]
%DIFDELCMD <     [%%%
\DIFdelFL{data/processed/
      }%DIFDELCMD < [{%%%
\textit{\DIFdelFL{10 capas .gpkg}} %DIFAUXCMD
\DIFdelFL{\textnormal{(Git~LFS)}}%DIFDELCMD < }%%%
\DIFdelFL{,
        draw=none, fill=none, font=}%DIFDELCMD < \footnotesize\itshape\ttfamily]
%DIFDELCMD <     ]
%DIFDELCMD <     [%%%
\DIFdelFL{maps/, fill=unamOro!8, draw=unamOro!60
      }%DIFDELCMD < [%%%
\DIFdelFL{exports/
        }%DIFDELCMD < [{%%%
\textit{\DIFdelFL{*.pdf · *.png}}%DIFAUXCMD
%DIFDELCMD < }%%%
\DIFdelFL{,
          draw=none, fill=none, font=}%DIFDELCMD < \footnotesize\itshape\ttfamily]
%DIFDELCMD <       ]
%DIFDELCMD <       [%%%
\DIFdelFL{fonts/}%DIFDELCMD < ]
%DIFDELCMD <       [%%%
\DIFdelFL{projects/
        }%DIFDELCMD < [%%%
\textit{\DIFdelFL{*.qgz}}%DIFAUXCMD
\DIFdelFL{,
          draw=none, fill=none, font=}%DIFDELCMD < \footnotesize\itshape\ttfamily]
%DIFDELCMD <       ]
%DIFDELCMD <       [%%%
\DIFdelFL{templates/
        }%DIFDELCMD < [{%%%
\textit{\DIFdelFL{*.qpt}} %DIFAUXCMD
\DIFdelFL{\textnormal{(paletas UNAM)}}%DIFDELCMD < }%%%
\DIFdelFL{,
          draw=none, fill=none, font=}%DIFDELCMD < \footnotesize\itshape\ttfamily]
%DIFDELCMD <       ]
%DIFDELCMD <     ]
%DIFDELCMD <     [%%%
\DIFdelFL{pipeline/
      }%DIFDELCMD < [%%%
\DIFdelFL{run\_pipeline.py,
        draw=none, fill=none, font=}%DIFDELCMD < \footnotesize\ttfamily\color{unamAzul}]
%DIFDELCMD <     ]
%DIFDELCMD <     [%%%
\DIFdelFL{setup/
      }%DIFDELCMD < [%%%
\DIFdelFL{install\_fonts.sh,
        draw=none, fill=none, font=}%DIFDELCMD < \footnotesize\ttfamily\color{unamAzul}]
%DIFDELCMD <       [%%%
\DIFdelFL{install\_fonts.bat,
        draw=none, fill=none, font=}%DIFDELCMD < \footnotesize\ttfamily\color{unamAzul}]
%DIFDELCMD <     ]
%DIFDELCMD <   ]
%DIFDELCMD < \end{forest}
\fi
%DIFDELCMD < %%%
\DIFdelendFL \DIFaddbeginFL \figplaceholder{Figures/TiKz_Libraries/Diagramas/Arboles/tg-arbol-directorios}
\DIFaddendFL \caption[Estructura de directorios de Transport-gis-zmvm-mjg]%
  {Árbol de directorios del repositorio \textbf{Transport-gis-zmvm-mjg}. Las carpetas \texttt{maps/} (oro) y la raíz (azul UNAM)
  son los nodos centrales de la producción cartográfica. Los archivos en azul son \textit{scripts} ejecutables;
  los de fondo plano son activos generados o de datos.}
\label{fig:tg-arbol}
\end{figure}

La Figura~\ref{fig:mapa-red-esquematica} presenta el resultado de este proceso aplicado a la red de transporte completa de la \zmvm{}:
un mapa de síntesis que integra los doce sistemas reconocidos por el modelo, proyectados en el espacio metropolitano real. Este mapa
constituye la representación cartográfica de referencia del estudio y es el punto de partida para la interpretación de los indicadores
presentados en el Capítulo~\ref{ch:resultados-vft}. Las versiones de alta resolución a página completa de todos los mapas producidos
por \textbf{Transport-gis} se encuentran en el Anexo~\ref{anx:cartografia-presentacion}.

\begin{figure}[htbp]
  \centering
  \figinclude[width=\textwidth]{Figures/Mapas/Mapa_Red_Esquematica}
  \caption[Red multimodal de transporte de la ZMVM]{Red esquemática multimodal de la \zmvm{} generada por \textbf{Transport-gis-zmvm-mjg}.
    Incluye los doce sistemas de transporte reconocidos por el modelo, proyectados en EPSG:32614 (UTM zona~14N) sobre los límites político-administrativos
    de la \cdmx{} y el Estado de México. \fuentefigura{Elaboración propia con QGIS~3.x, plantillas institucionales \textsc{unam};
    datos: \textsc{gtfs} \textsc{semovi} 2024 e \textsc{inegi} 2023.}}
  \label{fig:mapa-red-esquematica}
\end{figure}


%% ── 5-Resultados/5-1-Introduccion-Resultados.tex
\subsection*{5 1 Introduccion Resultados}
% !TeX root = ../tesis.tex
%%%%%%%%%%%%%%%%%%%%%%%%%%%%%%%%%%%%%%%%%%%%%%%%%%%%%%%%%%%%%%%%%%%
% SECCIÓN 5.1 — LA RED METROPOLITANA COMO OBJETO DE ANÁLISIS
%%%%%%%%%%%%%%%%%%%%%%%%%%%%%%%%%%%%%%%%%%%%%%%%%%%%%%%%%%%%%%%%%%%
\section{La Red Metropolitana como Objeto de Análisis}
\label{sec:intro-resultados}

La construcción del grafo topológico de la \zmvm constituye el punto de partida
de todos los indicadores que se presentan en este capítulo. Para transformar la
geografía estática de las rutas de transporte en un modelo operativo y ruteable,
el \textbf{Modelo VFT} implementa una arquitectura de procesamiento en tres
capas progresivas, cuyo flujo se sintetiza en la Figura~\ref{fig:arq-capas}.

%DIF <  --- Figura: arquitectura en 3 capas ---
%DIF >  --- Figura: pipeline de procesamiento VFT ---
\begin{figure}[H]
    \centering
    \DIFdelbeginFL %DIFDELCMD < \begin{subfigure}[b]{0.32\textwidth}
%DIFDELCMD <         \centering
%DIFDELCMD <         \figinclude[width=\textwidth]{Figures/Cap5/arq_capa1.png}
%DIFDELCMD <         \caption{%%%
\DIFdelFL{Capa 1: Adquisición y procesamiento geométrico}\DIFdelendFL \DIFaddbeginFL \figplaceholder{Figures/TiKz_Libraries/Diagramas/vft-pipeline-procesamiento}
    \caption[Pipeline de procesamiento del Modelo VFT]{\DIFaddFL{Las cuatro etapas de
    transformación que convierten los datos brutos de }\texttt{\DIFaddFL{Apimetro}} \DIFaddFL{en el
    grafo topológico operativo}\DIFaddendFL . \DIFaddbeginFL \DIFaddFL{El }\termino{Snapping Lógico Geográfico} \DIFaddFL{mediante
    árbol }\texttt{\DIFaddFL{KDTree}} \DIFaddFL{reduce las 22}{\DIFaddFL{,}}\DIFaddFL{766 entidades originales a
    11}{\DIFaddFL{,}}\DIFaddFL{115 nodos canónicos; la ponderación final asigna impedancias temporales
    a cada una de las 45}{\DIFaddFL{,}}\DIFaddFL{679 aristas del grafo dirigido
    $\mathcal{G}(\mathcal{V},\mathcal{E},w)$.}\DIFaddendFL }
    \DIFdelbeginFL %DIFDELCMD < \label{fig:arq-capa1}
%DIFDELCMD <     \end{subfigure}
%DIFDELCMD <     \hfill
%DIFDELCMD <     \begin{subfigure}[b]{0.32\textwidth}
%DIFDELCMD <         \centering
%DIFDELCMD <         \figinclude[width=\textwidth]{Figures/Cap5/arq_capa2.png}
%DIFDELCMD <         %%%
%DIFDELCMD < \caption{%
{%DIFAUXCMD
\DIFdelFL{Capa 2: Modelado topológico y dinámico.}}
        %DIFAUXCMD
%DIFDELCMD < \label{fig:arq-capa2}
%DIFDELCMD <     \end{subfigure}
%DIFDELCMD <     \hfill
%DIFDELCMD <     \begin{subfigure}[b]{0.32\textwidth}
%DIFDELCMD <         \centering
%DIFDELCMD <         \figinclude[width=\textwidth]{Figures/Cap5/arq_capa3.png}
%DIFDELCMD <         %%%
%DIFDELCMD < \caption{%
{%DIFAUXCMD
\DIFdelFL{Capa 3: Evaluación y salida de indicadores.}}
        %DIFAUXCMD
%DIFDELCMD < \label{fig:arq-capa3}
%DIFDELCMD <     \end{subfigure}
%DIFDELCMD <     %%%
%DIFDELCMD < \caption[Arquitectura en tres capas del Modelo VFT]{%
{%DIFAUXCMD
\DIFdelFL{Representación esquemática
    de las tres capas de procesamiento del Modelo VFT: adquisición geométrica,
    modelado topológico con impedancia temporal, y evaluación de indicadores.}}
    %DIFAUXCMD
\DIFdelendFL \label{fig:arq-capas}
    \fuentefigura{Elaboración propia.}
\end{figure}

% --- Estadísticas del grafo ---
\subsection{Caracterización del Grafo Construido}
\label{subsec:estadisticas-grafo}

La ingesta de datos desde \texttt{Apimetro} proporcionó \textbf{22{,}766 entidades
espaciales} que, tras el proceso de \termino{Snapping Lógico Geográfico} mediante
un árbol \texttt{KDTree} con tolerancia de 50~m, se consolidaron en un grafo
dirigido con las siguientes dimensiones:

\begin{itemize}
    \item \textbf{11{,}115 nodos topológicos} (estaciones operativas, eliminando
    nodos fantasma y redundancias geométricas).
    \item \textbf{45{,}679 aristas}, de las cuales 25{,}197 corresponden a
    segmentos de servicio real y 20{,}482 a aristas de transbordo peatonal
    generadas algorítmicamente.
    \item \textbf{11 sistemas de transporte} representados: \textsc{Metro},
    Metrobús, Trolebús, Tren Ligero, Tren Suburbano, Cablebús, Mexicable,
    Mexibús, \textsc{rtp}, Corredores Concesionados y Pumabús.
    \item Distribución modal: \textbf{12\% infraestructura masiva estructurada}
    (Metro, MB, TL, CBB, SUB) y \textbf{88\% transporte de superficie}
    (RTP, CC, Trolebús, Mexibús, Mexicable).
\end{itemize}

La Figura~\ref{fig:paneles-grafo} ilustra las cuatro perspectivas de
visualización del grafo resultante: la infraestructura agregada, la conectividad
interestación, y la direccionalidad de ida y regreso.

% --- Figura: 4 paneles del grafo ---
\begin{figure}[H]
    \centering
    \begin{subfigure}[b]{0.48\textwidth}
        \centering
        \figinclude[width=\textwidth]{Figures/Cap5/panel1_infraestructura.png}
        \caption{Panel 1: Infraestructura agregada de la red.}
        \label{fig:panel1}
    \end{subfigure}
    \hfill
    \begin{subfigure}[b]{0.48\textwidth}
        \centering
        \figinclude[width=\textwidth]{Figures/Cap5/panel2_conectividad.png}
        \caption{Panel 2: Conectividad interestación (aristas).}
        \label{fig:panel2}
    \end{subfigure}

    \vspace{0.5em}

    \begin{subfigure}[b]{0.48\textwidth}
        \centering
        \figinclude[width=\textwidth]{Figures/Cap5/panel3_ida.png}
        \caption{Panel 3: Direccionalidad de servicio (sentido ida).}
        \label{fig:panel3}
    \end{subfigure}
    \hfill
    \begin{subfigure}[b]{0.48\textwidth}
        \centering
        \figinclude[width=\textwidth]{Figures/Cap5/panel4_regreso.png}
        \caption{Panel 4: Direccionalidad de servicio (sentido regreso).}
        \label{fig:panel4}
    \end{subfigure}
    \caption[Visualización del grafo topológico de la \zmvm]{Cuatro perspectivas
    de visualización del grafo dirigido construido sobre la red multimodal de la
    \zmvm. El proceso de Snapping Lógico redujo los nodos en aproximadamente
    un 90\% respecto a la representación geométrica cruda.}
    \label{fig:paneles-grafo}
    \fuentefigura{Elaboración propia con \texttt{VFTModel} (corrida 2026-04-18).}
\end{figure}

% --- Tabla resumen de indicadores ---
\subsection{Estructura de Indicadores y Dependencias por Fase}
\label{subsec:tabla-indicadores}

Los ocho indicadores que componen el análisis topológico se organizan en tres
fases cuya ejecución es secuencialmente dependiente, como se muestra en el
Cuadro~\ref{tab:resumen-indicadores-cap5}.

\begin{table}[H]
    \centering
    \caption[Resumen de indicadores del Modelo VFT por fase]{Resumen de los
    indicadores del Modelo VFT, su fase de cálculo, dependencias y estado
    de avance en la investigación.}
    \label{tab:resumen-indicadores-cap5}
    \begin{tabularx}{\textwidth}{|l|l|X|l|}
        \hline
        \textbf{Indicador} & \textbf{Fase} & \textbf{Dependencias} & \textbf{Estado} \\
        \hline
        Validación de Impedancia    & 1 & Grafo base + fallbacks            & Calculado \\
        Cobertura Espacial ($C$)    & 1 & Geometría de estaciones           & Calculado \\
        Fuerza Capilar ($FC$)       & 1 & Grafo base                        & Calculado \\
        Factor de Desviación ($DF$) & 1 & Grafo base + impedancia           & Calculado \\
        \hline
        Coef. de Fricción Vial ($CF$)  & 2 & Fase 1 completa                & Calculado \\
        Penalización por Transf. ($T_b$) & 2 & Fase 1 + $CF$              & Formulado \\
        Node Score ($VFT\_score$)   & 2 & Fase 1 + datos MV1/MV2            & Planificado \\
        \hline
        Tiempo de Viaje Prom. ($T$) & 3 & Fases 1 y 2 completas             & Pendiente \\
        Centralidad Interm. ($B$)   & 3 & Fases 1 y 2 + $T$                 & Pendiente \\
        Vulnerabilidad ($\Delta E$) & 3 & Fases 1, 2 y $B$                  & Pendiente \\
        \hline
    \end{tabularx}
    \fuentefigura{Elaboración propia.}
\end{table}

%% ── 5-Resultados/5-2-1-Fase1/5-2-1-1-Impedancia.tex
\subsection*{5 2 1 1 Impedancia}
% !TeX root = ../../tesis.tex
% =============================================================================
%  5.2.1.1 VALIDACIÓN DEL MOTOR DE IMPEDANCIA
% =============================================================================
\subsection{Validación del Motor de Impedancia Temporal}
\label{subsec:validacion-impedancia}

Para verificar la consistencia del \textbf{Motor de Impedancia VFT}, se calculó
la velocidad efectiva implícita de cada sistema tras aplicar el Coeficiente de
Fricción Vial ($C_f$) y se contrastó con los rangos operativos nominales
documentados por los organismos operadores de la \zmvm. La velocidad efectiva
se obtiene a partir de la Ecuación~\eqref{eq:velocidad-efectiva}:

\begin{equation}
    V_{\text{efectiva}} = \frac{V_{\text{fallback}}}{C_f}
    \label{eq:velocidad-efectiva}
\end{equation}

donde $V_{\text{fallback}}$ representa la velocidad comercial de flujo libre del
sistema y $C_f$ es el coeficiente de fricción según el tipo de derecho de vía.

\colorbox{azulRevision}{El Cuadro~\ref{tab:validacion-impedancia} sintetiza los resultados} de la
corrida del modelo del 16 de abril de 2026, ejecutada sobre
\textbf{25{,}197 segmentos de servicio real}.

\begin{table}[H]
    \centering
    \caption[Validación de velocidades efectivas por sistema de transporte]{
    Contraste entre la velocidad implícita post-fricción calculada por el motor
    VFT y los rangos operativos nominales de cada sistema. Corrida: 2026-04-16.}
    \label{tab:validacion-impedancia}
    \begin{tabular}{|l|r|r|c|c|}
        \hline
        \textbf{Sistema} & \textbf{$V$ modelo (km/h)} & \textbf{$V$ implícita (km/h)} &
        \textbf{Rango esperado} & \textbf{Estado} \\
        \hline
        \colorbox{azulRevision}{Tren Suburbano}   & 65.0 & 65.0 & [55--75]  & \checkmark\ OK \\
        \colorbox{azulRevision}{Metro (STC)}      & 36.0 & 36.0 & [30--42]  & \checkmark\ OK \\
        \colorbox{azulRevision}{Tren Ligero}      & 22.0 & 22.0 & [18--26]  & \checkmark\ OK \\
        \colorbox{azulRevision}{Trolebús}         & 18.0 & 13.0 & [14--22]  & $\approx$ Límite \\
        \colorbox{azulRevision}{Cablebús}         & 20.0 & 17.4 & [16--24]  & \checkmark\ OK \\
        \colorbox{azulRevision}{Mexicable}        & 20.0 & 17.4 & [16--24]  & \checkmark\ OK \\
        \colorbox{azulRevision}{Metrobús}         & 16.3 & 14.1 & [13--20]  & \checkmark\ OK \\
        \colorbox{azulRevision}{Mexibús}          & 16.3 & 14.1 & [13--20]  & \checkmark\ OK \\
        \colorbox{azulRevision}{\textsc{rtp}}     & 14.0 &  8.0 & [10--18]  & $\approx$ Revisar \\
        \colorbox{azulRevision}{Corredores Conc.} & 11.0 &  8.0 & [8--14]   & \checkmark\ OK \\
        \colorbox{azulRevision}{Pumabús}          & 11.0 &  8.0 & [8--14]   & \checkmark\ OK \\
        \hline
    \end{tabular}
    \fuentefigura{Elaboración propia con \texttt{VFTModel} (corrida 2026-04-16).
    Fuentes de referencia: STC Metro \textsc{cdmx}, Metrobús \textsc{cdmx},
    \textsc{rtp} \textsc{cdmx}, \textsc{ste} Trolebús.}
\end{table}

La Figura~\ref{fig:impedancia-sistemas} presenta la distribución completa de
velocidades efectivas y tiempos por segmento para cada sistema.

\begin{figure}[H]
    \centering
    \figinclude[width=\textwidth]{Figures/Cap5/fig01_impedancia.png}
    \caption[Caracterización de la impedancia temporal por sistema]{(a)~Velocidad
    efectiva post-fricción por sistema, coloreada según el tipo de derecho de vía.
    (b)~Distribución de tiempos de viaje por segmento interestación (segmentos
    menores a 5~min, sin valores atípicos). Corrida: 2026-04-16.}
    \label{fig:impedancia-sistemas}
    \fuentefigura{Elaboración propia con \texttt{VFTModel}.}
\end{figure}

Los resultados son coherentes con la metodología adoptada. El caso más
sensible es el del \textbf{Trolebús}, cuya velocidad efectiva (\DIFdelbegin \DIFdel{~}\DIFdelend \DIFaddbegin \DIFadd{$\approx$\,}\DIFaddend 13.0~km/h con
$C_f = 1.38$) se sitúa ligeramente por debajo del límite inferior del rango
operativo documentado (14--22~km/h). Esta diferencia de aproximadamente
1~km/h (7\%) es atribuible a que el tipo de derecho de vía \textit{compartido}
no captura la preferencia operativa parcial que el trolebús mantiene en ciertas
vialidades de la \textsc{cdmx}. Para fines de análisis estructural de red,
esta desviación se encuentra dentro de los márgenes de calibración aceptables.

%% ── 5-Resultados/5-2-1-Fase1/5-2-1-2-Cobertura.tex
\subsection*{5 2 1 2 Cobertura}
% !TeX root = ../../tesis.tex
% =============================================================================
%  5.2.1.2 NIVEL DE COBERTURA ESPACIAL
% =============================================================================
\subsection{\bajoRevision{Nivel de Cobertura Espacial (\texorpdfstring{$C$}{C})}}
\label{subsec:cobertura-espacial}
\DIFdelbegin %DIFDELCMD < \notaRevision{Los resultados de esta subsección están sujetos a revisión con el
%DIFDELCMD < tutor. Los valores presentados corresponden a la corrida del modelo del
%DIFDELCMD < 18 de abril de 2026 con un isócrono de 800~m y pueden ajustarse tras
%DIFDELCMD < validación metodológica.}
%DIFDELCMD < %%%
\DIFdelend 

El \textbf{Nivel de Cobertura Espacial} ($C$) evalúa la equidad territorial de
la red de transporte masivo al cuantificar qué proporción del área de la
demarcación se encuentra dentro de un radio de caminata aceptable respecto
a alguna estación. Se define mediante la siguiente expresión:

\begin{equation}
    C = \frac{A_{\text{cubierta}}}{A_{\text{total}}} \times 100
    \label{eq:cobertura-res}
\end{equation}

donde $A_{\text{cubierta}}$ es el área resultante de la unión de los
isócronos peatonales (\termino{buffers} de 800~m) de todas las estaciones del sistema,
y $A_{\text{total}}$ es la superficie geodésica total de la demarcación
político-administrativa analizada.

El Cuadro~\ref{tab:peores-cdmx} identifica las cinco alcaldías de la
\textsc{cdmx} con menor cobertura del sistema de transporte masivo,
evidenciando las zonas de mayor exclusión territorial.

\begin{table}[H]
    \centering
    \caption[Alcaldías con menor cobertura de transporte masivo (CDMX)]{
    Cinco alcaldías de la Ciudad de México con menor porcentaje de cobertura
    territorial del sistema de transporte masivo. Isócrono: 800~m.
    Corrida: 2026-04-18.}
    \label{tab:peores-cdmx}
    \begin{tabular}{|l|r|r|r|}
        \hline
        \textbf{Demarcación} & \textbf{Área total (km²)} &
        \textbf{Área cubierta (km²)} & \textbf{Cobertura (\%)} \\
        \hline
        Milpa Alta              & 298.01 &  35.0    &  11.75 \\
        La Magdalena Contreras  &  63.37 &  16.33   &  25.76 \\
        Tlalpan                 & 314.25 &  94.31   &  30.01 \\
        Tláhuac                 &  85.78 &  43.85   &  51.12 \\
        Cuajimalpa de Morelos   &  71.10 &  37.12   &  52.20 \\
        \hline
    \end{tabular}
    \fuentefigura{Elaboración propia con \texttt{VFTModel} y datos de
    \texttt{Apimetro}.}
\end{table}

\begin{table}[H]
    \centering
    \caption[Municipios del EDOMEX con menor cobertura de transporte masivo]{
    Cinco municipios del Estado de México con menor porcentaje de cobertura
    territorial del sistema de transporte masivo. Isócrono: 800~m.
    Corrida: 2026-04-18.}
    \label{tab:peores-edomex}
    \begin{tabular}{|l|r|r|r|}
        \hline
        \textbf{Demarcación} & \textbf{Área total (km²)} &
        \textbf{Área cubierta (km²)} & \textbf{Cobertura (\%)} \\
        \hline
        Tonanitla       &   8.52 & 0.00 & 0.00 \\
        Capulhuac       &  21.50 & 0.00 & 0.00 \\
        Calimaya        & 104.26 & 0.00 & 0.00 \\
        Donato Guerra   & 181.36 & 0.00 & 0.00 \\
        Ayapango        &  50.83 & 0.00 & 0.00 \\
        \hline
    \end{tabular}
    \fuentefigura{Elaboración propia con \texttt{VFTModel} y datos de
    \texttt{Apimetro}.}
\end{table}

% =============================================================================
%  Desglose de Cobertura por Sistema de Transporte
% =============================================================================

\subsubsection{Sistema de Transporte Colectivo Metro (METRO)}
\label{subsubsec:cobertura-metro}

El Sistema de Transporte Colectivo Metro constituye la columna vertebral de la
movilidad masiva en la Ciudad de México. Debido a su naturaleza de red ferroviaria
subterránea y elevada con derecho de vía exclusivo, sus estaciones actúan como
nodos estructurales de alta capacidad que articulan los flujos intermodales de la
metrópoli. El análisis de su cobertura espacial mediante isócronos peatonales de
800~m permite identificar las demarcaciones que se benefician directamente de este
sistema, así como aquellas periferias y zonas centrales que presentan una
desconexión crítica respecto a la red de transporte más masiva de la entidad.

El Cuadro~\ref{tab:cobertura-metro-cdmx} detalla las cinco alcaldías de la Ciudad de
México con los niveles más altos de cobertura territorial respecto a la red del
Metro, evidenciando la concentración de la infraestructura ferroviaria en el núcleo
urbano denso de la cuenca.

\begin{table}[H]
    \centering
    \caption[Alcaldías con mayor cobertura del STC Metro (CDMX)]{
    Cinco alcaldías de la Ciudad de México con mayor porcentaje de cobertura
    territorial para el Sistema de Transporte Colectivo Metro. Isócrono: 800~m.
    Corrida: 2026-04-18.}
    \label{tab:cobertura-metro-cdmx}
    \begin{tabular}{|l|r|r|r|}
        \hline
        \textbf{Demarcación} & \textbf{Área total (km²)} & \textbf{Área cubierta (km²)} & \textbf{Cobertura (\%)} \\
        \hline
        Cuauhtémoc          & 32.50 &  28.35 &  87.23 \\
        Benito Juárez       & 26.68 &  21.34 &  79.99 \\
        Venustiano Carranza & 33.84 &  24.39 &  72.77 \\
        Iztacalco           & 23.08 &  11.66 &  50.52 \\
        Azcapotzalco        & 33.50 &  13.80 &  41.18 \\
        \hline
    \end{tabular}
    \fuentefigura{Elaboración propia con \texttt{VFTModel} y datos de \texttt{Apimetro}.}
\end{table}

Por su parte, la distribución geográfica de esta cobertura se visualiza en la
Figura~\ref{fig:mapa-cobertura-metro}, donde se aprecia la concentración de la
infraestructura en la zona central de la cuenca y el drástico desvanecimiento del
servicio hacia las demarcaciones del sur y del poniente.

\begin{figure}[H]
    \centering
    \figinclude[width=0.85\textwidth]{Figures/Cap5/Cobertura/Mapas/NB02_fig02_METRO_cobertura.png}
    \caption[Mapa de cobertura territorial del STC Metro en la CDMX]{Mapa del
    nivel de cobertura territorial del Sistema de Transporte Colectivo Metro
    por demarcación en la Ciudad de México, calculado a partir de \termino{buffers}
    de caminata peatonal de 800~m.}
    \label{fig:mapa-cobertura-metro}
    \fuentefigura{Elaboración propia con \texttt{VFTModel} utilizando visualización geoespacial.}
\end{figure}

\subsubsection{Metrobús (MB)}
\label{subsubsec:cobertura-mb}

El sistema Metrobús, que opera bajo el modelo de Autobuses de Tránsito Rápido
(\termino{Bus Rapid Transit}, BRT), se ha consolidado como la segunda red de
transporte masivo más importante de la Ciudad de México. Dada su mayor flexibilidad
constructiva frente a los sistemas ferroviarios pesados, su trazado ha logrado
expandirse por corredores viales estratégicos, complementando y alimentando a la
red del Metro. No obstante, al evaluar su cobertura mediante isócronos peatonales
de 800~m, persisten notables brechas territoriales, dejando a varias demarcaciones
periféricas con acceso limitado o nulo a esta red troncal.

El Cuadro~\ref{tab:cobertura-mb-cdmx} presenta las cinco alcaldías con mayor
proporción de su territorio cubierto por el Metrobús, destacando la concentración
del servicio en las zonas centro, norte y oriente de la ciudad, siguiendo los ejes
viales primarios.

\begin{table}[H]
    \centering
    \caption[Alcaldías con mayor cobertura de Metrobús (CDMX)]{
    Cinco alcaldías de la Ciudad de México con mayor porcentaje de cobertura
    territorial para el sistema Metrobús (MB). Isócrono: 800~m.
    Corrida: 2026-04-18.}
    \label{tab:cobertura-mb-cdmx}
    \begin{tabular}{|l|r|r|r|}
        \hline
        \textbf{Demarcación} & \textbf{Área total (km²)} & \textbf{Área cubierta (km²)} & \textbf{Cobertura (\%)} \\
        \hline
        Cuauhtémoc          & 32.50 &  26.74 &  82.27 \\
        Benito Juárez       & 26.68 &  16.21 &  60.74 \\
        Iztacalco           & 23.08 &  13.95 &  60.44 \\
        Venustiano Carranza & 33.84 &  16.10 &  47.58 \\
        Gustavo A. Madero   & 87.84 &  38.75 &  44.12 \\
        \hline
    \end{tabular}
    \fuentefigura{Elaboración propia con \texttt{VFTModel} y datos de \texttt{Apimetro}.}
\end{table}

Visualmente, el alcance espacial del sistema se ilustra en la
Figura~\ref{fig:mapa-cobertura-mb}. El mapa confirma cómo las líneas de BRT se
concentran fuertemente en las zonas centro, norte y oriente —siguiendo los ejes
viales primarios—, mientras que las barreras topográficas y de traza urbana
mantienen a las demarcaciones del surponiente con una cobertura residual o
inexistente.

\begin{figure}[H]
    \centering
    \figinclude[width=0.85\textwidth]{Figures/Cap5/Cobertura/Mapas/NB02_fig02_MB_cobertura.png}
    \caption[Mapa de cobertura territorial del Metrobús en la CDMX]{Mapa del
    nivel de cobertura territorial del sistema Metrobús por demarcación en la
    Ciudad de México, calculado a partir de \termino{buffers} de caminata
    peatonal de 800~m.}
    \label{fig:mapa-cobertura-mb}
    \fuentefigura{Elaboración propia con \texttt{VFTModel} utilizando visualización geoespacial.}
\end{figure}

\subsubsection{Cablebús (CBB)}
\label{subsubsec:cobertura-cbb}

El Cablebús es un sistema de teleférico urbano diseñado para resolver problemas de
accesibilidad en zonas de alta marginación y topografía compleja, actuando como
alimentador de alta capacidad hacia el Metro. Su cobertura se concentra en
Gustavo A. Madero (12.88\,\%), Miguel Hidalgo (11.85\,\%) e Iztapalapa
(11.27\,\%), con presencia marginal en el resto de la ciudad
(Cuadro~\ref{tab:cobertura-cbb-cdmx}).

\begin{table}[H]
    \centering
    \caption[Alcaldías con mayor cobertura de Cablebús (CDMX)]{
    Cinco alcaldías de la Ciudad de México con mayor porcentaje de cobertura
    territorial para el sistema Cablebús (CBB). Isócrono: 800~m.
    Corrida: 2026-04-18.}
    \label{tab:cobertura-cbb-cdmx}
    \begin{tabular}{|l|r|r|r|}
        \hline
        \textbf{Demarcación} & \textbf{Área total (km²)} & \textbf{Área cubierta (km²)} & \textbf{Cobertura (\%)} \\
        \hline
        Gustavo A. Madero   &  87.84 & 11.31 &  12.88 \\
        Miguel Hidalgo      &  46.39 &  5.50 &  11.85 \\
        Iztapalapa          & 113.07 & 12.74 &  11.27 \\
        Álvaro Obregón      &  95.82 &  4.38 &   4.57 \\
        Venustiano Carranza &  33.84 &  0.00 &   0.00 \\
        \hline
    \end{tabular}
    \fuentefigura{Elaboración propia con \texttt{VFTModel} y datos de \texttt{Apimetro}.}
\end{table}

La Figura~\ref{fig:mapa-cobertura-cbb} ilustra esta focalización espacial. El mapa
revela cómo las áreas de influencia del Cablebús se restringen a corredores muy
específicos (como en Gustavo A. Madero, Iztapalapa y Álvaro Obregón), contrastando
drásticamente con el vacío de infraestructura aérea en el resto de la cuenca.

\begin{figure}[H]
    \centering
    \figinclude[width=0.85\textwidth]{Figures/Cap5/Cobertura/Mapas/NB02_fig02_CBB_cobertura.png}
    \caption[Mapa de cobertura territorial del Cablebús en la CDMX]{Mapa del
    nivel de cobertura territorial del sistema Cablebús por demarcación en la
    Ciudad de México, calculado a partir de \termino{buffers} de caminata
    peatonal de 800~m.}
    \label{fig:mapa-cobertura-cbb}
    \fuentefigura{Elaboración propia con \texttt{VFTModel} utilizando visualización geoespacial.}
\end{figure}

\subsubsection{Red de Transporte de Pasajeros (RTP)}
\label{subsubsec:cobertura-rtp}

La Red de Transporte de Pasajeros (\textsc{rtp}) desempeña un rol fundamental en
la equidad espacial del sistema de movilidad de la Ciudad de México. A diferencia
de los sistemas confinados, la \textsc{rtp} opera como un servicio de superficie
con vocación social y de alimentación, diseñado específicamente para penetrar en
zonas periféricas, de topografía compleja y de menor densidad, donde la
implementación de infraestructura pesada es inviable. Su extensión territorial es
notablemente mayor que la de cualquier otro sistema de la metrópoli.

El Cuadro~\ref{tab:cobertura-rtp-cdmx} presenta las cinco alcaldías con mayor
proporción de cobertura por parte de la \textsc{rtp}, confirmando la concentración
del servicio en el centro y norte de la ciudad, donde la densidad de rutas y paradas
alcanza sus niveles más altos.

\begin{table}[H]
    \centering
    \caption[Alcaldías con mayor cobertura de RTP (CDMX)]{
    Cinco alcaldías de la Ciudad de México con mayor porcentaje de cobertura
    territorial para la Red de Transporte de Pasajeros (RTP). Isócrono: 800~m.
    Corrida: 2026-04-18.}
    \label{tab:cobertura-rtp-cdmx}
    \begin{tabular}{|l|r|r|r|}
        \hline
        \textbf{Demarcación} & \textbf{Área total (km²)} & \textbf{Área cubierta (km²)} & \textbf{Cobertura (\%)} \\
        \hline
        Miguel Hidalgo      &  46.39 &  41.90 &  90.32 \\
        Cuauhtémoc          &  32.50 &  29.32 &  90.22 \\
        Azcapotzalco        &  33.50 &  30.08 &  89.81 \\
        Gustavo A. Madero   &  87.84 &  75.86 &  86.37 \\
        Iztapalapa          & 113.07 &  93.66 &  82.83 \\
        \hline
    \end{tabular}
    \fuentefigura{Elaboración propia con \texttt{VFTModel} y datos de \texttt{Apimetro}.}
\end{table}

La Figura~\ref{fig:mapa-cobertura-rtp} muestra la dispersión territorial de este
sistema. A diferencia del Metro o del Metrobús, el mapa de la \textsc{rtp} revela
una capilaridad mucho mayor hacia el sur y poniente de la ciudad, confirmando su
rol estratégico como conector de las zonas más alejadas de la cuenca con la red
troncal.

\begin{figure}[H]
    \centering
    \figinclude[width=0.85\textwidth]{Figures/Cap5/Cobertura/Mapas/NB02_fig02_RTP_cobertura.png}
    \caption[Mapa de cobertura territorial de la RTP en la CDMX]{Mapa del
    nivel de cobertura territorial de la Red de Transporte de Pasajeros
    (\textsc{rtp}) por demarcación en la Ciudad de México, calculado a partir
    de \termino{buffers} de caminata peatonal de 800~m.}
    \label{fig:mapa-cobertura-rtp}
    \fuentefigura{Elaboración propia con \texttt{VFTModel} utilizando visualización geoespacial.}
\end{figure}

\subsubsection{Trolebús (TROLE)}
\label{subsubsec:cobertura-trole}

El Trolebús —operado por el Servicio de Transportes Eléctricos (\textsc{ste})—
es uno de los ejes históricos de electromovilidad de cero emisiones en la capital,
con presencia relevante en el centro y oriente de la ciudad. El
Cuadro~\ref{tab:cobertura-trole-cdmx} registra coberturas superiores al 50\,\%
en cuatro alcaldías centrales, encabezadas por Cuauhtémoc (69.90\,\%).

\textbf{Nota metodológica:} Es importante precisar que, al momento de la
ejecución de este análisis espacial (18 de abril de 2026), el Portal de Datos
Abiertos de la Ciudad de México aún no había liberado los trazados geoespaciales
ni el padrón de paradas oficiales correspondientes a la \textbf{Línea 13} del
Trolebús. Por consiguiente, dicha línea no se encuentra contabilizada en los
isócronos de cobertura presentados en esta sección.

\begin{table}[H]
    \centering
    \caption[Alcaldías con mayor cobertura de Trolebús (CDMX)]{
    Cinco alcaldías de la Ciudad de México con mayor porcentaje de cobertura
    territorial para el sistema Trolebús (TROLE). Isócrono: 800~m.
    Corrida: 2026-04-18.}
    \label{tab:cobertura-trole-cdmx}
    \begin{tabular}{|l|r|r|r|}
        \hline
        \textbf{Demarcación} & \textbf{Área total (km²)} & \textbf{Área cubierta (km²)} & \textbf{Cobertura (\%)} \\
        \hline
        Cuauhtémoc    &  32.50 & 22.72 &  69.90 \\
        Benito Juárez &  26.68 & 16.50 &  61.85 \\
        Azcapotzalco  &  33.50 & 17.48 &  52.18 \\
        Iztacalco     &  23.08 & 11.93 &  51.69 \\
        Coyoacán      & 125.17 & 25.17 &  46.71 \\
        \hline
    \end{tabular}
    \fuentefigura{Elaboración propia con \texttt{VFTModel} y datos de \texttt{Apimetro}.}
\end{table}

Por su parte, la Figura~\ref{fig:mapa-cobertura-trole} ilustra la huella espacial
del Trolebús. El mapa visibiliza la fuerte concentración del servicio en alcaldías
como Benito Juárez, Cuauhtémoc, Iztacalco e Iztapalapa (impulsada por el Trolebús
Elevado), contrastando drásticamente con la carencia de infraestructura eléctrica
hacia las periferias montañosas.

\begin{figure}[H]
    \centering
    \figinclude[width=0.85\textwidth]{Figures/Cap5/Cobertura/Mapas/NB02_fig02_TROLE_cobertura.png}
    \caption[Mapa de cobertura territorial del Trolebús en la CDMX]{Mapa del
    nivel de cobertura territorial del sistema Trolebús por demarcación en la
    Ciudad de México, calculado a partir de \termino{buffers} de caminata
    peatonal de 800~m.}
    \label{fig:mapa-cobertura-trole}
    \fuentefigura{Elaboración propia con \texttt{VFTModel} utilizando visualización geoespacial.}
\end{figure}

\subsubsection{Corredores Concesionados (CC)}
\label{subsubsec:cobertura-cc}

Los Corredores Concesionados constituyen la red de rutas de transporte de
superficie operadas por empresas privadas con concesión del Gobierno de la Ciudad
de México. Con base en su lógica de trazado sobre los ejes viales primarios y
secundarios de la ciudad consolidada, este sistema presenta uno de los niveles de
cobertura territorial más altos de la metrópoli, especialmente en las alcaldías de
mayor densidad urbana del centro y el oriente. Más que un sistema residual, los
Corredores Concesionados actúan como la red capilar que penetra en los intersticios
que los sistemas troncales no pueden servir directamente.

El Cuadro~\ref{tab:cobertura-cc-cdmx} presenta las cinco alcaldías con mayor
cobertura territorial de este sistema, donde se aprecia la saturación práctica del
servicio en el núcleo urbano denso de la capital.

\begin{table}[H]
    \centering
    \caption[Alcaldías con mayor cobertura de Corredores Concesionados (CDMX)]{
    Cinco alcaldías de la Ciudad de México con mayor porcentaje de cobertura
    territorial para el sistema de Corredores Concesionados (CC). Isócrono: 800~m.
    Corrida: 2026-04-18.}
    \label{tab:cobertura-cc-cdmx}
    \begin{tabular}{|l|r|r|r|}
        \hline
        \textbf{Demarcación} & \textbf{Área total (km²)} & \textbf{Área cubierta (km²)} & \textbf{Cobertura (\%)} \\
        \hline
        Cuauhtémoc          & 32.50 & 32.47 &  99.90 \\
        Benito Juárez       & 26.68 & 26.32 &  98.65 \\
        Iztacalco           & 23.08 & 22.25 &  96.42 \\
        Miguel Hidalgo      & 46.39 & 44.33 &  95.57 \\
        Venustiano Carranza & 33.84 & 27.77 &  82.06 \\
        \hline
    \end{tabular}
    \fuentefigura{Elaboración propia con \texttt{VFTModel} y datos de \texttt{Apimetro}.}
\end{table}

La Figura~\ref{fig:mapa-cobertura-cc} confirma la distribución omnipresente de
este sistema en la ciudad consolidada. El mapa ilustra cómo los Corredores
Concesionados recubren de manera prácticamente total las alcaldías centrales,
mientras que el servicio se va diluyendo hacia las periferias con suelo de
conservación o menor densidad.

\begin{figure}[H]
    \centering
    \figinclude[width=0.85\textwidth]{Figures/Cap5/Cobertura/Mapas/NB02_fig02_CC_cobertura.png}
    \caption[Mapa de cobertura territorial de Corredores Concesionados en la CDMX]{
    Mapa del nivel de cobertura territorial del sistema de Corredores
    Concesionados por demarcación en la Ciudad de México, calculado a partir de
    \termino{buffers} de caminata peatonal de 800~m.}
    \label{fig:mapa-cobertura-cc}
    \fuentefigura{Elaboración propia con \texttt{VFTModel} utilizando visualización geoespacial.}
\end{figure}

\subsubsection{Tren Ligero (TL)}
\label{subsubsec:cobertura-tl}

El Tren Ligero es un sistema de transporte férreo de capacidad media operado por
el Servicio de Transportes Eléctricos (\textsc{ste}). Su trazado consta de una
única línea que conecta la estación terminal Tasqueña de la Línea 2 del Metro con
el centro de la alcaldía Xochimilco. Debido a esta configuración lineal y su
confinamiento exclusivo en el sur de la cuenca, su área de influencia geográfica
es altamente focalizada.

El Cuadro~\ref{tab:cobertura-tl-cdmx} presenta las alcaldías con mayor nivel de
cobertura por parte del Tren Ligero. Únicamente Coyoacán, Xochimilco y Tlalpan
registran acceso al sistema dentro del isócrono peatonal de 800~m, mientras que
el resto de las demarcaciones carece por completo de este servicio.

\begin{table}[H]
    \centering
    \caption[Alcaldías con mayor cobertura de Tren Ligero (CDMX)]{
    Cinco alcaldías de la Ciudad de México con mayor porcentaje de cobertura
    territorial para el sistema Tren Ligero (TL). Isócrono: 800~m.
    Corrida: 2026-04-18.}
    \label{tab:cobertura-tl-cdmx}
    \begin{tabular}{|l|r|r|r|}
        \hline
        \textbf{Demarcación} & \textbf{Área total (km²)} & \textbf{Área cubierta (km²)} & \textbf{Cobertura (\%)} \\
        \hline
        Coyoacán            &  53.88 & 9.31 &  17.28 \\
        Xochimilco          & 114.03 & 6.82 &   5.98 \\
        Tlalpan             & 314.25 & 4.07 &   1.30 \\
        Venustiano Carranza &  33.84 & 0.00 &   0.00 \\
        Tláhuac             &  85.78 & 0.00 &   0.00 \\
        \hline
    \end{tabular}
    \fuentefigura{Elaboración propia con \texttt{VFTModel} y datos de \texttt{Apimetro}.}
\end{table}

La Figura~\ref{fig:mapa-cobertura-tl} ilustra espacialmente esta restricción
territorial. El mapa confirma cómo el servicio es exclusivo del corredor sur,
actuando principalmente como alimentador de la red masiva central para los
habitantes de las alcaldías sureñas.

\begin{figure}[H]
    \centering
    \figinclude[width=0.85\textwidth]{Figures/Cap5/Cobertura/Mapas/NB02_fig02_TL_cobertura.png}
    \caption[Mapa de cobertura territorial del Tren Ligero en la CDMX]{Mapa del
    nivel de cobertura territorial del sistema Tren Ligero por demarcación en
    la Ciudad de México, calculado a partir de \termino{buffers} de caminata
    peatonal de 800~m.}
    \label{fig:mapa-cobertura-tl}
    \fuentefigura{Elaboración propia con \texttt{VFTModel} utilizando visualización geoespacial.}
\end{figure}

\subsubsection{Mexibús (MXB)}
\label{subsubsec:cobertura-mxb}

El sistema Mexibús opera sobre corredores viales del Estado de México con una
presencia marginal en la Ciudad de México. Su cobertura máxima dentro de los
límites capitalinos no supera el 6.58\,\% (Iztacalco), siendo nula en las
demarcaciones del sur y oriente (Cuadro~\ref{tab:cobertura-mxb-cdmx}). Para
la visualización cartográfica de este sistema en el contexto de la red
metropolitana véase la Figura~\ref{fig:anx-mapa-cobertura-total} del
Anexo~\ref{anx:cartografia-presentacion}.

\begin{table}[H]
    \centering
    \caption[Alcaldías con mayor cobertura de Mexibús (CDMX)]{
    Cinco alcaldías de la Ciudad de México con mayor porcentaje de cobertura
    territorial para el sistema Mexibús (MXB). Isócrono: 800~m.
    Corrida: 2026-04-18.}
    \label{tab:cobertura-mxb-cdmx}
    \begin{tabular}{|l|r|r|r|}
        \hline
        \textbf{Demarcación} & \textbf{Área total (km²)} & \textbf{Área cubierta (km²)} & \textbf{Cobertura (\%)} \\
        \hline
        Iztacalco           &  23.08 & 1.52 & 6.58 \\
        Venustiano Carranza &  33.84 & 1.90 & 5.62 \\
        Gustavo A. Madero   &  87.84 & 3.23 & 3.67 \\
        Xochimilco          & 114.03 & 0.00 & 0.00 \\
        Tláhuac             &  85.78 & 0.00 & 0.00 \\
        \hline
    \end{tabular}
    \fuentefigura{Elaboración propia con \texttt{VFTModel} y datos de \texttt{Apimetro}.}
\end{table}

\subsubsection{Mexicable (MEXICABLE)}
\label{subsubsec:cobertura-mexicable}

El Mexicable es un teleférico urbano del Estado de México cuya incursión en la
Ciudad de México es mínima: únicamente Gustavo A. Madero registra cobertura
efectiva (3.19\,\%), mientras el resto del territorio capitalino presenta valores
nulos (Cuadro~\ref{tab:cobertura-mexicable-cdmx}). Para la visualización
cartográfica de este sistema en el contexto de la red metropolitana véase la
Figura~\ref{fig:anx-mapa-cobertura-total} del
Anexo~\ref{anx:cartografia-presentacion}.

\begin{table}[H]
    \centering
    \caption[Alcaldías con mayor cobertura de Mexicable (CDMX)]{
    Cinco alcaldías de la Ciudad de México con mayor porcentaje de cobertura
    territorial para el sistema Mexicable. Isócrono: 800~m.
    Corrida: 2026-04-18.}
    \label{tab:cobertura-mexicable-cdmx}
    \begin{tabular}{|l|r|r|r|}
        \hline
        \textbf{Demarcación} & \textbf{Área total (km²)} & \textbf{Área cubierta (km²)} & \textbf{Cobertura (\%)} \\
        \hline
        Gustavo A. Madero   &  87.84 & 2.80 & 3.19 \\
        Venustiano Carranza &  33.84 & 0.00 & 0.00 \\
        Xochimilco          & 114.03 & 0.00 & 0.00 \\
        Tláhuac             &  85.78 & 0.00 & 0.00 \\
        Tlalpan             & 314.25 & 0.00 & 0.00 \\
        \hline
    \end{tabular}
    \fuentefigura{Elaboración propia con \texttt{VFTModel} y datos de \texttt{Apimetro}.}
\end{table}

\subsubsection{Tren Suburbano (SUB)}
\label{subsubsec:cobertura-sub}

El Tren Suburbano conecta la terminal Buenavista con Cuautitlán (Estado de
México), limitando su área de influencia en la Ciudad de México al corredor norte.
Solo Cuauhtémoc (6.18\,\%) y Azcapotzalco (5.99\,\%) registran cobertura
efectiva, con el resto del territorio capitalino sin acceso directo al sistema
(Cuadro~\ref{tab:cobertura-sub-cdmx}).

\begin{table}[H]
    \centering
    \caption[Alcaldías con mayor cobertura de Tren Suburbano (CDMX)]{
    Cinco alcaldías de la Ciudad de México con mayor porcentaje de cobertura
    territorial para el sistema Tren Suburbano (SUB). Isócrono: 800~m.
    Corrida: 2026-04-18.}
    \label{tab:cobertura-sub-cdmx}
    \begin{tabular}{|l|r|r|r|}
        \hline
        \textbf{Demarcación} & \textbf{Área total (km²)} & \textbf{Área cubierta (km²)} & \textbf{Cobertura (\%)} \\
        \hline
        Cuauhtémoc          &  32.50 & 2.01 & 6.18 \\
        Azcapotzalco        &  33.50 & 2.01 & 5.99 \\
        Xochimilco          & 114.03 & 0.00 & 0.00 \\
        Venustiano Carranza &  33.84 & 0.00 & 0.00 \\
        Tláhuac             &  85.78 & 0.00 & 0.00 \\
        \hline
    \end{tabular}
    \fuentefigura{Elaboración propia con \texttt{VFTModel} y datos de \texttt{Apimetro}.}
\end{table}

La Figura~\ref{fig:mapa-cobertura-sub} confirma la huella territorial de este
sistema. El mapa evidencia cómo la cobertura del Tren Suburbano apenas alcanza la
alcaldía Cuauhtémoc y el extremo norte de Azcapotzalco, mientras que el resto del
territorio capitalino permanece fuera de su radio de influencia peatonal.

\begin{figure}[H]
    \centering
    \figinclude[width=0.85\textwidth]{Figures/Cap5/Cobertura/Mapas/NB02_fig02_SUB_cobertura.png}
    \caption[Mapa de cobertura territorial del Tren Suburbano en la CDMX]{Mapa
    del nivel de cobertura territorial del sistema Tren Suburbano por demarcación
    en la Ciudad de México, calculado a partir de \termino{buffers} de caminata
    peatonal de 800~m.}
    \label{fig:mapa-cobertura-sub}
    \fuentefigura{Elaboración propia con \texttt{VFTModel} utilizando visualización geoespacial.}
\end{figure}

\subsubsection{Tren Interurbano (INTERURBANO)}
\label{subsubsec:cobertura-interurbano}

El Tren Interurbano ``El Insurgente'', de vocación regional Toluca-\cdmx{},
concentra su presencia en el sector poniente con cobertura efectiva únicamente en
Cuajimalpa de Morelos (2.23\,\%) y Álvaro Obregón (0.44\,\%), mientras el resto
de la ciudad carece de acceso directo al sistema
(Cuadro~\ref{tab:cobertura-interurbano-cdmx}). Para la visualización cartográfica
de este sistema en el contexto de la red metropolitana véase la
Figura~\ref{fig:anx-mapa-cobertura-total} del
Anexo~\ref{anx:cartografia-presentacion}.

\begin{table}[H]
    \centering
    \caption[Alcaldías con mayor cobertura de Tren Interurbano (CDMX)]{
    Cinco alcaldías de la Ciudad de México con mayor porcentaje de cobertura
    territorial para el sistema Tren Interurbano (INTERURBANO). Isócrono: 800~m.
    Corrida: 2026-04-18.}
    \label{tab:cobertura-interurbano-cdmx}
    \begin{tabular}{|l|r|r|r|}
        \hline
        \textbf{Demarcación} & \textbf{Área total (km²)} & \textbf{Área cubierta (km²)} & \textbf{Cobertura (\%)} \\
        \hline
        Cuajimalpa de Morelos &  71.10 & 1.59 & 2.23 \\
        Álvaro Obregón        &  95.82 & 0.42 & 0.44 \\
        Venustiano Carranza   &  33.84 & 0.00 & 0.00 \\
        Xochimilco            & 114.03 & 0.00 & 0.00 \\
        Tláhuac               &  85.78 & 0.00 & 0.00 \\
        \hline
    \end{tabular}
    \fuentefigura{Elaboración propia con \texttt{VFTModel} y datos de \texttt{Apimetro}.}
\end{table}

\subsubsection{Pumabús (PB)}
\label{subsubsec:cobertura-pb}

El sistema Pumabús es un servicio de transporte interno operado por la Universidad
Nacional Autónoma de México (UNAM), de carácter público y gratuito, cuya operación
se circunscribe al campus de Ciudad Universitaria. Si bien el modelo registra una
cobertura del 20.13\,\% en Coyoacán ---el porcentaje más alto entre los sistemas
de escala campus---, este valor refleja la concentración geométrica del isócrono
dentro del polígono universitario y no un alcance funcional sobre la trama urbana
circundante. El resto de las demarcaciones capitalinas registra coberturas menores
al 2.5\,\% o nulas, confirmando su carácter de sistema cautivo del recinto
universitario (Cuadro~\ref{tab:cobertura-pb-cdmx}). Para la visualización
cartográfica de este sistema en el contexto de la red metropolitana véase la
Figura~\ref{fig:anx-mapa-cobertura-total} del
Anexo~\ref{anx:cartografia-presentacion}.

\begin{table}[H]
    \centering
    \caption[Alcaldías con mayor cobertura de Pumabús (CDMX)]{
    Cinco alcaldías de la Ciudad de México con mayor porcentaje de cobertura
    territorial para el sistema Pumabús (PB). Isócrono: 800~m.
    Corrida: 2026-04-18.}
    \label{tab:cobertura-pb-cdmx}
    \begin{tabular}{|l|r|r|r|}
        \hline
        \textbf{Demarcación} & \textbf{Área total (km²)} & \textbf{Área cubierta (km²)} & \textbf{Cobertura (\%)} \\
        \hline
        Coyoacán            &  53.88 & 10.85 &  20.13 \\
        Álvaro Obregón      &  95.82 &  2.25 &   2.34 \\
        Tlalpan             & 314.25 &  0.25 &   0.08 \\
        Venustiano Carranza &  33.84 &  0.00 &   0.00 \\
        Xochimilco          & 114.03 &  0.00 &   0.00 \\
        \hline
    \end{tabular}
    \fuentefigura{Elaboración propia con \texttt{VFTModel} y datos de \texttt{Apimetro}.}
\end{table}


% =============================================================================
%  Vista cartográfica de la cobertura total de la red
% =============================================================================
\subsubsection{Vista Cartográfica de la Cobertura Total}
\label{subsubsec:cobertura-total-cartografia}

El análisis pormenorizado por sistema presentado en las secciones anteriores
permite construir una lectura integral del territorio metropolitano en términos
de accesibilidad al transporte masivo. Las
Figuras~\ref{fig:mapa-cobertura-red-total} y~\ref{fig:mapa-cobertura-calor}
sintetizan esta lectura mediante dos representaciones cartográficas complementarias
producidas con el sistema \textbf{Transport-gis-zmvm-mjg}: la primera muestra el
alcance territorial acumulado de los isócronos de todos los sistemas; la segunda
expresa la densidad de esa cobertura como mapa de calor, revelando los núcleos de
máxima accesibilidad y los vacíos periféricos. Las versiones a página completa se
encuentran en el Anexo~\ref{anx:cartografia-presentacion}
(Figuras~\ref{fig:anx-mapa-cobertura-total}
y~\ref{fig:anx-mapa-cobertura-calor}).

\begin{figure}[H]
    \centering
    \figinclude[width=0.90\textwidth]{Figures/Mapas/Mapa_Cobertura_Red_Total}
    \caption[Cobertura espacial total de la red multimodal — isócronas 800~m]{
    Cobertura espacial integrada de los doce sistemas de transporte masivo de la
    \zmvm{}, calculada a partir de isócronos peatonales de 800~m. Las zonas sin
    cobertura se concentran en el sur y oriente de la cuenca metropolitana,
    evidenciando las demarcaciones con mayor exclusión territorial del sistema.}
    \label{fig:mapa-cobertura-red-total}
    \fuentefigura{Elaboración propia con \texttt{Transport-gis-zmvm-mjg};
    datos: GTFS SEMOVI 2024 e INEGI 2023.}
\end{figure}

\begin{figure}[H]
    \centering
    \figinclude[width=0.90\textwidth]{Figures/Mapas/Mapa_Cobertura_Mapa_Calor}
    \caption[Densidad de cobertura de la red multimodal — mapa de calor]{
    Densidad de la cobertura espacial de la red multimodal expresada como mapa
    de calor. Los núcleos de máxima accesibilidad se ubican en las alcaldías
    centrales (Cuauhtémoc, Benito Juárez, Iztacalco), mientras que las periferias
    sur y oriente presentan densidades cercanas a cero, confirmando la inequidad
    territorial diagnosticada.}
    \label{fig:mapa-cobertura-calor}
    \fuentefigura{Elaboración propia con \texttt{Transport-gis-zmvm-mjg};
    datos: GTFS SEMOVI 2024 e INEGI 2023.}
\end{figure}

%% ── 5-Resultados/5-2-1-Fase1/5-2-1-3-FuerzaCapilar.tex
\subsection*{5 2 1 3 FuerzaCapilar}
% !TeX root = ../../tesis.tex
% =============================================================================
%  5.2.1.3 FUERZA CAPILAR Y NODOS DE TRANSFERENCIA
% =============================================================================
\subsection{\bajoRevision{Fuerza Capilar y Nodos de Transferencia (\texorpdfstring{$FC$}{FC})}}
\label{subsec:fuerza-capilar}
\DIFdelbegin %DIFDELCMD < \notaRevision{Los resultados de esta subsección están sujetos a revisión con el
%DIFDELCMD < comité tutor. Los valores presentados corresponden a la corrida del modelo del
%DIFDELCMD < 18 de abril de 2026.}
%DIFDELCMD < %%%
\DIFdelend 

La \textbf{Fuerza Capilar} ($FC$) cuantifica el nivel de convergencia de flujos
de transporte en cada nodo o agrupación de nodos del grafo. El indicador se
evalúa desde dos perspectivas complementarias:

\textbf{A. Fuerza Capilar Geográfica ($FC_H$) — Macro-Hubs:} agrupa
topológicamente estaciones físicamente colindantes (radio de 100~m) en un
subgrafo $H$ que consolida los flujos de sus estaciones constituyentes,
representando el fenómeno urbano del transbordo físico (CETRAMs):

\begin{equation}
    FC_{H} = \sum_{j \in H} \bigl(k_{\text{in}}(j) + k_{\text{out}}(j)\bigr)
    - \text{Aristas Internas}
    \label{eq:fc-geografica}
\end{equation}

El Cuadro~\ref{tab:fuerza-capilar-hubs} detalla el comportamiento numérico de los
cinco principales Macro-Hubs de transferencia identificados en la red con base en
esta formulación, mientras que la Figura~\ref{fig:barplot-macro-hubs} permite
contrastar las magnitudes relativas de la concentración de flujos en los diez
nodos más significativos de la metrópoli.

\begin{table}[H]
    \centering
    \caption[Top 5 Macro-Hubs con mayor Fuerza Capilar Geográfica ($FC_H$)]{
    Los cinco Macro-Hubs de transferencia (nodos agrupados topológicamente a
    menos de 100~m) con mayor Fuerza Capilar Geográfica en la red metropolitana.
    Corrida: 2026-04-18.}
    \label{tab:fuerza-capilar-hubs}
    \fontsize{8.5pt}{10.5pt}\selectfont
    \begin{tabularx}{\textwidth}{|c|X|c|c|c|c|}
        \hline
        \textbf{ID Hub} & \textbf{Principales Estaciones Agrupadas} & \textbf{Sistemas} & \textbf{Nodos} & \textbf{Conexiones} \textit{(In / Out)} & \textbf{$FC_H$} \\
        \hline
        2174 & Zona Hospitales Sur (San Fernando, Chimalcoyotl, Coapa, Hosp. Manuel Gea Gtz., INER, Urgencias Respiratorias) & CC, RTP & 21 & 66 / 64 & \textbf{130} \\
        \hline
        136 & Eje Huipulco -- Estadio Azteca (Acoxpa, Huipulco, Estadio Azteca, Renato Leduc, San Alejandro) & CC, RTP, TL & 61 & 61 / 50 & \textbf{111} \\
        \hline
        1597 & Corredor Tlalpan Sur (Calle La Rosa, Tezoquipa, Galeana, Jojutla) & CC, RTP & 10 & 49 / 51 & \textbf{100} \\
        \hline
        259 & Nodo Barranca del Muerto (Barranca del Muerto, Alpes, Periférico Sur, José María Velasco, Miguel Ocaranza) & CC, RTP & 21 & 50 / 42 & \textbf{92} \\
        \hline
        129 & CETRAM Chapultepec (Metro Chapultepec, Acapulco, Agustín Melgar, Sonora, Veracruz, Tampico) & CC, METRO, RTP, TROLE & 29 & 37 / 40 & \textbf{77} \\
        \hline
    \end{tabularx}
    \vspace{1ex}
    \fuentefigura{Elaboración propia con \texttt{VFTModel} y datos de \texttt{Apimetro}.}
\end{table}

\begin{figure}[H]
    \centering
    \figinclude[width=0.85\textwidth]{Figures/Cap5/FuerzaCapilar/Tabla/NB03_fig01_top10_macrohubs_barplot.png}
    \caption[Distribución de los principales Macro-Hubs por Fuerza Capilar]{
    Gráfico de barras comparativo de los diez principales Macro-Hubs de la red
    metropolitana, ordenados con base en su nivel acumulado de Fuerza Capilar
    Geográfica ($FC_H$).}
    \label{fig:barplot-macro-hubs}
    \fuentefigura{Elaboración propia con \texttt{VFTModel} a partir de los resultados analíticos del modelo.}
\end{figure}

Por último, la distribución espacial y la cobertura geográfica de estos nodos de
transferencia masiva dentro del entorno urbano de la \zmvm{} se \DIFdelbegin \DIFdel{presenta mapeada
}\DIFdelend \DIFaddbegin \DIFadd{presentan mapeadas
}\DIFaddend de forma detallada en la Figura~\ref{fig:mapa-fuerza-capilar}.

\begin{figure}[H]
    \centering
    \figinclude[width=0.85\textwidth]{Figures/Cap5/FuerzaCapilar/Mapa/NB03_mapa01_macrohubs.png}
    \caption[Mapa de Macro-Hubs y Fuerza Capilar Geográfica]{
    Representación geoespacial de la \termino{Fuerza Capilar Geográfica}
    ($FC_H$). El mapa destaca la concentración de flujos de transporte en los
    principales nodos y zonas de transferencia física (Macro-Hubs) de la \zmvm.}
    \label{fig:mapa-fuerza-capilar}
    \fuentefigura{Elaboración propia con \texttt{VFTModel} utilizando visualización geoespacial.}
\end{figure}

% =============================================================================
%  5.2.1.4 FUERZA CAPILAR SISTÉMICA (MATEMÁTICA)
% =============================================================================
\subsection{\bajoRevision{Fuerza Capilar Sistémica (Matemática)}}
\label{subsec:fuerza-capilar-matematica}

\textbf{B. Fuerza Capilar Sistémica o Matemática ($FC$):} A diferencia de la
aproximación geográfica orientada a macro-hubs, esta perspectiva evalúa de manera
estricta la convergencia de flujos directamente sobre cada nodo topológico
individual. El análisis permite identificar los puntos exactos de la
infraestructura matemática de la red donde se concentra la mayor densidad de
aristas de entrada y salida, aislando singularidades críticas en la traza sin
aplicar tolerancias de proximidad espacial ni fusión de estaciones colindantes.

El Cuadro~\ref{tab:fuerza-capilar-matematica} detalla el comportamiento de los
cinco nodos individuales que registraron la mayor Fuerza Capilar pura dentro del
grafo de transporte metropolitano.

\begin{table}[H]
    \centering
    \caption[Top 5 Nodos con mayor Fuerza Capilar Sistémica (Matemática)]{
    Los cinco nodos individuales con mayor Fuerza Capilar Sistémica (Matemática)
    en la red metropolitana, evaluados de manera estricta sin agrupación
    geográfica. Corrida: 2026-04-18.}
    \label{tab:fuerza-capilar-matematica}
    \fontsize{9.5pt}{11.5pt}\selectfont
    \begin{tabular}{|c|l|c|c|c|}
        \hline
        \textbf{ID Nodo} & \textbf{Estación (Nodo Topológico)} & \textbf{Sistemas} & \textbf{Conexiones} \textit{(In / Out)} & \textbf{$FC$} \\
        \hline
        7270  & Luis Murillo                          & RTP & 21 / 20 & \textbf{41} \\
        \hline
        2817  & Calz. de Tlalpan -- Glorieta Huipulco & RTP & 19 / 15 & \textbf{34} \\
        \hline
        10211 & Sobre Calz. de Tlalpan y Luis Murillo & CC  & 16 / 17 & \textbf{33} \\
        \hline
        10717 & Tlalpan y Luis Murillo                & CC  & 18 / 15 & \textbf{33} \\
        \hline
        5540  & Estadio Azteca 2                      & RTP & 14 / 18 & \textbf{32} \\
        \hline
    \end{tabular}
    \vspace{1ex}
    \fuentefigura{Elaboración propia con \texttt{VFTModel} y datos de \texttt{Apimetro}.}
\end{table}

Las magnitudes y diferencias relativas entre los diez elementos individuales con
mayor concentración de conexiones se visualizan en el gráfico de barras de la
Figura~\ref{fig:barplot-nodos-matematicos}.

\begin{figure}[H]
    \centering
    \figinclude[width=0.85\textwidth]{Figures/Cap5/FuerzaCapilar/Tabla/NB03_fig02_top10_nodos_matematicos_barplot.png}
    \caption[Distribución de los principales nodos individuales por Fuerza Capilar]{
    Gráfico de barras comparativo de los diez nodos individuales con mayor
    Fuerza Capilar Sistémica (Matemática) en la red metropolitana.}
    \label{fig:barplot-nodos-matematicos}
    \fuentefigura{Elaboración propia con \texttt{VFTModel} a partir de los resultados analíticos del modelo.}
\end{figure}

La distribución geográfica e implicación espacial de estos puntos específicos de
alta carga conectiva dentro de la red urbana consolidada se presenta en la
Figura~\ref{fig:mapa-nodos-matematicos}.

\begin{figure}[H]
    \centering
    \figinclude[width=0.85\textwidth]{Figures/Cap5/FuerzaCapilar/Mapa/NB03_mapa02_nodos_matematicos.png}
    \caption[Mapa de Nodos Individuales y Fuerza Capilar Sistémica]{
    Representación geoespacial de la Fuerza Capilar Sistémica (Matemática) a
    nivel de nodo individual en la \zmvm{}.}
    \label{fig:mapa-nodos-matematicos}
    \fuentefigura{Elaboración propia con \texttt{VFTModel} utilizando visualización geoespacial.}
\end{figure}

\subsubsection{Análisis del Eje Periférico — Infraestructura Preexistente}
\label{subsubsec:periferico-fc}

Como diagnóstico previo a la propuesta del corredor anillar, se aislaron los nodos
y Macro-Hubs cuyos nombres contienen la denominación ``Periférico'', con el objetivo
de identificar los puntos de mayor tensión vehicular e intermodal sobre esta
vialidad fundamental. Este análisis se divide en dos enfoques: la concentración
geográfica (transbordos físicos en Macro-Hubs) y la concentración matemática
estricta (nodos topológicos individuales).

\paragraph{A. Fuerza Capilar Geográfica (Macro-Hubs) en el Periférico}

Al aplicar un radio de tolerancia de 100~m para agrupar las estaciones colindantes
sobre el eje, emergen los principales complejos de transbordo y alimentación. El
Cuadro~\ref{tab:periferico-hubs-geo} detalla los cinco Macro-Hubs de mayor Fuerza
Capilar Geográfica ($FC_H$), destacando la importancia de Barranca del Muerto y la
intersección con la zona de hospitales/Tlalpan como los anclajes más fuertes del
corredor.

\begin{table}[H]
    \centering
    \caption[Top 5 Macro-Hubs con mayor Fuerza Capilar sobre el Eje Periférico]{
    Cinco Macro-Hubs de transferencia con mayor Fuerza Capilar Geográfica
    ($FC_H$) ubicados sobre el Anillo Periférico. Corrida: 2026-04-18.}
    \label{tab:periferico-hubs-geo}
    \fontsize{8.5pt}{10.5pt}\selectfont
    \begin{tabularx}{\textwidth}{|c|X|c|c|c|c|}
        \hline
        \textbf{ID Hub} & \textbf{Principales Estaciones Agrupadas} & \textbf{Sistemas} & \textbf{Nodos} & \textbf{Conexiones} \textit{(In / Out)} & \textbf{$FC_H$} \\
        \hline
        259 & Nodo Barranca del Muerto (Barranca del Muerto, Alpes, Periférico Sur, Miguel Ocaranza) & CC, RTP & 21 & 50 / 42 & \textbf{92} \\
        \hline
        374 & Intersección Tlalpan--Periférico (Renato Leduc, Huipulco, Periférico, Av. Pedregal) & CC, RTP & 18 & 34 / 33 & \textbf{67} \\
        \hline
        453 & Nodo San Antonio (San Antonio, Periférico Sur, Puente Periférico Norte) & CC, RTP & 7 & 23 / 23 & \textbf{46} \\
        \hline
        402 & Eje Oriente Tepalcates (Constitución de Apatzingán, Canal de San Juan, Tepalcates) & CC, MB, RTP & 8 & 20 / 26 & \textbf{46} \\
        \hline
        291 & Intersección Xochimilco--Tepepan (16 de Septiembre, Kioto, Tec de Monterrey, Periférico Participación) & CC, RTP, TL & 14 & 23 / 22 & \textbf{45} \\
        \hline
    \end{tabularx}
    \vspace{1ex}
    \fuentefigura{Elaboración propia con \texttt{VFTModel} y datos de \texttt{Apimetro}.}
\end{table}

La Figura~\ref{fig:barplot-hubs-periferico} compara las magnitudes relativas de
estos agrupamientos geográficos, mientras que la Figura~\ref{fig:mapa-hubs-periferico}
espacializa su huella sobre la vialidad, revelando los polos naturales de atracción
de demanda que la futura infraestructura anillar deberá integrar.

\begin{figure}[H]
    \centering
    \figinclude[width=0.85\textwidth]{Figures/Cap5/FuerzaCapilar/Tabla/NB03_fig04_top10_macrohubs_periferico_barplot.png}
    \caption[Fuerza Capilar de Macro-Hubs en Anillo Periférico]{Gráfico de barras
    de los diez principales Macro-Hubs geográficos ubicados en el Anillo Periférico,
    ordenados por su nivel de Fuerza Capilar.}
    \label{fig:barplot-hubs-periferico}
    \fuentefigura{Elaboración propia con \texttt{VFTModel}.}
\end{figure}

\begin{figure}[H]
    \centering
    \figinclude[width=0.85\textwidth]{Figures/Cap5/FuerzaCapilar/Mapa/NB03_mapa03_macrohubs_periferico.png}
    \caption[Distribución espacial de Macro-Hubs en el Periférico]{Mapa de la
    Fuerza Capilar Geográfica ($FC_H$) en los Macro-Hubs del Anillo Periférico.}
    \label{fig:mapa-hubs-periferico}
    \fuentefigura{Elaboración propia con \texttt{VFTModel} utilizando visualización geoespacial.}
\end{figure}

\paragraph{B. Fuerza Capilar Sistémica (Nodos Matemáticos) en el Periférico}

Al desagregar los Hubs y analizar la Fuerza Capilar Sistémica ($FC$) estricta a
nivel de nodo individual, es posible identificar los paraderos específicos que
soportan la mayor carga de conexiones en la infraestructura actual. El
Cuadro~\ref{tab:periferico-nodos-math} presenta los cinco paraderos matemáticos
de mayor convergencia de flujo.

\begin{table}[H]
    \centering
    \caption[Nodos con mayor Fuerza Capilar Sistémica sobre el Eje Periférico]{
    Cinco nodos topológicos individuales con mayor Fuerza Capilar Matemática
    ($FC$) ubicados sobre el Anillo Periférico. Corrida: 2026-04-18.}
    \label{tab:periferico-nodos-math}
    \fontsize{9.5pt}{11.5pt}\selectfont
    \begin{tabular}{|c|l|c|r|r|c|}
        \hline
        \textbf{ID Nodo} & \textbf{Estación (Nodo Topológico)} & \textbf{Sistema} & \textbf{Entr.} & \textbf{Sal.} & \textbf{$FC$} \\
        \hline
        707  & Anillo Periférico y Blvd. Adolfo Ruiz Cortines & RTP & 14 & 16 & \textbf{30} \\
        \hline
        706  & Anillo Periférico y Blvd. Adolfo Ruiz Cortines & RTP & 13 & 13 & \textbf{26} \\
        \hline
        8538 & Periférico -- Barranca del Muerto              &  CC & 12 &  9 & \textbf{21} \\
        \hline
        602  & Anillo Periférico Canal de Garay -- Bilbao     & RTP & 10 & 10 & \textbf{20} \\
        \hline
        8510 & Periférico                                     & RTP &  9 & 11 & \textbf{20} \\
        \hline
    \end{tabular}
    \vspace{1ex}
    \fuentefigura{Elaboración propia con \texttt{VFTModel} y datos de \texttt{Apimetro}.}
\end{table}

La jerarquía analítica de estos puntos singulares se constata en la
Figura~\ref{fig:barplot-nodos-periferico}, seguida de la
Figura~\ref{fig:mapa-nodos-periferico}, que detalla la distribución geográfica a
escala de paradero individual, útil para el micro-posicionamiento de futuras
estaciones anillares.

\begin{figure}[H]
    \centering
    \figinclude[width=0.85\textwidth]{Figures/Cap5/FuerzaCapilar/Tabla/NB03_fig05_top10_nodos_matematicos_periferico_barplot.png}
    \caption[Fuerza Capilar de nodos individuales en Anillo Periférico]{
    Gráfico de barras comparativo de los diez principales nodos matemáticos
    individuales ubicados en el Anillo Periférico.}
    \label{fig:barplot-nodos-periferico}
    \fuentefigura{Elaboración propia con \texttt{VFTModel}.}
\end{figure}

\begin{figure}[H]
    \centering
    \figinclude[width=0.85\textwidth]{Figures/Cap5/FuerzaCapilar/Mapa/NB03_mapa04_nodos_matematicos_periferico.png}
    \caption[Distribución espacial de nodos matemáticos en el Periférico]{Mapa
    de la Fuerza Capilar Sistémica (Matemática) a nivel de nodo individual sobre
    el corredor del Anillo Periférico.}
    \label{fig:mapa-nodos-periferico}
    \fuentefigura{Elaboración propia con \texttt{VFTModel} utilizando visualización geoespacial.}
\end{figure}

\vspace{2ex}
\begin{quote}
    \small\textit{Nota aclaratoria:} los datos presentados en esta sección
    reflejan exclusivamente la infraestructura de transporte \textbf{preexistente}
    sobre el Anillo Periférico. No se ha trazado ni modelado aún la geometría
    del corredor anillar propuesto; la sección permanece en plano especulativo
    para identificar los puntos naturales de anclaje de la futura línea a nivel
    capilar.
\end{quote}

La Figura~\ref{fig:mapa-fc-presentacion} consolida, mediante cartografía de
presentación, la distribución geoespacial de los principales nodos \termino{hub}
de la red y su intensidad capilar. La versión a página completa se encuentra en
la Figura~\ref{fig:anx-mapa-fuerza-capilar} del
Anexo~\ref{anx:cartografia-presentacion}.

\begin{figure}[H]
    \centering
    \figinclude[width=0.90\textwidth]{Figures/Mapas/Mapa_Fuerza_Capilar}
    \caption[Cartografía de presentación — Fuerza Capilar y nodos \textit{hub} principales]{
    Representación cartográfica de los veinte nodos con mayor Fuerza Capilar en
    la red de transporte de la \zmvm{}. Los nodos señalados corresponden a los
    puntos de mayor influencia redistributiva de flujos intermodales, con especial
    concentración en el corredor sur de la cuenca metropolitana, sobre el eje
    que recorre el Anillo Periférico.}
    \label{fig:mapa-fc-presentacion}
    \fuentefigura{Elaboración propia con \texttt{Transport-gis-zmvm-mjg};
    datos: GTFS SEMOVI 2024 e INEGI 2023.}
\end{figure}

%% ── 5-Resultados/5-2-1-Fase1/5-2-1-4-FactorDesviacion.tex
\subsection*{5 2 1 4 FactorDesviacion}
% !TeX root = ../../tesis.tex
% =============================================================================
%  5.2.1.5 FACTOR DE DESVIACIÓN
% =============================================================================
\subsection{\bajoRevision{Factor de Desviación (\texorpdfstring{$DF$}{DF})}}
\label{subsec:detour-factor}
\notaRevision{Los resultados de esta subsección están sujetos a revisión con el
comité tutor. Los valores presentados corresponden a una muestra de 100~rutas
aleatorias ejecutadas el 27 de abril de 2026.}

El \textbf{Factor de Desviación} ($DF$) mide la eficiencia topológica de la
red comparando la distancia real recorrida a través de la infraestructura
con\DIFdelbegin \DIFdel{tra}\DIFdelend  la distancia teórica en línea recta (Haversine) entre el mismo par
origen-destino:

\begin{equation}
    DF = \frac{D_{\text{red}} + D_{\text{caminata}}}{D_{\text{haversine}}(O, D)}
    \label{eq:detour-factor-res}
\end{equation}

\begin{itemize}
    \item $D_{\text{red}}$: distancia acumulada sobre el trazo real de la ruta
    óptima en tiempo.
    \item $D_{\text{caminata}}$: suma de la primera y última milla peatonal.
    \item $D_{\text{haversine}}(O,D)$: distancia en línea recta entre origen y
    destino calculada con la fórmula de Haversine.
\end{itemize}

Un valor de $DF = 1.0$ indica una ruta perfectamente directa; valores superiores
a $1.5$ revelan ineficiencias topológicas significativas.

% -----------------------------------------------------------------------------
\subsubsection{Estadísticas Globales de la Red}
\label{subsubsec:df-estadisticas}

Sobre una muestra aleatoria de \textbf{100 pares origen-destino} (semilla
reproducible $= 42$), el análisis arrojó los siguientes estadísticos globales:

\begin{table}[H]
    \centering
    \caption[Estadísticos globales del Factor de Desviación — muestra de 100 rutas]{
    Estadísticos descriptivos del $DF$ sobre 100 pares O-D de la red de la \zmvm.
    Corrida: 2026-04-27.}
    \label{tab:df-estadisticas}
    \begin{tabular}{|l|r|}
        \hline
        \textbf{Estadístico} & \textbf{Valor} \\
        \hline
        $n$ (rutas calculadas)     & 100   \\
        Media ($\bar{DF}$)         & 1.545 \\
        Desviación estándar        & 0.517 \\
        Mínimo                     & 0.97  \\
        Percentil 25               & 1.27  \\
        Mediana                    & 1.40  \\
        Percentil 75               & 1.67  \\
        Máximo                     & 4.34  \\
        \hline
    \end{tabular}
    \fuentefigura{Elaboración propia con \texttt{VFTModel} (corrida 2026-04-27).}
\end{table}

La Figura~\ref{fig:df-caso-inicial} muestra la distribución geoespacial de los
100 pares O-D y las rutas calculadas sobre el territorio de la \zmvm{}.

\begin{figure}[H]
    \centering
    \figinclude[width=0.85\textwidth]{Figures/Cap5/FactorDesviacion/Mapa/NB04_mapa01_caso_inicial.png}
    \caption[Distribución de la muestra O-D para el Factor de Desviación]{
    Distribución de los 100 pares origen-destino de la muestra aleatoria y sus
    rutas óptimas en la \zmvm. Corrida: 2026-04-27.}
    \label{fig:df-caso-inicial}
    \fuentefigura{Elaboración propia con \texttt{VFTModel} (corrida 2026-04-27).}
\end{figure}

% -----------------------------------------------------------------------------
\subsubsection{Casos de Estudio Seleccionados}
\label{subsubsec:df-casos}

Se analizaron doce rutas que evalúan la conectividad desde las periferias sur y
oriente hacia distintos destinos de la metrópoli. Las
Figuras~\ref{fig:df-casos-1-4}--\ref{fig:df-casos-9-12} presentan los trazados
cartográficos de cada caso; el Cuadro~\ref{tab:df-casos} consolida los valores
numéricos al final de esta sección.

% --- Casos 1–4 ---------------------------------------------------------------
\begin{figure}[H]
    \centering
    \begin{subfigure}[t]{0.48\textwidth}
        \figinclude[width=\textwidth]{Figures/Cap5/FactorDesviacion/Mapa/NB04_mapa04_caso01.png}
        \caption{Caso~1: Milpa Alta → Autódromo. $DF = 1.29$.}
        \label{fig:df-c01}
    \end{subfigure}
    \hfill
    \begin{subfigure}[t]{0.48\textwidth}
        \figinclude[width=\textwidth]{Figures/Cap5/FactorDesviacion/Mapa/NB04_mapa05_caso02.png}
        \caption{Caso~2: Milpa Alta → Reforma (CBD). $DF = 1.37$.}
        \label{fig:df-c02}
    \end{subfigure}

    \vspace{0.5em}
    \begin{subfigure}[t]{0.48\textwidth}
        \figinclude[width=\textwidth]{Figures/Cap5/FactorDesviacion/Mapa/NB04_mapa06_caso03.png}
        \caption{Caso~3: Milpa Alta → Santa Fe (CBD). $DF = 1.27$.}
        \label{fig:df-c03}
    \end{subfigure}
    \hfill
    \begin{subfigure}[t]{0.48\textwidth}
        \figinclude[width=\textwidth]{Figures/Cap5/FactorDesviacion/Mapa/NB04_mapa07_caso04.png}
        \caption{Caso~4: Milpa Alta → FES Acatlán. $DF = 1.35$.}
        \label{fig:df-c04}
    \end{subfigure}
    \caption[Casos de estudio del Factor de Desviación: 1–4]{Rutas de los casos 1 al 4.
    Origen común: Milpa Alta hacia distintos destinos del norte y poniente.}
    \label{fig:df-casos-1-4}
    \fuentefigura{Elaboración propia con \texttt{VFTModel} (corrida 2026-04-27).}
\end{figure}

% --- Casos 5–8 ---------------------------------------------------------------
\begin{figure}[H]
    \centering
    \begin{subfigure}[t]{0.48\textwidth}
        \figinclude[width=\textwidth]{Figures/Cap5/FactorDesviacion/Mapa/NB04_mapa08_caso05.png}
        \caption{Caso~5: Milpa Alta → C.U. $DF = 1.32$.}
        \label{fig:df-c05}
    \end{subfigure}
    \hfill
    \begin{subfigure}[t]{0.48\textwidth}
        \figinclude[width=\textwidth]{Figures/Cap5/FactorDesviacion/Mapa/NB04_mapa09_caso06.png}
        \caption{Caso~6: Milpa Alta → Coyoacán. $DF = 1.34$.}
        \label{fig:df-c06}
    \end{subfigure}

    \vspace{0.5em}
    \begin{subfigure}[t]{0.48\textwidth}
        \figinclude[width=\textwidth]{Figures/Cap5/FactorDesviacion/Mapa/NB04_mapa10_caso07.png}
        \caption{Caso~7: Santa Catarina → C.U. $DF = 1.82$.}
        \label{fig:df-c07}
    \end{subfigure}
    \hfill
    \begin{subfigure}[t]{0.48\textwidth}
        \figinclude[width=\textwidth]{Figures/Cap5/FactorDesviacion/Mapa/NB04_mapa11_caso08.png}
        \caption{Caso~8: Santa Catarina → FES Acatlán. $DF = 1.42$.}
        \label{fig:df-c08}
    \end{subfigure}
    \caption[Casos de estudio del Factor de Desviación: 5–8]{Rutas de los casos 5 al 8.
    Destinos en el sur de la ciudad y primeros casos desde Santa Catarina.}
    \label{fig:df-casos-5-8}
    \fuentefigura{Elaboración propia con \texttt{VFTModel} (corrida 2026-04-27).}
\end{figure}

% --- Casos 9–12 --------------------------------------------------------------
\begin{figure}[H]
    \centering
    \begin{subfigure}[t]{0.48\textwidth}
        \figinclude[width=\textwidth]{Figures/Cap5/FactorDesviacion/Mapa/NB04_mapa12_caso09.png}
        \caption{Caso~9: Milpa Alta → Santa Catarina. $DF = 2.52$.}
        \label{fig:df-c09}
    \end{subfigure}
    \hfill
    \begin{subfigure}[t]{0.48\textwidth}
        \figinclude[width=\textwidth]{Figures/Cap5/FactorDesviacion/Mapa/NB04_mapa13_caso10.png}
        \caption{Caso~10: Milpa Alta → Ajusco. $DF = 1.42$.}
        \label{fig:df-c10}
    \end{subfigure}

    \vspace{0.5em}
    \begin{subfigure}[t]{0.48\textwidth}
        \figinclude[width=\textwidth]{Figures/Cap5/FactorDesviacion/Mapa/NB04_mapa14_caso11.png}
        \caption{Caso~11: Ajusco → Preparatoria 1. $DF = 1.43$.}
        \label{fig:df-c11}
    \end{subfigure}
    \hfill
    \begin{subfigure}[t]{0.48\textwidth}
        \figinclude[width=\textwidth]{Figures/Cap5/FactorDesviacion/Mapa/NB04_mapa15_caso12.png}
        \caption{Caso~12: P. Los Olivos → Centro Tláhuac. $DF = 1.06$.}
        \label{fig:df-c12}
    \end{subfigure}
    \caption[Casos de estudio del Factor de Desviación: 9–12]{Rutas de los casos 9 al 12.
    El caso~9 (sur-oriente transversal) registra el mayor $DF$ del conjunto.}
    \label{fig:df-casos-9-12}
    \fuentefigura{Elaboración propia con \texttt{VFTModel} (corrida 2026-04-27).}
\end{figure}

% --- Extremos globales de la muestra aleatoria -------------------------------
Para situar los casos de estudio en el contexto de la muestra completa, las
Figuras~\ref{fig:df-mejor-caso} y~\ref{fig:df-peor-caso} ilustran los extremos
globales obtenidos en los 100 pares O-D aleatorios.

\begin{figure}[H]
    \centering
    \begin{subfigure}[t]{0.48\textwidth}
        \figinclude[width=\textwidth]{Figures/Cap5/FactorDesviacion/Mapa/NB04_mapa02_ruta_mas_eficiente.png}
        \caption{Mejor caso: Plutarco Elías Calles → Fuente de Loreto. $DF = 0.97$.}
        \label{fig:df-mejor-caso}
    \end{subfigure}
    \hfill
    \begin{subfigure}[t]{0.48\textwidth}
        \figinclude[width=\textwidth]{Figures/Cap5/FactorDesviacion/Mapa/NB04_mapa03_ruta_mas_ineficiente.png}
        \caption{Peor caso: Carlos A. Vidal → 20 de Noviembre. $DF = 4.34$.}
        \label{fig:df-peor-caso}
    \end{subfigure}
    \caption[Extremos del Factor de Desviación en la muestra aleatoria]{
    Rutas con menor y mayor Factor de Desviación en la muestra aleatoria de 100 pares O-D.}
    \label{fig:df-extremos}
    \fuentefigura{Elaboración propia con \texttt{VFTModel} (corrida 2026-04-27).}
\end{figure}

% --- Tabla resumen de los 12 casos -------------------------------------------
\begin{table}[H]
    \centering
    \caption[Factor de Desviación — resumen de los doce casos de estudio]{
    Resultados del Factor de Desviación para los doce casos estratégicos de la
    \zmvm{} ilustrados en las figuras anteriores. Corrida: 2026-04-27.}
    \label{tab:df-casos}
    \small
    \begin{tabularx}{\textwidth}{|l|X|X|r|r|r|}
        \hline
        \textbf{\#} & \textbf{Origen} & \textbf{Destino} &
        \textbf{$D_{\text{red}}$ (km)} & \textbf{$D_{\text{Hav.}}$ (km)} & \textbf{$DF$} \\
        \hline
         1 & Milpa Alta        & Autódromo / Centro-Oriente   & 29.18 & 22.59 & 1.29 \\
         2 & Milpa Alta        & Paseo de la Reforma (CBD)    & 37.73 & 27.50 & 1.37 \\
         3 & Milpa Alta        & Santa Fe (CBD)               & 36.42 & 28.62 & 1.27 \\
         4 & Milpa Alta        & FES Acatlán                  & 48.34 & 35.73 & 1.35 \\
         5 & Milpa Alta        & Ciudad Universitaria         & 26.25 & 19.96 & 1.32 \\
         6 & Milpa Alta        & Coyoacán                     & 25.24 & 18.86 & 1.34 \\
         7 & Santa Catarina    & Ciudad Universitaria         & 42.06 & 23.10 & 1.82 \\
         8 & Santa Catarina    & FES Acatlán                  & 45.64 & 32.16 & 1.42 \\
         9 & Milpa Alta        & Santa Catarina (Sur-Oriente) & 37.22 & 14.79 & 2.52 \\
        10 & Milpa Alta        & Ajusco                       & 29.19 & 20.56 & 1.42 \\
        11 & Ajusco            & Preparatoria 1               & 15.91 & 11.16 & 1.43 \\
        12 & Parque Los Olivos & Centro de Tláhuac            &  2.20 &  2.08 & 1.06 \\
        \hline
        \multicolumn{5}{|l|}{\textbf{Mejor caso global (muestra aleatoria):}
        Plutarco \DIFdelbeginFL \DIFdelFL{Elias }\DIFdelendFL \DIFaddbeginFL \DIFaddFL{Elías }\DIFaddendFL Calles → Fuente de Loreto} & \textbf{0.97} \\
        \multicolumn{5}{|l|}{\textbf{Peor caso global (muestra aleatoria):}
        Carlos A. Vidal → 20 de Noviembre} & \textbf{4.34} \\
        \hline
    \end{tabularx}
    \fuentefigura{Elaboración propia con \texttt{VFTModel} (corrida 2026-04-27).}
\end{table}

El análisis revela que los viajes periféricos sur-oriente (Caso~9: $DF = 2.52$,
Figura~\ref{fig:df-c09}) registran la mayor ineficiencia entre los casos de
estudio, obligando al usuario a recorrer 37.2~km de red para cubrir una distancia
lineal de 14.8~km. Esta situación responde a la ausencia de conectividad
transversal en el arco sur-oriente de la ciudad, que impone al tránsito
perimetral el ingreso al centro metropolitano como nodo obligado, confirmando la
hipótesis sobre la necesidad de un corredor anillar en el Periférico.

La Figura~\ref{fig:mapa-df-presentacion} sintetiza estos hallazgos mediante
cartografía de presentación del Factor de Desviación sobre el conjunto de nodos
de la red. El gradiente de color verde-a-rojo permite identificar directamente las
zonas donde la infraestructura existente impone los mayores rodeos, reforzando de
forma visual el argumento sobre la necesidad del corredor anillar en el arco
sur-oriente. La versión a página completa se encuentra en la
Figura~\ref{fig:anx-mapa-factor-desviacion} del
Anexo~\ref{anx:cartografia-presentacion}.

\begin{figure}[H]
    \centering
    \figinclude[width=0.90\textwidth]{Figures/Mapas/Mapa_DF}
    \caption[Cartografía de presentación — Factor de Desviación por nodo]{
    Factor de Desviación ($DF$) por nodo en la red de transporte de la \zmvm{}.
    El gradiente verde-a-rojo indica la relación entre la distancia real de viaje
    y la distancia euclidiana: los nodos en rojo señalan las zonas donde las
    rutas disponibles imponen desviaciones significativas, penalizando con mayor
    intensidad a la periferia sur y oriente de la cuenca metropolitana.}
    \label{fig:mapa-df-presentacion}
    \fuentefigura{Elaboración propia con \texttt{Transport-gis-zmvm-mjg};
    datos: GTFS SEMOVI 2024 e INEGI 2023.}
\end{figure}

%% ── 5-Resultados/5-2-2-Resultados-VFT-Fase2.tex
\subsection*{5 2 2 Resultados VFT Fase2}
% !TeX root = ../tesis.tex
%%%%%%%%%%%%%%%%%%%%%%%%%%%%%%%%%%%%%%%%%%%%%%%%%%%%%%%%%%%%%%%%%%%
% SECCIÓN 5.2-2 — RESULTADOS VFT: FASE 2
% Ponderación dinámica de la red
%%%%%%%%%%%%%%%%%%%%%%%%%%%%%%%%%%%%%%%%%%%%%%%%%%%%%%%%%%%%%%%%%%%
\section{\bajoRevision{Fase 2: Ponderación Dinámica de la Red}}
\label{sec:resultados-fase2}

\notaRevision{Los indicadores de esta fase se encuentran en revisión con el
comité tutor. Los coeficientes de fricción vial y la penalización por
transferencia han sido calculados e integrados al grafo; el \termino{Node Score}
se encuentra en etapa de planificación, pendiente de la disponibilidad de las
vistas materializadas MV1 y MV2 de \textbf{Apimetro}.}

La Fase~2 del Modelo VFT enriquece el grafo construido en la Fase~1 con dos
capas de pesos dinámicos que aproximan las condiciones reales de operación de la
red: el \textbf{Coeficiente de Fricción Vial} ($CF$), que penaliza el tiempo de
traslado en función del nivel de interferencia del tráfico urbano sobre cada
tipo de vía, y la \textbf{Penalización por Transferencia} ($T_b$), que incorpora
el costo de caminata y espera en cada transbordo intermodal. Estos dos pesos
conforman el componente dinámico de la \textbf{Ecuación VFT}:

\begin{equation}
    w(u,v) = \frac{D_{uv}}{V_{\text{modo}}} \cdot CF_{\text{vía}} + T_b
    \label{eq:vft-arista}
\end{equation}

\begin{itemize}
    \item $D_{uv}$: distancia geodésica del segmento calculada mediante la
    fórmula de \termino{Haversine} (metros).
    \item $V_{\text{modo}}$: velocidad comercial ajustada del sistema de
    transporte (km/h), extraída del \geojson{} de \textbf{Apimetro} o
    sustituida por el valor de \termino{fallback} calibrado por sistema.
    \item $CF_{\text{vía}}$: coeficiente de fricción vial según el tipo de
    derecho de vía del segmento (adimensional, $\geq 1$).
    \item $T_b$: penalización por transferencia peatonal en minutos; se aplica
    exclusivamente en las aristas de transbordo generadas por el
    \termino{Snapping Lógico Geográfico}.
\end{itemize}

Para las aristas de servicio normal, el peso de enrutamiento es
$\mathit{weight} = (D_{uv}/V_{\text{modo}}) \cdot CF_{\text{vía}}$; el
$T_b$ queda registrado como atributo informativo\DIFaddbegin \DIFadd{, }\DIFaddend pero no penaliza el
segmento. La penalización por transferencia actúa únicamente sobre aristas de
tipo \texttt{transfer}, como se describe en la Sección~\ref{subsec:res-penalizacion-transferencia}.

% ----------------------------------------------------------------
\subsection{Coeficiente de Fricción Vial (\texorpdfstring{$CF$}{CF})}
\label{subsec:coeficiente-friccion}
% ----------------------------------------------------------------

El \textbf{Coeficiente de Fricción Vial} ($CF$) traduce el tipo de
\termino{derecho de vía} de cada segmento en un penalizador multiplicativo
sobre el tiempo de traslado. La premisa es que circular por un túnel del
Metro —aislado del tráfico en superficie— implica un costo temporal muy
distinto al de un camión de la \textsc{rtp} inmerso en el flujo vehicular de
Insurgentes o Periférico.

\subsubsection{Constante de Saturación Vial}
\label{subsubsec:beta-saturacion}

El parámetro base del coeficiente es la constante de saturación $\beta$,
derivada del \termino{TomTom Traffic Index} para la \zmvm:

\begin{equation}
    \beta_{\text{ZMVM}} = 0.759
    \label{eq:beta-zmvm}
\end{equation}

Este valor indica que la ciudad opera, en condiciones promedio, al 76~\% de
su capacidad vial. Es decir, un vehículo en tráfico mixto experimenta, en
promedio, un sobretiempo del 76~\% respecto al recorrido en flujo libre.

\subsubsection{Fórmula y Valores por Derecho de Vía}
\label{subsubsec:formula-cf}

El coeficiente $CF$ se obtiene a partir de un parámetro de escala $\alpha$
que modula la exposición de cada tipo de vía a la saturación urbana:

\begin{equation}
    CF = 1 + \alpha \cdot \beta_{\text{ZMVM}}
    \label{eq:cf-formula}
\end{equation}

\begin{itemize}
    \item $\alpha \in \{0.0,\ 0.2,\ 0.5,\ 1.0\}$: fracción de la saturación
    urbana que afecta el segmento, determinada por su tipo de derecho de vía.
    \item $\beta_{\text{ZMVM}} = 0.759$: constante de saturación de la \zmvm
    (\termino{TomTom Traffic Index}).
    \item $CF \geq 1$: un valor de 1.0 indica flujo ideal sin interferencia;
    valores superiores representan el penalizador proporcional al congestionamiento.
\end{itemize}

El Cuadro~\ref{tab:cf-por-via} sintetiza los cuatro tipos de derecho de vía
modelados, los sistemas de transporte correspondientes y el valor numérico del
coeficiente resultante.

\begin{table}[H]
    \centering
    \caption[Coeficiente de Fricción Vial por tipo de derecho de vía]{Coeficiente
    de Fricción Vial ($CF$) por tipo de derecho de vía en el Modelo VFT.
    $\alpha$ escala el impacto de la saturación urbana
    $\beta_{\text{ZMVM}} = 0.759$.}
    \label{tab:cf-por-via}
    \begin{tabularx}{\textwidth}{|l|c|c|X|}
        \hline
        \textbf{Derecho de vía} & \textbf{$\alpha$} & \textbf{$CF$} &
        \textbf{Sistemas representativos} \\
        \hline
        Exclusivo   & 0.0 & 1.000 & Metro, Tren Suburbano \\
        Confinado   & 0.2 & 1.152 & Metrobús, Cablebús, Mexicable \\
        Compartido  & 0.5 & 1.380 & Trolebús, Tren Ligero \\
        Mixto       & 1.0 & 1.759 & \textsc{rtp}, Corredores Concesionados,
                                    Mexibús, Pumabús \\
        \hline
    \end{tabularx}
    \fuentefigura{Elaboración propia con base en \texttt{impedance.py}
    (VFTModel, corrida 2026-04-18). Los valores $\alpha$ tienen como referencia
    normativa el \textit{Manual de Calles: Diseño Vial para Ciudades Mexicanas}
    \parencite{sedatu2019manual}.}
\end{table}

El rango del coeficiente $CF \in [1.000,\ 1.759]$ implica que el tiempo de
traslado en un segmento de vía mixta (RTP, Corredores Concesionados) puede ser
hasta un 75.9~\% mayor que el mismo recorrido bajo condiciones de flujo ideal.
Esta diferencia no refleja únicamente la velocidad operacional del vehículo —que
ya está calibrada en $V_{\text{modo}}$—\DIFaddbegin \DIFadd{, }\DIFaddend sino el efecto sistemático de las
detenciones por semáforos, cruces peatonales y congestión que son inherentes a
los segmentos de superficie compartida.

En el caso del sistema de derecho de vía desconocido, el Modelo VFT asigna
por defecto el valor de vía \texttt{mixto} ($CF = 1.759$), aplicando el
principio de escenario más conservador para no subestimar los tiempos de traslado.

% ----------------------------------------------------------------
\subsection{Penalización por Transferencia (\texorpdfstring{$T_b$}{Tb})}
\label{subsec:res-penalizacion-transferencia}
% ----------------------------------------------------------------

La \textbf{Penalización por Transferencia} ($T_b$) modela el costo temporal \DIFaddbegin \DIFadd{en el
}\DIFaddend que un pasajero incurre al cambiar de sistema o línea dentro de la red. A diferencia
del tiempo de traslado —que depende de la distancia y la velocidad del vehículo—,
el costo de transbordo consta de dos componentes independientes: el tiempo de
caminata entre plataformas y el tiempo de espera en el andén de destino.

\subsubsection{Fórmula de la Penalización}
\label{subsubsec:formula-tb}

Para cada arista de transbordo generada por el \termino{Snapping Lógico Geográfico}
entre los nodos $u$ (origen del transbordo) y $v$ (destino del transbordo), el
peso de la arista es:

\begin{equation}
    T_b(u, v) = t_{\text{caminata}}(u, v) + \frac{F_v}{2}
    \label{eq:tb-formula}
\end{equation}

\begin{itemize}
    \item $t_{\text{caminata}}(u, v)$: tiempo de caminata entre los dos nodos
    a una velocidad peatonal normativa de 5~km/h (1.38~m/s), equivalente al
    límite superior establecido en el \textit{Manual de Calles SEDATU}
    \parencite{sedatu2019manual} para infraestructura de transbordo sin detenciones
    en semáforo.
    \item $F_v$: frecuencia de servicio del sistema destino $v$ (minutos entre
    unidades); si no está reportada en el \geojson{}, se utiliza el valor
    de \termino{fallback} calibrado por sistema.
    \item $F_v / 2$: tiempo de espera esperado, calculado como la mitad de la
    frecuencia bajo la hipótesis de llegada aleatoria del pasajero al andén
    (\termino{random arrival assumption}).
\end{itemize}

El tiempo de caminata se calcula como:

\begin{equation}
    t_{\text{caminata}}(u, v) = \frac{D_{uv}}{V_{\text{peatonal}}}
    = \frac{D_{uv}}{83.33\ \text{m/min}}
    \label{eq:t-caminata}
\end{equation}

donde $D_{uv}$ es la distancia geodésica entre los dos nodos (Haversine) en
metros, y $V_{\text{peatonal}} = 5000\ \text{m} / 60\ \text{min} = 83.33\ \text{m/min}$.

\subsubsection{Valores de Frecuencia de Fallback}
\label{subsubsec:fallback-frecuencia}

El Cuadro~\ref{tab:fallback-frecuencia} presenta los valores de frecuencia de
\termino{fallback} calibrados para cada sistema de transporte, junto con el
costo de abordaje resultante ($F_v / 2$). Estos valores representan intervalos
promedio en condiciones de operación normal en hora no pico.

\begin{table}[H]
    \centering
    \caption[Frecuencias de \textit{fallback} y costo de abordaje por sistema]{Frecuencias
    de servicio de \termino{fallback} y penalización de abordaje resultante
    ($T_b^{\text{espera}} = F_v / 2$) por sistema de transporte en el Modelo VFT.}
    \label{tab:fallback-frecuencia}
    \begin{tabular}{|l|c|c|}
        \hline
        \textbf{Sistema} & \textbf{$F_v$ (min)} & \textbf{$F_v / 2$ (min)} \\
        \hline
        \textsc{Metro} (STC)          & 3.0  & 1.50 \\
        Metrobús (MB)                 & 5.0  & 2.50 \\
        Tren Ligero (TL)              & 10.0 & 5.00 \\
        Trolebús (TROLE)              & 8.0  & 4.00 \\
        \textsc{rtp}                  & 15.0 & 7.50 \\
        Tren Suburbano (SUB)          & 12.0 & 6.00 \\
        Cablebús (CBB)                & 10.0 & 5.00 \\
        Corredores Concesionados (CC) & 15.0 & 7.50 \\
        Mexibús                       & 6.0  & 3.00 \\
        \hline
    \end{tabular}
    \fuentefigura{Elaboración propia con base en \texttt{graph\_builder.py}
    (VFTModel, corrida 2026-04-18). Los valores corresponden a frecuencias
    promedio en horario valle; las frecuencias en hora pico son reportadas por
    los operadores como atributo del \geojson{} y tienen precedencia sobre los
    \termino{fallbacks}.}
\end{table}

\subsubsection{Decisión de Diseño: Boarding Cost Informativo}
\label{subsubsec:boarding-cost-diseno}

Una precisión metodológica relevante: el costo de abordaje $F_v / 2$ \textbf{no
se suma al peso de las aristas de servicio normal}. Solo actúa como componente
del peso de las aristas de transbordo peatonal. Esta decisión responde a que el
pasajero ya está dentro del vehículo en los segmentos normales; el tiempo de
espera por frecuencia solo es pertinente en el momento de abordar un nuevo
sistema. Sumarlo en cada arista normal equivaldría a penalizar al pasajero en
cada estación intermedia, sobreestimando sistemáticamente los tiempos de ruta.

En consecuencia, el Modelo VFT registra $F_v / 2$ como atributo informativo
en las aristas de servicio (\texttt{boarding\_cost\_min}) para análisis
descriptivos, pero el atributo \texttt{weight} de estas aristas solo contiene
el tiempo de traslado con fricción.

% ----------------------------------------------------------------
\subsection{\bajoRevision{Node Score — Estado de Implementación}}
\label{subsec:node-score}
% ----------------------------------------------------------------

\notaRevision{El \termino{Node Score} ($\text{VFT\_score}$) se encuentra en etapa
de planificación. Su cálculo requiere los datos de afluencia por línea (MV1) y
los polígonos AGEB con atributos censales (MV2) de \textbf{Apimetro}, aún no
disponibles en la fecha de cierre de este documento. La arquitectura del indicador,
sus fórmulas y las dependencias de datos se presentan a continuación como
fundamento para la etapa de implementación.}

El \termino{Node Score} ($\text{VFT\_score}_i$) constituye la segunda capa de
pesos del Modelo VFT: mientras que los pesos de arista modelan el costo temporal
de recorrer la red, el \termino{Node Score} asigna a cada nodo una medida de
\textbf{importancia compuesta} que combina propiedades topológicas de la red con
información de demanda observada y potencial. Este indicador es el que alimenta
el análisis de vulnerabilidad de la Fase~3.

\subsubsection{Fórmula del Node Score}
\label{subsubsec:formula-vft-score}

\begin{equation}
    \text{VFT\_score}_i = \alpha_1 \hat{k}_i + \alpha_2 \hat{a}_i + \alpha_3 \hat{p}_i
    \label{eq:vft-score}
\end{equation}

donde los tres componentes normalizados al rango $[0, 1]$ son:

\begin{itemize}
    \item $\hat{k}_i$: \textbf{Fuerza Capilar normalizada} (min-max sobre todos
    los nodos del grafo). Calculada en la Fase~1 para los 11{,}115 nodos.
    \item $\hat{a}_i$: \textbf{Afluencia estimada por estación normalizada}.
    Derivada de la afluencia diaria por línea ($A_\ell$, fuente MV1) desagregada
    internamente por VFT mediante ponderación por grado topológico:
    \begin{equation}
        a_i = \sum_{\ell \ni i} A_\ell \cdot
        \frac{g_i}{\displaystyle\sum_{j \in \ell} g_j}
        \label{eq:desagregacion-afluencia}
    \end{equation}
    donde $g_i = \deg(i)$ es el grado del nodo $i$ en el grafo NetworkX.
    \item $\hat{p}_i$: \textbf{Población captada en buffer de 800~m normalizada}.
    Calculada intersectando un radio de 800~m alrededor de cada estación con los
    polígonos AGEB (fuente MV2), ponderando cada AGEB por la fracción de su área
    dentro del buffer:
    \begin{equation}
        p_i = \sum_{j \in \mathcal{A}_i} \text{pob}_j \cdot
        \frac{\text{área}(B_i \cap \text{AGEB}_j)}{\text{área}(\text{AGEB}_j)}
        \label{eq:poblacion-captada}
    \end{equation}
    donde $B_i$ es el buffer circular de 800~m en proyección métrica EPSG:6372
    (México Centro) y $\mathcal{A}_i$ el conjunto de AGEBs que lo intersectan.
\end{itemize}

La normalización de cada componente usa el método \termino{min-max}:
\begin{equation}
    \hat{x}_i = \frac{x_i - \min(\mathbf{x})}{\max(\mathbf{x}) - \min(\mathbf{x})}
    \label{eq:normalizacion-minmax}
\end{equation}

\subsubsection{Determinación de los Pesos \texorpdfstring{$\alpha$}{\textalpha}}
\label{subsubsec:pesos-alpha}

Los pesos $\alpha_1, \alpha_2, \alpha_3$ se determinarán en dos etapas
progresivas una vez que los datos de MV1 y MV2 estén disponibles:

\begin{enumerate}
    \item \textbf{Validación inicial (pesos iguales):} $\alpha_1 = \alpha_2 =
    \alpha_3 = 1/3$. Objetivo: verificar que nodos de alta importancia reconocida
    (Pantitlán, Insurgentes, Pino Suárez, Balderas) aparecen en el top-10 del
    \termino{ranking}. Si no lo hacen, el problema se ubica en los datos de entrada.
    \item \textbf{Pesos por entropía (CRITIC):} los pesos finales se derivan de
    la variabilidad observada de cada componente sobre todos los nodos:
    \begin{equation}
        \alpha_n = \frac{\sigma_n}{\displaystyle\sum_{m=1}^{3} \sigma_m}
        \label{eq:alpha-entropia}
    \end{equation}
    donde $\sigma_n = \text{std}(\hat{x}_n)$. Un componente más discriminante
    (mayor dispersión entre nodos) recibe mayor peso, penalizando automáticamente
    los componentes redundantes.
\end{enumerate}

\subsubsection{Dependencias de Datos y Estado}
\label{subsubsec:node-score-estado}

El Cuadro~\ref{tab:node-score-dependencias} resume los componentes del
\termino{Node Score}, su fuente de datos y su estado de disponibilidad a la
fecha de cierre del documento.

\begin{table}[H]
    \centering
    \caption[Dependencias del Node Score por componente]{Componentes del
    $\text{VFT\_score}$, fuentes de datos requeridas y estado de implementación.}
    \label{tab:node-score-dependencias}
    \begin{tabularx}{\textwidth}{|l|X|l|l|}
        \hline
        \textbf{Componente} & \textbf{Fuente} & \textbf{Responsable} & \textbf{Estado} \\
        \hline
        $\hat{k}_i$ Fuerza Capilar    & VFT interno (\path{capillary_strength.py}) & VFTModel & Calculado \\
        $\hat{a}_i$ Afluencia         & MV1 Apimetro — afluencia por línea & Apimetro & Pendiente \\
        $\hat{p}_i$ Población captada & MV2 Apimetro — AGEBs Censo INEGI 2020 & Apimetro & Pendiente \\
        Pesos $\alpha$ (entropía)     & Derivados de $\sigma(\hat{k}, \hat{a}, \hat{p})$ & VFTModel & Pendiente MV1+MV2 \\
        \hline
    \end{tabularx}
    \fuentefigura{Elaboración propia con base en
    \texttt{weighted\_graph\_enrichment.md} y
    \texttt{apimetro\_requerimientos\_mvs.md} (VFTModel, 2026-04-28).}
\end{table}

El diseño del \termino{Node Score} garantiza la \textbf{independencia entre
componentes}: la desagregación de afluencia usa el grado topológico crudo
$g_i = \deg(i)$ —no la Fuerza Capilar $k_i$— para evitar que $\hat{k}_i$
domine doblemente el cómputo. Antes de asignar los pesos finales por entropía,
se verificará la correlación $\rho(\hat{k}, \hat{a})$: si $\rho \geq 0.75$,
los dos componentes se fusionarán en un índice combinado $\delta_i = \beta
\hat{k}_i + (1-\beta)\hat{a}_i$ antes de calcular el \termino{score} final.

%% ── Figures/TiKz_Libraries/Diagramas/Arboles/tg-arbol-directorios.tex
\subsection*{tg arbol directorios}
%DIF >  Árbol: Estructura de directorios de Transport-gis-zmvm-mjg
%DIF >  Usado en: 4-Capas-Codigo/4-3-Visualizacion.tex  (fig:tg-arbol)
%DIF >  Requiere: \usepackage{forest} — declarado en Latex/Comands.tex
\DIFaddbegin \iffalse
\begin{forest}
  \DIFadd{for tree = }{
    \DIFadd{font   = }\footnotesize\ttfamily\DIFadd{,
    grow'  = 0,
    child anchor  = west,
    parent anchor = east,
    anchor = west,
    draw   = gray!45,
    fill   = gray!7,
    rounded corners = 2pt,
    inner xsep = 6pt,
    inner ysep = 3pt,
    l sep  = 20pt,
    s sep  = 5pt,
    edge path = }{
      \noexpand\path [\DIFadd{draw=gray!50, thin, rounded corners=2pt}]
        \DIFadd{(!u.parent anchor) -- +(8pt,0) |- (.child anchor)
        }\forestoption{edge label}\DIFadd{;
    }}\DIFadd{,
  }}
  [\DIFadd{Transport-gis-zmvm-mjg/, fill=unamAzul!12, draw=unamAzul!70,
                             font=}\small\bfseries\ttfamily
    [\DIFadd{analysis/
      }[\DIFadd{notebooks/}]
      [\DIFadd{scripts/
        }[\DIFadd{apimetro\_client.py,
          draw=none, fill=none, font=}\footnotesize\ttfamily\color{unamAzul}]
        [\DIFadd{vft\_runner.py,
          draw=none, fill=none, font=}\footnotesize\ttfamily\color{unamAzul}]
      ]
    ]
    [\DIFadd{data/processed/
      }[{\textit{\DIFadd{10 capas .gpkg}} \DIFadd{\textnormal{(Git~LFS)}}}\DIFadd{,
        draw=none, fill=none, font=}\footnotesize\itshape\ttfamily]
    ]
    [\DIFadd{maps/, fill=unamOro!8, draw=unamOro!60
      }[\DIFadd{exports/
        }[{\textit{\DIFadd{*.pdf · *.png}}}\DIFadd{,
          draw=none, fill=none, font=}\footnotesize\itshape\ttfamily]
      ]
      [\DIFadd{fonts/}]
      [\DIFadd{projects/
        }[\textit{\DIFadd{*.qgz}}\DIFadd{,
          draw=none, fill=none, font=}\footnotesize\itshape\ttfamily]
      ]
      [\DIFadd{templates/
        }[{\textit{\DIFadd{*.qpt}} \DIFadd{\textnormal{(paletas UNAM)}}}\DIFadd{,
          draw=none, fill=none, font=}\footnotesize\itshape\ttfamily]
      ]
    ]
    [\DIFadd{pipeline/
      }[\DIFadd{run\_pipeline.py,
        draw=none, fill=none, font=}\footnotesize\ttfamily\color{unamAzul}]
    ]
    [\DIFadd{setup/
      }[\DIFadd{install\_fonts.sh,
        draw=none, fill=none, font=}\footnotesize\ttfamily\color{unamAzul}]
      [\DIFadd{install\_fonts.bat,
        draw=none, fill=none, font=}\footnotesize\ttfamily\color{unamAzul}]
    ]
  ]
\end{forest}
\fi
 \DIFaddend

%% ── Figures/TiKz_Libraries/Diagramas/macro-arquitectura-stack.tex
\subsection*{macro arquitectura stack}
%DIF >  Diagrama: Macro-Arquitectura del Stack de Investigación (3 capas apiladas)
%DIF >  Usado en: 4-Capas-Codigo/4-1-Apimetro.tex  (fig:macro-arquitectura-stack)
%DIF >  Fix 2026-08-28: font=\small en estilo capa para que \hfill funcione; más espaciado vertical
\DIFaddbegin \begin{tikzpicture}[
    capa/.style = {
        rectangle, draw=unamAzul!60,
        text width=11cm, minimum height=2.0cm,
        align=left, inner sep=14pt,
        rounded corners=4pt, font=\small
    },
    flecha/.style = {Stealth-Stealth, thick, color=unamAzul!55},
    etiq/.style   = {font=\scriptsize\itshape, color=gray!70,
                     align=center, text width=2.2cm}
]

%% — Capas (base más oscura, tope más clara) —
\node[capa, fill=unamAzul!28] (c1) {
    \textbf{Capa 1 — Apimetro} \hfill \texttt{Go · PostgreSQL/PostGIS}\\[5pt]
    Ingesta \gtfs{} · ETL GeoJSON · \api{} REST · Esquema relacional
};
\node[capa, fill=unamAzul!16, above=1.1cm of c1] (c2) {
    \textbf{Capa 2 — VFTModel} \hfill \texttt{Python · NetworkX}\\[5pt]
    Motor analítico · Constructor del grafo · 8 indicadores topológicos
};
\node[capa, fill=unamAzul!7,  above=1.1cm of c2] (c3) {
    \textbf{Capa 3 — Transport-gis-zmvm-mjg} \hfill \texttt{QGIS · Python}\\[5pt]
    Capas GeoPackage · Plantillas \textsc{unam} · Cartografía de presentación
};

%% — Flechas entre capas con etiquetas de intercambio —
\draw[flecha]
    ($(c1.north) + (0.9, 0)$) -- node[etiq, left=4pt] {\gtfs{} / GeoJSON}
    ($(c2.south) + (0.9, 0)$);

\draw[flecha]
    ($(c2.north) + (0.9, 0)$) -- node[etiq, left=4pt] {grafo / resultados}
    ($(c3.south) + (0.9, 0)$);

%% — Flecha lateral "FLUJO DE DATOS" —
\draw[-Stealth, thick, color=unamAzul!50]
    ($(c1.south east) + (0.55, 0)$) -- ($(c3.north east) + (0.55, 0)$)
    node[midway, right=4pt,
         font=\footnotesize\bfseries\color{unamAzul!65},
         rotate=90, anchor=south] {FLUJO DE DATOS};

\end{tikzpicture}
 \DIFaddend

%% ── Figures/TiKz_Libraries/Diagramas/tg-capas-geopackage.tex
\subsection*{tg capas geopackage}
%DIF >  Diagrama: Capas GeoPackage de Transport-gis organizadas por categoría
%DIF >  Usado en: 4-Capas-Codigo/4-3-Visualizacion.tex  (fig:tg-capas)
%DIF >  Fix 2026-08-28: node distance aumentado de 0.12cm a 0.28cm para dar respiración al texto
\DIFaddbegin \begin{tikzpicture}[
  hdr1/.style = {rectangle, rounded corners=3pt,
                 draw=gray!60, fill=gray!55,
                 minimum width=4.2cm, minimum height=0.80cm,
                 align=center, font=\small\bfseries\color{white}},
  lyr1/.style = {rectangle, rounded corners=2pt,
                 draw=gray!45, fill=gray!8,
                 minimum width=4.2cm, minimum height=0.72cm,
                 align=center, font=\footnotesize},
  hdr2/.style = {rectangle, rounded corners=3pt,
                 draw=unamAzul, fill=unamAzul!70,
                 minimum width=4.2cm, minimum height=0.80cm,
                 align=center, font=\small\bfseries\color{white}},
  lyr2/.style = {rectangle, rounded corners=2pt,
                 draw=unamAzul!50, fill=unamAzul!8,
                 minimum width=4.2cm, minimum height=0.72cm,
                 align=center, font=\footnotesize},
  hdr3/.style = {rectangle, rounded corners=3pt,
                 draw=unamOro!80, fill=unamOro!75,
                 minimum width=4.2cm, minimum height=0.80cm,
                 align=center, font=\small\bfseries\color{white}},
  lyr3/.style = {rectangle, rounded corners=2pt,
                 draw=unamOro!60, fill=unamOro!10,
                 minimum width=4.2cm, minimum height=0.72cm,
                 align=center, font=\footnotesize},
  node distance = 0.28cm
]
%% — Columna 1: Contexto territorial —
\node[hdr1] (h1) at (0,0) {Contexto territorial};
\node[lyr1, below=of h1] (c1a) {\texttt{LimitesPoliticos.gpkg}};
\node[lyr1, below=of c1a] (c1b) {\texttt{PoligonosEdoMex.gpkg}};
\node[lyr1, below=of c1b] (c1c) {\texttt{PoligonoEstudio1.gpkg}};
\node[lyr1, below=of c1c] (c1d) {\texttt{MedioFisicoNatural.gpkg}};

%% — Columna 2: Red e infraestructura —
\node[hdr2] (h2) at (4.8,0) {Red e infraestructura};
\node[lyr2, below=of h2] (c2a) {\texttt{Transporte.gpkg}};
\node[lyr2, below=of c2a] (c2b) {\texttt{InfraEstructura.gpkg}};
\node[lyr2, below=of c2b] (c2c) {\texttt{Vialidades.gpkg}};

%% — Columna 3: Indicadores del modelo —
\node[hdr3] (h3) at (9.6,0) {Indicadores del modelo};
\node[lyr3, below=of h3] (c3a) {\texttt{Cobertura800m.gpkg}};
\node[lyr3, below=of c3a] (c3b) {\texttt{FuerzaCapilar.gpkg}};
\node[lyr3, below=of c3b] (c3c) {\texttt{FactorDesviacion.gpkg}};
\end{tikzpicture}
 \DIFaddend

%% ── Figures/TiKz_Libraries/Diagramas/tg-pipeline-cartografico.tex
\subsection*{tg pipeline cartografico}
%DIF >  Diagrama: Pipeline de producción cartográfica de Transport-gis
%DIF >  Usado en: 4-Capas-Codigo/4-3-Visualizacion.tex  (fig:tg-pipeline)
\DIFaddbegin \begin{tikzpicture}[
  srcbox/.style = {rectangle, rounded corners=3pt, draw=gray!55, fill=gray!8,
                   minimum width=3.0cm, minimum height=0.90cm,
                   align=center, font=\small},
  extbox/.style = {rectangle, rounded corners=3pt, draw=unamAzul, fill=unamAzul!10,
                   minimum width=3.2cm, minimum height=0.90cm,
                   align=center, font=\small},
  intbox/.style = {rectangle, rounded corners=3pt, draw=unamAzul!60, fill=unamGrisClaro,
                   minimum width=6.4cm, minimum height=1.05cm,
                   align=center, font=\small},
  outbox/.style = {rectangle, rounded corners=3pt, draw=unamOro, fill=unamOro!15,
                   minimum width=3.2cm, minimum height=0.90cm,
                   align=center, font=\small},
  ->, >=Stealth, thick
]
%% — Fuentes externas —
\node[srcbox] (gtfs)  at (0.0, 9.2) {\gtfs{} SEMOVI~2024};
\node[srcbox] (inegi) at (5.4, 9.2) {INEGI~2023 (cartografía)};

%% — Componentes externos —
\node[extbox] (api) at (0.0, 7.5) {\textbf{Apimetro}\\{\footnotesize API REST (Go)}};
\node[extbox] (vft) at (5.4, 7.5) {\textbf{VFTModel}\\{\footnotesize (Python / NetworkX)}};

%% — Transport-gis: análisis —
\node[intbox] (scripts) at (2.7, 5.5)
  {{\footnotesize\texttt{apimetro\_client.py} $\cdot$ \texttt{vft\_runner.py}}\\
   {\footnotesize\texttt{pipeline/run\_pipeline.py}}};

%% — Transport-gis: datos —
\node[intbox] (gpkg) at (2.7, 4.0)
  {{\small\texttt{data/processed/}}\\
   {\footnotesize 10 capas \texttt{.gpkg} (Git~LFS)}};

%% — Transport-gis: QGIS —
\node[intbox] (qgis) at (2.7, 2.5)
  {{\small QGIS~3.x — \texttt{maps/projects/*.qgz}}\\
   {\footnotesize\texttt{maps/templates/*.qpt} (paletas UNAM)}};

%% — Transport-gis: exports —
\node[intbox] (exports) at (2.7, 1.0)
  {{\small\texttt{maps/exports/}}\\
   {\footnotesize\texttt{*.pdf} $\cdot$ \texttt{*.png}}};

%% — Destino final —
\node[outbox] (tesis) at (2.7, -0.5)
  {\textbf{Tesis LaTeX}\\{\footnotesize\texttt{Figures/}}};

%% — Flechas: fuentes → externos —
\draw[color=gray!60] (gtfs)  -- (api);
\draw[color=gray!60] (inegi) -- (vft);

%% — Flechas: externos → scripts (convergentes, con codo exterior) —
\draw[color=unamAzul!70] (api.south) -- ++(0,-0.55) -| (scripts.north west);
\draw[color=unamAzul!70] (vft.south) -- ++(0,-0.55) -| (scripts.north east);

%% — Flujo interno —
\draw[color=unamAzul!70] (scripts) -- (gpkg);
\draw[color=unamAzul!70] (gpkg)    -- (qgis);
\draw[color=unamAzul!70] (qgis)    -- (exports);

%% — Exports → LaTeX —
\draw[color=unamOro] (exports) -- (tesis);

%% — Recuadro del repositorio Transport-gis (en fondo) —
\begin{scope}[on background layer]
  \node[draw=unamAzul!40, dashed, rounded corners=6pt, fill=unamAzul!3,
        fit={(scripts)(gpkg)(qgis)(exports)},
        inner sep=14pt] (tgbox) {};
\end{scope}
%% — Label DENTRO del recuadro, esquina sup-izq —
\node[font=\small\bfseries\color{unamAzul}, anchor=north west]
  at ($(tgbox.north west) + (4pt,-3pt)$) {Transport-gis-zmvm-mjg};

\end{tikzpicture}
 \DIFaddend

%% ── Figures/TiKz_Libraries/Diagramas/vft-arquitectura-hexagonal.tex
\subsection*{vft arquitectura hexagonal}
%DIF >  Diagrama: Arquitectura hexagonal de tres capas de VFTModel
%DIF >  Usado en: 4-Capas-Codigo/4-2-VFTModel.tex  (fig:vft-arquitectura)
\DIFaddbegin \begin{tikzpicture}[
    capa/.style = {
        rectangle, draw=unamAzul, rounded corners=5pt,
        text width=11cm, align=left,
        minimum height=1.6cm, inner sep=10pt
    },
    flecha/.style = {Stealth-Stealth, thick, color=unamAzul!55}
]
    \node[capa, fill=unamAzul!5] (api) {
        \textbf{Capa de \api{}} \hfill \texttt{src/api/}\\
        {\small Peticiones \textsc{http} · Validación Pydantic · Respuestas \textsc{json}}
    };
    \node[capa, fill=unamAzul!12, below=0.45cm of api] (core) {
        \textbf{Capa de Dominio} \hfill \texttt{src/core/}\\
        {\small Algoritmos topológicos · Modelos matemáticos · Constructor del grafo}
    };
    \node[capa, fill=unamAzul!22, below=0.45cm of core] (infra) {
        \textbf{Capa de Infraestructura} \hfill \texttt{src/infrastructure/}\\
        {\small Clientes \textsc{http} · Archivos \geojson{} locales · Colas asíncronas}
    };
    \draw[flecha] (api)  -- (core);
    \draw[flecha] (core) -- (infra);
\end{tikzpicture}
 \DIFaddend

%% ── Figures/TiKz_Libraries/Diagramas/vft-arquitectura-ports-adapters.tex
\subsection*{vft arquitectura ports adapters}
%DIF >  Diagrama: Arquitectura Ports-and-Adapters de VFTModel
%DIF >  Usado en: 4-Capas-Codigo/4-2-VFTModel.tex  (fig:VFTModel-arquitectura)
%DIF >  Fix 2026-08-28: eliminados actores externos como cajas separadas (diagrama era >18cm);
%DIF >    información de sistemas externos integrada dentro de los adaptadores.
\DIFaddbegin \begin{tikzpicture}[
    dominio/.style = {
        rectangle, draw=unamAzul, fill=unamAzul!18,
        text width=3.8cm, align=center,
        minimum height=2.4cm, rounded corners=5pt,
        font=\small\bfseries
    },
    adaptador/.style = {
        rectangle, draw=unamAzul!70, fill=unamAzul!6,
        text width=3.2cm, align=center,
        minimum height=1.8cm, rounded corners=3pt,
        font=\small
    },
    flecha/.style = {Stealth-Stealth, thick, color=unamAzul!55}
]

%% — Dominio central —
\node[dominio] (dom) {
    \textbf{Dominio}\\
    \texttt{src/core/}\\[3pt]
    {\footnotesize Algoritmos \textsc{vft}\\Constructor del grafo}
};

%% — Adaptadores alrededor del dominio —
\node[adaptador, left=0.8cm of dom] (adp-api) {
    \textbf{Adaptador \api{}}\\
    \texttt{src/api/}\\[3pt]
    {\footnotesize FastAPI \textsc{rest} $\cdot$ Go}
};
\node[adaptador, right=0.8cm of dom] (adp-viz) {
    \textbf{Adaptador Viz}\\[3pt]
    {\footnotesize Jupyter $\cdot$ \textsc{qgis}\\Transport-gis}
};
\node[adaptador, below=0.9cm of dom] (adp-per) {
    \textbf{Adaptador Persistencia}\\
    \texttt{src/infrastructure/}\\[3pt]
    {\footnotesize Redis $\cdot$ Celery\\PostgreSQL/PostGIS}
};

%% — Flechas entre adaptadores y dominio —
\draw[flecha] (adp-api) -- (dom);
\draw[flecha] (dom)     -- (adp-viz);
\draw[flecha] (dom)     -- (adp-per);

%% — Borde Docker (en fondo) —
\begin{scope}[on background layer]
    \node[draw=unamAzul!35, dashed, rounded corners=8pt, fill=unamAzul!3,
          fit={(adp-api)(adp-viz)(adp-per)(dom)},
          inner sep=16pt] (docker) {};
\end{scope}

%% — Etiqueta Docker centrada en la parte superior —
\node[font=\scriptsize\bfseries\color{unamAzul!65}, anchor=north]
    at (docker.north) {Docker};

\end{tikzpicture}
 \DIFaddend

%% ── Figures/TiKz_Libraries/Diagramas/vft-flujo-procesamiento.tex
\subsection*{vft flujo procesamiento}
%DIF >  Diagrama: Flujo de procesamiento de VFTModel (6 etapas encadenadas)
%DIF >  Usado en: 4-Capas-Codigo/4-2-VFTModel.tex  (fig:vft-flujo)
%DIF >  Nota: el \resizebox{\linewidth}{!}{...} permanece en el capítulo, envolviendo este \input
\DIFaddbegin \begin{tikzpicture}[
    node distance = 0.5cm and 0.35cm,
    caja/.style = {
        rectangle, draw=unamAzul, fill=unamAzul!7,
        text width=1.95cm, align=center,
        minimum height=1.3cm, rounded corners=4pt,
        font=\small
    },
    flecha/.style = {-Stealth, thick, color=unamAzul!70},
    etiq/.style = {
        font=\scriptsize\itshape, color=gray,
        align=center, text width=1.95cm
    }
]
    \node[caja] (apy) {\textbf{Apimetro}\\{\tiny Go Backend}};
    \node[caja, right=of apy]  (val) {\textbf{Validación}\\{\tiny Pydantic}};
    \node[caja, right=of val]  (grf) {\textbf{Constructor}\\{\tiny del Grafo}};
    \node[caja, right=of grf]  (imp) {\textbf{Impedancia}\\{\tiny Motor VFT}};
    \node[caja, right=of imp]  (ana) {\textbf{Analizadores}\\{\tiny Topológicos}};
    \node[caja, right=of ana]  (res) {\textbf{Resultados}\\{\tiny JSON / CSV}};

    \node[etiq, below=0.18cm of apy]  {Estaciones\\y Líneas\\\geojson{}};
    \node[etiq, below=0.18cm of val]  {Contratos\\de datos};
    \node[etiq, below=0.18cm of grf]  {Grafo\\NetworkX};
    \node[etiq, below=0.18cm of imp]  {Pesos en\\minutos};
    \node[etiq, below=0.18cm of ana]  {8 Indicadores};
    \node[etiq, below=0.18cm of res]  {{\upshape\api{}} REST};

    \draw[flecha] (apy) -- (val);
    \draw[flecha] (val) -- (grf);
    \draw[flecha] (grf) -- (imp);
    \draw[flecha] (imp) -- (ana);
    \draw[flecha] (ana) -- (res);
\end{tikzpicture}
 \DIFaddend

%% ── Figures/TiKz_Libraries/Diagramas/vft-pipeline-procesamiento.tex
\subsection*{vft pipeline procesamiento}
%DIF >  Diagrama: Pipeline de procesamiento interno del Modelo VFT (4 etapas)
%DIF >  Usado en: 5-Resultados/5-1-Introduccion-Resultados.tex  (fig:arq-capas)
%DIF >  Reemplaza: 3 subfiguras arq_capa1/2/3.png con captions incorrectos
\DIFaddbegin \begin{tikzpicture}[
    node distance = 0.6cm and 0.5cm,
    etapa/.style = {
        rectangle, draw=unamAzul, fill=unamAzul!8,
        text width=2.8cm, align=center,
        minimum height=1.5cm, rounded corners=4pt,
        font=\small\bfseries
    },
    dato/.style = {
        font=\scriptsize\itshape, color=unamAzul!75,
        align=center, text width=2.8cm
    },
    cifra/.style = {
        font=\footnotesize\bfseries, color=unamAzul,
        align=center, text width=2.8cm
    },
    flecha/.style = {-Stealth, thick, color=unamAzul!65}
]

%% — Etapas —
\node[etapa] (ing) {Ingesta\\GTFS};
\node[etapa, right=of ing]  (snap) {Snapping\\Lógico};
\node[etapa, right=of snap] (grf)  {Construcción\\del Grafo};
\node[etapa, right=of grf]  (imp)  {Ponderación\\Impedancia};

%% — Cifras sobre cada etapa —
\node[cifra, above=0.18cm of ing]  {22{,}766 entidades};
\node[cifra, above=0.18cm of snap] {11{,}115 nodos};
\node[cifra, above=0.18cm of grf]  {45{,}679 aristas};
\node[cifra, above=0.18cm of imp]  {grafo $G$ dirigido};

%% — Detalles bajo cada etapa —
\node[dato, below=0.18cm of ing]  {Apimetro REST\\GeoJSON};
\node[dato, below=0.18cm of snap] {KDTree · $d \leq 50$~m\\eliminación duplicados};
\node[dato, below=0.18cm of grf]  {NetworkX\\25{,}197 servicio\\20{,}482 transbordo};
\node[dato, below=0.18cm of imp]  {Motor VFT\\pesos en minutos};

%% — Flechas —
\draw[flecha] (ing)  -- (snap);
\draw[flecha] (snap) -- (grf);
\draw[flecha] (grf)  -- (imp);

%% — Nota conceptual: abstracción Red Física → Grafo Matemático —
\node[font=\scriptsize\color{gray!70}, anchor=north]
    at ($(ing.south)!0.5!(imp.south) + (0,-1.15cm)$)
    {Red Física $\longrightarrow$ Grafo Matemático $\mathcal{G}(\mathcal{V},\,\mathcal{E},\,w)$};

\end{tikzpicture}
 \DIFaddend
